% concatenated from main.tex, preamble.tex, title.tex, abstract.tex
\documentclass[11pt,twoside,a4paper]{article}

\usepackage[british]{babel}
% The Babel package presented in the introduction allows to use special characters and also translates some elements within the document. This package also automatically activates the appropriate hyphenation rules for the language you choose.
% You can activate the babel package by adding the next command to the preamble: \usepackage[language]{babel}

%%%%%%%%%%%%%%%%%%%%%%%%%%%%%%%%%%%%%%%%%%%%%%%%%%%%%%%%%%%%%%%%%%%%%%%%%%%%%%%%%%%%%%%%%%%%%%%%%%%%%%%%%%%%%%%%%%%%%%%%%%%%%%%%%%%%%%%%%%%%%%%%%%%%%%%%%%%%%%

\usepackage[T1]{fontenc}

%%%%%%%%%%%%%%%%%%%%%%%%%%%%%%%%%%%%%%%%%%%%%%%%%%%%%%%%%%%%%%%%%%%%%%%%%%%%%%%%%%%%%%%%%%%%%%%%%%%%%%%%%%%%%%%%%%%%%%%%%%%%%%%%%%%%%%%%%%%%%%%%%%%%%%%%%%%%%%

\usepackage[utf8]{inputenc}
% inputenc allows the user to input accented characters directly from the keyboard
% With \usepackage[<encoding>]{inputenc} you can directly input accented and other characters. What's important is that <encoding> matches the encoding with which the file has been written and this depends on your operating system and the settings of your text editor.

% fontenc is oriented to output, that is, what fonts to use for printing characters.
% With \usepackage[T1]{fontenc} you choose an output font encoding that has support for the accented characters used by the most widespread European languages (German, French, Italian, Polish and others), which is important because otherwise TeX would not correctly hyphenate words containing accented letters.

% The two packages are not connected, though it is best to call fontenc first and then inputenc.
% https://tex.stackexchange.com/questions/44694/fontenc-vs-inputenc/44699

%%%%%%%%%%%%%%%%%%%%%%%%%%%%%%%%%%%%%%%%%%%%%%%%%%%%%%%%%%%%%%%%%%%%%%%%%%%%%%%%%%%%%%%%%%%%%%%%%%%%%%%%%%%%%%%%%%%%%%%%%%%%%%%%%%%%%%%%%%%%%%%%%%%%%%%%%%%%%%

% Line numbers


% https://ctan.math.illinois.edu/macros/latex/contrib/vruler/vruler.pdf

%\usepackage{vruler}
%\setvruler
%\setvruler[11pt][2][1][2][0][20pt][20pt][4.2bp][\textheight]
%\setvruler[11pt][1][1][1][0][20pt][20pt][4.2bp][\textheight]

% OR

%\usepackage{tikz,everypage}
%\AtBeginDocument{%
%  \AddEverypageHook{%
%    \begin{tikzpicture}[remember picture,overlay]
%      \path (current page.north west) --  (current page.south west) \foreach \i in {1,...,\fakelinenos} { node [pos={(\i-.5)/\fakelinenos}, xshift=\fakelinenoshift, line number style] {\i} }  ;
%    \end{tikzpicture}%
%  }%
%}
%\tikzset{%
%  line numbers/.store in=\fakelinenos,
%  line numbers=50,
%  line number shift/.store in=\fakelinenoshift,
%  line number shift=5mm,
%  line number style/.style={text=gray},
%}

% OR

%\usepackage{lineno}
% The package extends the lineno package so that even lines within a \footnote are numbered and may be referred to using \linelabel, \ref, etc.
%\linenumbers
% The package lineno should be put before csquotes
% https://tex.stackexchange.com/questions/447006/lineno-package-in-latex-causes-warning-message

%%%%%%%%%%%%%%%%%%%%%%%%%%%%%%%%%%%%%%%%%%%%%%%%%%%%%%%%%%%%%%%%%%%%%%%%%%%%%%%%%%%%%%%%%%%%%%%%%%%%%%%%%%%%%%%%%%%%%%%%%%%%%%%%%%%%%%%%%%%%%%%%%%%%%%%%%%%%%%

\usepackage{csquotes}
% When using babel or polyglossia with biblatex, loading csquotes is recommended to ensure that quoted texts are typeset according to the rules of your main language.

%%%%%%%%%%%%%%%%%%%%%%%%%%%%%%%%%%%%%%%%%%%%%%%%%%%%%%%%%%%%%%%%%%%%%%%%%%%%%%%%%%%%%%%%%%%%%%%%%%%%%%%%%%%%%%%%%%%%%%%%%%%%%%%%%%%%%%%%%%%%%%%%%%%%%%%%%%%%%%

\usepackage[backend=biber, maxbibnames=10, maxcitenames=4, maxalphanames=4, bibencoding=utf8, style=alphabetic, isbn=false, url=false, doi=true, giveninits=true]{biblatex}
% This sets the font size to 10pt with a 10pt baseline skip
\renewcommand*{\bibfont}{\fontsize{10}{10}\selectfont}
\addbibresource{bibliography.bib}
% Package for bibliography
% https://tex.stackexchange.com/questions/25701/bibtex-vs-biber-and-biblatex-vs-natbib
\renewbibmacro{in:}{}
%\renewbibmacro*{in}{}
% Remove the "in(:)" in the bibliography
% print url if no doi
\renewbibmacro*{doi+eprint+url}{%
  \printfield{doi}%
  \newunit\newblock%
  \iftoggle{bbx:eprint}{%
    \usebibmacro{eprint}%
  }{}%
  \newunit\newblock%
  \iffieldundef{doi}{%
  \usebibmacro{url+urldate}}%
  {}%
}
% To print URL only if DOI is not present
% https://tex.stackexchange.com/questions/424774/print-url-only-if-doi-not-present
% To configure biber instead of bibtex, see
% https://tex.stackexchange.com/questions/154751/biblatex-with-biber-configuring-my-editor-to-avoid-undefined-citations
% remove months from all entries
\AtEveryBibitem{\clearfield{month}}
% remove notes from all entries
\AtEveryBibitem{\clearfield{note}}
% How to properly quote arxiv entries
% https://tex.stackexchange.com/questions/436684/correct-arxiv-hyperlinks-with-biblatex-and-mendeley
% @online{borza2025,
%   author      = {Samuël Borza and Wilhelm Klingenberg and Patrick Wood},
%   title       = {Optimal transport on the sub-Lorentzian Heisenberg group},
%   year        = {2025},
%   version     = 1,
%   eprinttype  = {arxiv},
%   eprintclass = {math.MG},
%   eprint      = {2504.03062}
% }

%%%%%%%%%%%%%%%%%%%%%%%%%%%%%%%%%%%%%%%%%%%%%%%%%%%%%%%%%%%%%%%%%%%%%%%%%%%%%%%%%%%%%%%%%%%%%%%%%%%%%%%%%%%%%%%%%%%%%%%%%%%%%%%%%%%%%%%%%%%%%%%%%%%%%%%%%%%%%%

\usepackage{setspace}
% Reduce spacing in bibliography using biblatex
%\singlespacing
%\onehalfspacing
%\doublespacing

%%%%%%%%%%%%%%%%%%%%%%%%%%%%%%%%%%%%%%%%%%%%%%%%%%%%%%%%%%%%%%%%%%%%%%%%%%%%%%%%%%%%%%%%%%%%%%%%%%%%%%%%%%%%%%%%%%%%%%%%%%%%%%%%%%%%%%%%%%%%%%%%%%%%%%%%%%%%%%

%\usepackage{mathptmx}
%\usepackage{mathpazo}
\usepackage{lmodern}
% High quality family of fonts.

%%%%%%%%%%%%%%%%%%%%%%%%%%%%%%%%%%%%%%%%%%%%%%%%%%%%%%%%%%%%%%%%%%%%%%%%%%%%%%%%%%%%%%%%%%%%%%%%%%%%%%%%%%%%%%%%%%%%%%%%%%%%%%%%%%%%%%%%%%%%%%%%%%%%%%%%%%%%%%

\usepackage[indentafter]{titlesec}
\titleformat{name=\section}{}{\thetitle.}{0.8em}{\centering\scshape}
\titleformat{name=\subsection}[runin]{}{\thetitle.}{0.5em}{\bfseries}[.]
\titleformat{name=\subsubsection}[runin]{}{\thetitle.}{0.5em}{\itshape}[.]
% \titleformat{name=\paragraph,numberless}[runin]{}{}{0em}{}[.]
% \titlespacing{\paragraph}{0em}{0em}{0.5em}
\titleformat{name=\subparagraph,numberless}[runin]{}{}{0em}{}[.]
\titlespacing{\subparagraph}{0em}{0em}{0.5em}
% Use the titlesec package to make section titles look like the style of the amsart document class

%%%%%%%%%%%%%%%%%%%%%%%%%%%%%%%%%%%%%%%%%%%%%%%%%%%%%%%%%%%%%%%%%%%%%%%%%%%%%%%%%%%%%%%%%%%%%%%%%%%%%%%%%%%%%%%%%%%%%%%%%%%%%%%%%%%%%%%%%%%%%%%%%%%%%%%%%%%%%%

\usepackage{url}
% The command \url is a form of verbatim command that allows linebreaks at certain characters or combinations of characters, accepts reconfiguration, and can usually be used in the argument to another command. (The \urldef command provides robust commands that serve in cases when \url doesn't work in an argument.) The command is intended for email addresses, hypertext links, directories/paths, etc., which normally have no spaces, so by default the package ignores spaces in its argument.

%%%%%%%%%%%%%%%%%%%%%%%%%%%%%%%%%%%%%%%%%%%%%%%%%%%%%%%%%%%%%%%%%%%%%%%%%%%%%%%%%%%%%%%%%%%%%%%%%%%%%%%%%%%%%%%%%%%%%%%%%%%%%%%%%%%%%%%%%%%%%%%%%%%%%%%%%%%%%%

\usepackage{mathtools}
% Mathtools provides a series of packages designed to enhance the appearance of documents containing a lot of mathematics.
% The mathtools package is an extension package to amsmath.
% Moreover it fixes some deficiencies of the amsmath package and provides additional useful commands such as \coloneqq.
% Amsmath is the principal package in the AMS-LATEX distribution. It adapts for use in LATEX most of the mathematical features found in AMS-TEX; it is highly recommended as an adjunct to serious mathematical typesetting in LATEX.

% Absolute values
% https://tex.stackexchange.com/questions/43008/absolute-value-symbols

\DeclarePairedDelimiter\abs{\lvert}{\rvert}%
\DeclarePairedDelimiter\norm{\lVert}{\rVert}%

% Swap the definition of \abs* and \norm*, so that \abs
% and \norm resizes the size of the brackets, and the 
% starred version does not.
\makeatletter
\let\oldabs\abs
\def\abs{\@ifstar{\oldabs}{\oldabs*}}
%
\let\oldnorm\norm
\def\norm{\@ifstar{\oldnorm}{\oldnorm*}}
\makeatother

%%%%%%%%%%%%%%%%%%%%%%%%%%%%%%%%%%%%%%%%%%%%%%%%%%%%%%%%%%%%%%%%%%%%%%%%%%%%%%%%%%%%%%%%%%%%%%%%%%%%%%%%%%%%%%%%%%%%%%%%%%%%%%%%%%%%%%%%%%%%%%%%%%%%%%%%%%%%%%

\usepackage{amssymb}
% amssymb provides an extended symbol collection. For example, after loading amssymb you have the following additional binary relation symbols: \barwedge, \boxdot, \boxminus, \boxplus, \boxtimes, \Cap, \Cup (and many more), the arrow \leadsto, and some other symbols such as \Box and \Diamond. Another useful feature is the \mathbb command to produce blackboard bold characters.
% Since amssymb internally loads amsfonts, it's enough to load the former.
% amsfonts provides an extended set of fonts for use in mathematics, including: extra mathematical symbols; blackboard bold letters (uppercase only); fraktur letters; subscript sizes of bold math italic and bold Greek letters; subscript sizes of large symbols such as sum and product; added sizes of the Computer Modern small caps font; cyrillic fonts (from the University of Washington); Euler mathematical fonts.

%%%%%%%%%%%%%%%%%%%%%%%%%%%%%%%%%%%%%%%%%%%%%%%%%%%%%%%%%%%%%%%%%%%%%%%%%%%%%%%%%%%%%%%%%%%%%%%%%%%%%%%%%%%%%%%%%%%%%%%%%%%%%%%%%%%%%%%%%%%%%%%%%%%%%%%%%%%%%%

\usepackage{amsthm}
% amsthm helps to define theorem-like structures.
% The package facilitates the kind of theorem setup typically needed in American Mathematical Society publications.

%%%%%%%%%%%%%%%%%%%%%%%%%%%%%%%%%%%%%%%%%%%%%%%%%%%%%%%%%%%%%%%%%%%%%%%%%%%%%%%%%%%%%%%%%%%%%%%%%%%%%%%%%%%%%%%%%%%%%%%%%%%%%%%%%%%%%%%%%%%%%%%%%%%%%%%%%%%%%%

\usepackage{mathrsfs}
% Support use of the Raph Smith’s Formal Script font in mathematics. Provides a \mathscr command, rather than overwriting the standard \mathcal command, as in calrsfs.

%%%%%%%%%%%%%%%%%%%%%%%%%%%%%%%%%%%%%%%%%%%%%%%%%%%%%%%%%%%%%%%%%%%%%%%%%%%%%%%%%%%%%%%%%%%%%%%%%%%%%%%%%%%%%%%%%%%%%%%%%%%%%%%%%%%%%%%%%%%%%%%%%%%%%%%%%%%%%%

\usepackage{dsfont}
% A font based on Computer Modern Roman useful for typesetting the mathematical symbols for the natural numbers (N), whole numbers (Z), rational numbers (Q), real numbers (R) and complex numbers (C); coverage includes all Roman capital letters, ‘1’, ‘h’ and ‘k’. 

%%%%%%%%%%%%%%%%%%%%%%%%%%%%%%%%%%%%%%%%%%%%%%%%%%%%%%%%%%%%%%%%%%%%%%%%%%%%%%%%%%%%%%%%%%%%%%%%%%%%%%%%%%%%%%%%%%%%%%%%%%%%%%%%%%%%%%%%%%%%%%%%%%%%%%%%%%%%%%

\usepackage{enumitem}
% This package provides user control over the layout of the three basic list environments: enumerate, itemize and description. It supersedes both enumerate and mdwlist (providing well-structured replacements for all their funtionality), and in addition provides functions to compute the layout of labels, and to ‘clone’ the standard environments, to create new environments with counters of their own.
% The enumitem package is far more flexible when compared with enumerate. The latter provides an optional argument where you can specify the item number format using a generic representation while the former provides a key-value interface where one can specify number format representation in addition to a host of other list-related settings (or perform these settings globally).
% In short, enumitem provides a modern, user-friendly, key-value interface to many list manipulations, while enumerate is only geared towards short-hand list representation.
% Here you may have to adjust these parameters as per your needs:
% \topsep: space between first item and preceding paragraph.
% \partopsep: extra space added to \topsep when environment starts a new paragraph.
% \itemsep: space between successive items.
% In the above code the parameters are set locally. If you want, you can make them global with the help of \setlist[enumerate]{topsep=0pt,itemsep=-1ex,partopsep=1ex,parsep=1ex}. For details, refer the enumitem manual.
% You can suppress all of these spaces by using nolistsep. We use topsep=-\parskip to counter the space introduced by \usepackage[parfill]{parskip}.
\setlist[enumerate]{noitemsep, partopsep=0pt, topsep=0pt, parsep=0pt, itemsep=0pt}
\setlist[itemize]{noitemsep, partopsep=0pt, topsep=0pt, parsep=0pt, itemsep=0pt}

%%%%%%%%%%%%%%%%%%%%%%%%%%%%%%%%%%%%%%%%%%%%%%%%%%%%%%%%%%%%%%%%%%%%%%%%%%%%%%%%%%%%%%%%%%%%%%%%%%%%%%%%%%%%%%%%%%%%%%%%%%%%%%%%%%%%%%%%%%%%%%%%%%%%%%%%%%%%%%

\usepackage{interval}
% When typing an open interval as $]a,b[$, a closing bracket is being used in place of an opening fence and vice versa. This leads to the wrong spacing in, say, $]-a,b[$ or $A\in]a,b[=B$. The package attempts to solve this using:
% \interval{a}{b} -> [a,b]
% \interval[open]{a}{b} -> ]a,b[
% \interval[open left]{a}{b} -> ]a,b]

% The package also supports fence scaling and ensures that the enclosing fences will end up having the proper closing and opening types. TEX maths does not do this job properly.
\intervalconfig{soft open fences}
% This is just a fast way of saying left open fence=(, right open fence=)

% \ointerval[hoptionsi]{hstarti}{hendi}
% is short for \interval[open,hoptionsi]{hstarti}{hendi}
% \linterval[hoptionsi]{hstarti}{hendi}
% is short for \interval[open left,hoptionsi]{hstarti}{hendi}
% \rinterval[hoptionsi]{hstarti}{hendi}
% is short for \interval[open right,hoptionsi]{hstarti}{hendi}

%%%%%%%%%%%%%%%%%%%%%%%%%%%%%%%%%%%%%%%%%%%%%%%%%%%%%%%%%%%%%%%%%%%%%%%%%%%%%%%%%%%%%%%%%%%%%%%%%%%%%%%%%%%%%%%%%%%%%%%%%%%%%%%%%%%%%%%%%%%%%%%%%%%%%%%%%%%%%%

\usepackage[normalem]{ulem}
% The ulem package provides various types of underlining that can stretch between words and be broken across lines.
% Package ulem redefines \em/\emph. Emphasized text is no longer in italics, but the text is underlined. This can be turned off by option normalem: \usepackage[normalem]{ulem}

%%%%%%%%%%%%%%%%%%%%%%%%%%%%%%%%%%%%%%%%%%%%%%%%%%%%%%%%%%%%%%%%%%%%%%%%%%%%%%%%%%%%%%%%%%%%%%%%%%%%%%%%%%%%%%%%%%%%%%%%%%%%%%%%%%%%%%%%%%%%%%%%%%%%%%%%%%%%%%
% \renewcommand{\baselinestretch}{1.02} 
\usepackage[stretch=10]{microtype}
% Microtype plays with ever-so-slightly shrinking and stretching of the fonts and with the extent to which text protrudes into the margins in a way that yields results that look better, that have fewer instances of hyphenation, and fewer overfull hboxes. It doesn't work with latex, you have to use pdflatex instead.
% Use \usepackage[stretch=10]{microtype} to allow font expansion up to 1% (default is 2%).

%%%%%%%%%%%%%%%%%%%%%%%%%%%%%%%%%%%%%%%%%%%%%%%%%%%%%%%%%%%%%%%%%%%%%%%%%%%%%%%%%%%%%%%%%%%%%%%%%%%%%%%%%%%%%%%%%%%%%%%%%%%%%%%%%%%%%%%%%%%%%%%%%%%%%%%%%%%%%%

\flushbottom
% Makes the last line of a page flush with the bottom margin for neatness
%\raggedbottom

%%%%%%%%%%%%%%%%%%%%%%%%%%%%%%%%%%%%%%%%%%%%%%%%%%%%%%%%%%%%%%%%%%%%%%%%%%%%%%%%%%%%%%%%%%%%%%%%%%%%%%%%%%%%%%%%%%%%%%%%%%%%%%%%%%%%%%%%%%%%%%%%%%%%%%%%%%%%%%

\emergencystretch=1em
% This allows for an additional line-breaking pass with the amount of "tolerable" white space per line increased by 1em.
% Helps avoid overfull hbox (in bibliography for example)

%%%%%%%%%%%%%%%%%%%%%%%%%%%%%%%%%%%%%%%%%%%%%%%%%%%%%%%%%%%%%%%%%%%%%%%%%%%%%%%%%%%%%%%%%%%%%%%%%%%%%%%%%%%%%%%%%%%%%%%%%%%%%%%%%%%%%%%%%%%%%%%%%%%%%%%%%%%%%%

\usepackage[affil-it, noblocks]{authblk}

%\usepackage{etoolbox}
%\makeatletter
%\patchcmd{\@maketitle}
%{\ifx\@empty\@dedicatory}
%{\ifx\@empty\@date \else {\vskip3ex %\centering\footnotesize\@date\par\vskip1ex}\fi
%	\ifx\@empty\@dedicatory}
%{}{}
%\patchcmd{\@maketitle}
%{\ifx\@empty\@date\else \@footnotetext{\@setdate}\fi}
%{}{}{}
%\makeatother
% Patch to make the date appear at the top of the page when using the amsart documentclass, under the author instead of letting amsart place the date in a footnote at the bottom of the first page

%%%%%%%%%%%%%%%%%%%%%%%%%%%%%%%%%%%%%%%%%%%%%%%%%%%%%%%%%%%%%%%%%%%%%%%%%%%%%%%%%%%%%%%%%%%%%%%%%%%%%%%%%%%%%%%%%%%%%%%%%%%%%%%%%%%%%%%%%%%%%%%%%%%%%%%%%%%%%%

\newcommand\numberthis{\addtocounter{equation}{1}\tag{\theequation}}
% In align*, use this to number a particular line
% Rather than using align, and \nonumber-ing every other line

%%%%%%%%%%%%%%%%%%%%%%%%%%%%%%%%%%%%%%%%%%%%%%%%%%%%%%%%%%%%%%%%%%%%%%%%%%%%%%%%%%%%%%%%%%%%%%%%%%%%%%%%%%%%%%%%%%%%%%%%%%%%%%%%%%%%%%%%%%%%%%%%%%%%%%%%%%%%%%

\usepackage{xcolor}

%%%%%%%%%%%%%%%%%%%%%%%%%%%%%%%%%%%%%%%%%%%%%%%%%%%%%%%%%%%%%%%%%%%%%%%%%%%%%%%%%%%%%%%%%%%%%%%%%%%%%%%%%%%%%%%%%%%%%%%%%%%%%%%%%%%%%%%%%%%%%%%%%%%%%%%%%%%%%%

\usepackage[hidelinks,draft=false]{hyperref}
% The hyperref package is used to handle cross-referencing commands in LATEX to produce hypertext links in the document.
% The hyperref documentation says: "Make sure it comes last of your loaded packages". The reason is that it redefines many LaTeX commands. It's a rule of thumb that helps to avoid errors. However, there are exceptions, for instance the amsrefs user's guide notes that amsrefs has to be loaded after hyperref.
% https://tex.stackexchange.com/questions/1863/which-packages-should-be-loaded-after-hyperref-instead-of-before?noredirect=1&lq=1
\hypersetup{
  colorlinks   = true, %Colours links instead of ugly boxes
  urlcolor     = blue, %Colour for external hyperlinks
  linkcolor    = blue, %Colour of internal links
  citecolor   = red %Colour of citations
}

%%%%%%%%%%%%%%%%%%%%%%%%%%%%%%%%%%%%%%%%%%%%%%%%%%%%%%%%%%%%%%%%%%%%%%%%%%%%%%%%%%%%%%%%%%%%%%%%%%%%%%%%%%%%%%%%%%%%%%%%%%%%%%%%%%%%%%%%%%%%%%%%%%%%%%%%%%%%%%

\usepackage[nameinlink,noabbrev,capitalize]{cleveref}
% cleveref must be loaded after hyperref
\crefname{equation}{}{}
%cleveref output (1) instead of eq. (1) when referencing equations only

%%%%%%%%%%%%%%%%%%%%%%%%%%%%%%%%%%%%%%%%%%%%%%%%%%%%%%%%%%%%%%%%%%%%%%%%%%%%%%%%%%%%%%%%%%%%%%%%%%%%%%%%%%%%%%%%%%%%%%%%%%%%%%%%%%%%%%%%%%%%%%%%%%%%%%%%%%%%%%

\usepackage[plain]{fullpage}
% This package is called fullpage, and it sets the text height and width and the margins such that the page is almost full. It only leaves a uniform margin on all four sides.

%%%%%%%%%%%%%%%%%%%%%%%%%%%%%%%%%%%%%%%%%%%%%%%%%%%%%%%%%%%%%%%%%%%%%%%%%%%%%%%%%%%%%%%%%%%%%%%%%%%%%%%%%%%%%%%%%%%%%%%%%%%%%%%%%%%%%%%%%%%%%%%%%%%%%%%%%%%%%%

% https://tex.stackexchange.com/questions/156049/how-can-i-get-cleveref-enumitem-to-print-not-item-e-but-proposition-1
\newlist{theoenum}{enumerate}{1} % also creates a counter called 'propenumi'
\setlist[theoenum]{label=\normalfont(\roman*), ref=\theproposition~\normalfont(\roman*), noitemsep, partopsep=0pt, topsep=0pt, parsep=0pt, itemsep=0pt}
\crefalias{theoenumi}{theorem}

\usepackage{tocloft}
%https://tex.stackexchange.com/questions/142445/how-do-i-center-the-table-of-contents-title-using-tocloft
%centered contents
\renewcommand{\cfttoctitlefont}{\hspace*{\fill}\bfseries\Large}
\renewcommand{\cftaftertoctitle}{\hspace*{\fill}}
% Provides control over the typography of the Table of Contents, List of Figures and List of Tables, and the ability to create new ‘List of ...’
% https://ctan.org/pkg/tocloft?lang=en

%%%%%%%%%%%%%%%%%%%%%%%%%%%%%%%%%%%%%%%%%%%%%%%%%%%%%%%%%%%%%%%%%%%%%%%%%%%%%%%%%%%%%%%%%%%%%%%%%%%%%%%%%%%%%%%%%%%%%%%%%%%%%%%%%%%%%%%%%%%%%%%%%%%%%%%%%%%%%%

%svg command
%/Applications/Inkscape.app/Contents/MacOS/inkscape -D trigpgreater2.svg  -o trigpgreater2.pdf --export-latex

%%%%%%%%%%%%%%%%%%%%%%%%%%%%%%%%%%%%%%%%%%%%%%%%%%%%%%%%%%%%%%%%%%%%%%%%%%%%%%%%%%%%%%%%%%%%%%%%%%%%%%%%%%%%%%%%%%%%%%%%%%%%%%%%%%%%%%%%%%%%%%%%%%%%%%%%%%%%%%

\usepackage{titlefoot}
% Add keywords at the bottom of a document
% https://tex.stackexchange.com/questions/462440/add-keywords-at-the-bottom-of-a-document

%%%%%%%%%%%%%%%%%%%%%%%%%%%%%%%%%%%%%%%%%%%%%%%%%%%%%%%%%%%%%%%%%%%%%%%%%%%%%%%%%%%%%%%%%%%%%%%%%%%%%%%%%%%%%%%%%%%%%%%%%%%%%%%%%%%%%%%%%%%%%%%%%%%%%%%%%%%%%%

% OTHER PACKAGES

%%%%%%%%%%%%%%%%%%%%%%%%%%%%%%%%%%%%%%%%%%%%%%%%%%%%%%%%%%%%%%%%%%%%%%%%%%%%%%%%%%%%%%%%%%%%%%%%%%%%%%%%%%%%%%%%%%%%%%%%%%%%%%%%%%%%%%%%%%%%%%%%%%%%%%%%%%%%%%

% \usepackage[top=2.5cm, bottom=2.5cm, left=2.5cm, right=2.5cm]{geometry}
% The package provides an easy and flexible user interface to customize page layout, implementing auto-centering and auto-balancing mechanisms so that the users have only to give the least description for the page layout. For example, if you want to set each margin 2cm without header space, what you need is just 
% \usepackage[margin=2cm,nohead]{geometry}
% The package knows about all the standard paper sizes, so that the user need not know what the nominal ‘real’ dimensions of the paper are, just its standard name (such as a4, letter, etc.).

%%%%%%%%%%%%%%%%%%%%%%%%%%%%%%%%%%%%%%%%%%%%%%%%%%%%%%%%%%%%%%%%%%%%%%%%%%%%%%%%%%%%%%%%%%%%%%%%%%%%%%%%%%%%%%%%%%%%%%%%%%%%%%%%%%%%%%%%%%%%%%%%%%%%%%%%%%%%%%

% \usepackage{listings}
% The package enables the user to typeset programs (programming code) within LATEX.

%%%%%%%%%%%%%%%%%%%%%%%%%%%%%%%%%%%%%%%%%%%%%%%%%%%%%%%%%%%%%%%%%%%%%%%%%%%%%%%%%%%%%%%%%%%%%%%%%%%%%%%%%%%%%%%%%%%%%%%%%%%%%%%%%%%%%%%%%%%%%%%%%%%%%%%%%%%%%%

% \usepackage{graphicx}
% The package builds upon the graphics package, providing a key-value interface for optional arguments to the \includegraphics command. This interface provides facilities that go far beyond what the graphics package offers on its own.
% To have the images displayed even on draft mode, use \usepackage[final]{graphicx}

%%%%%%%%%%%%%%%%%%%%%%%%%%%%%%%%%%%%%%%%%%%%%%%%%%%%%%%%%%%%%%%%%%%%%%%%%%%%%%%%%%%%%%%%%%%%%%%%%%%%%%%%%%%%%%%%%%%%%%%%%%%%%%%%%%%%%%%%%%%%%%%%%%%%%%%%%%%%%%

% \usepackage{siunitx}
% For typesetting units and numbers

%%%%%%%%%%%%%%%%%%%%%%%%%%%%%%%%%%%%%%%%%%%%%%%%%%%%%%%%%%%%%%%%%%%%%%%%%%%%%%%%%%%%%%%%%%%%%%%%%%%%%%%%%%%%%%%%%%%%%%%%%%%%%%%%%%%%%%%%%%%%%%%%%%%%%%%%%%%%%%

% \usepackage[sc]{mathpazo}
% To use the palatino font.
% Depending on taste, you may want to use [osf] instead of [sc] to get old style numerals as well as the real small caps and better kerning. I for one find old style numerals prettier and classier than lining figures in text mode (using [osf] will keep lining figures in math mode).

%%%%%%%%%%%%%%%%%%%%%%%%%%%%%%%%%%%%%%%%%%%%%%%%%%%%%%%%%%%%%%%%%%%%%%%%%%%%%%%%%%%%%%%%%%%%%%%%%%%%%%%%%%%%%%%%%%%%%%%%%%%%%%%%%%%%%%%%%%%%%%%%%%%%%%%%%%%%%%

% \usepackage{float}
% Improves the interface for defining floating objects such as figures and tables.

%%%%%%%%%%%%%%%%%%%%%%%%%%%%%%%%%%%%%%%%%%%%%%%%%%%%%%%%%%%%%%%%%%%%%%%%%%%%%%%%%%%%%%%%%%%%%%%%%%%%%%%%%%%%%%%%%%%%%%%%%%%%%%%%%%%%%%%%%%%%%%%%%%%%%%%%%%%%%%

% TABLES
% https://tex.stackexchange.com/questions/12672/which-tabular-packages-do-which-tasks-and-which-packages-conflict
% \usepackage{tabular}
% Basic table package.
% \usepackage{booktabs}
% The booktabs package supports professional looking tables; better vertical spacing; better rules; specifically designed for tables without vertical lines (the norm for publication-quality tables).
% The booktabs package creates much nicer looking tables than the standard latex tables.
% \usepackage{cellspace}
% The cellspace package very useful to add, for example, to booktabs.
% This package is intended to allow automatic spacing out of the lines of an array. People often complain about text touching the \hline of a tabular when it is too high or too deep.
%\usepackage{array}
% An extended implementation of the array and tabular environments which extends the options for column formats.
% \usepackage{tabularx}
% The tabularx package provides a column type which expands to fill the specified width of the table.
% \usepackage{tabulary}
% The tabulary package provides column types which are proportional to the natural width of their contents.

%%%%%%%%%%%%%%%%%%%%%%%%%%%%%%%%%%%%%%%%%%%%%%%%%%%%%%%%%%%%%%%%%%%%%%%%%%%%%%%%%%%%%%%%%%%%%%%%%%%%%%%%%%%%%%%%%%%%%%%%%%%%%%%%%%%%%%%%%%%%%%%%%%%%%%%%%%%%%%

%\usepackage{amsaddr}
% This package is intended to be used with the amsart documentclass only. It lets you move the authors’ affiliations either just below the authors’ names on the front page or as footnotes on the first page. The email addresses are always listed as a footnote on the front page.

% To use this package, you need to include the directive \usepackage{amsaddr} in your .tex file. Alternatively, you may also use it with the foot option \usepackage[foot]{amsaddr}. When using the foot option, the adresses are put as footnotes on the first page of the document.

%%%%%%%%%%%%%%%%%%%%%%%%%%%%%%%%%%%%%%%%%%%%%%%%%%%%%%%%%%%%%%%%%%%%%%%%%%%%%%%%%%%%%%%%%%%%%%%%%%%%%%%%%%%%%%%%%%%%%%%%%%%%%%%%%%%%%%%%%%%%%%%%%%%%%%%%%%%%%%

%\usepackage[alphabetic, abbrev]{amsrefs}
% Amsrefs is a LATEX package for bibliographies that provides an archival data format similar to the format of BibTEX database files, but adapted to make direct processing by LATEX easier. The package can be used either in conjunction with BibTEX or as a replacement for BibTEX.

%%%%%%%%%%%%%%%%%%%%%%%%%%%%%%%%%%%%%%%%%%%%%%%%%%%%%%%%%%%%%%%%%%%%%%%%%%%%%%%%%%%%%%%%%%%%%%%%%%%%%%%%%%%%%%%%%%%%%%%%%%%%%%%%%%%%%%%%%%%%%%%%%%%%%%%%%%%%%%

% \usepackage{pgfplots}
% PGFPlots draws high-quality function plots in normal or logarithmic scaling with a user-friendly interface directly in TEX. The user supplies axis labels, legend entries and the plot coordinates for one or more plots and PGFPlots applies axis scaling, computes any logarithms and axis ticks and draws the plots, supporting line plots, scatter plots, piecewise constant plots, bar plots, area plots, mesh-- and surface plots and some more.

%%%%%%%%%%%%%%%%%%%%%%%%%%%%%%%%%%%%%%%%%%%%%%%%%%%%%%%%%%%%%%%%%%%%%%%%%%%%%%%%%%%%%%%%%%%%%%%%%%%%%%%%%%%%%%%%%%%%%%%%%%%%%%%%%%%%%%%%%%%%%%%%%%%%%%%%%%%%%%

% \usepackage{pgfplotstable}
% Pgfplotstable displays numerical tables rounded to desired precision in various display formats (for example scientific format, fixed point format or integer), using TEX’s mathematical facilities for pretty printing. Furthermore, it provides methods for table postprocessing.

%%%%%%%%%%%%%%%%%%%%%%%%%%%%%%%%%%%%%%%%%%%%%%%%%%%%%%%%%%%%%%%%%%%%%%%%%%%%%%%%%%%%%%%%%%%%%%%%%%%%%%%%%%%%%%%%%%%%%%%%%%%%%%%%%%%%%%%%%%%%%%%%%%%%%%%%%%%%%%

% \usepackage{pst-ode}
% The package defines \pstODEsolve for solving initial value problems for sets of Ordinary Differential Equations (ODE) using the Runge-Kutta-Fehlberg (RKF45) method with automatic step size adjustment.

%%%%%%%%%%%%%%%%%%%%%%%%%%%%%%%%%%%%%%%%%%%%%%%%%%%%%%%%%%%%%%%%%%%%%%%%%%%%%%%%%%%%%%%%%%%%%%%%%%%%%%%%%%%%%%%%%%%%%%%%%%%%%%%%%%%%%%%%%%%%%%%%%%%%%%%%%%%%%%

%\usepackage[parfill]{parskip} 
% For no indentation and space between paragraphs.

%%%%%%%%%%%%%%%%%%%%%%%%%%%%%%%%%%%%%%%%%%%%%%%%%%%%%%%%%%%%%%%%%%%%%%%%%%%%%%%%%%%%%%%%%%%%%%%%%%%%%%%%%%%%%%%%%%%%%%%%%%%%%%%%%%%%%%%%%%%%%%%%%%%%%%%%%%%%%%

%\setlength{\parindent}{0pt}
%\setlength{\parskip}{0.5em plus 3pt minus 3pt}
% use these to make paragraphs not indented,
% and to separate consecutive paragraphs

%%%%%%%%%%%%%%%%%%%%%%%%%%%%%%%%%%%%%%%%%%%%%%%%%%%%%%%%%%%%%%%%%%%%%%%%%%%%%%%%%%%%%%%%%%%%%%%%%%%%%%%%%%%%%%%%%%%%%%%%%%%%%%%%%%%%%%%%%%%%%%%%%%%%%%%%%%%%%%

% \usepackage{subcaption}
% \usepackage[margin=15mm,hang]{caption}
% The caption package provides many ways to customise the captions in floating environments like figure and table, and cooperates with many other packages. Facilities include rotating captions, sideways captions, continued captions (for tables or figures that come in several parts).
% \usepackage{subcaption}
% This package supports typesetting of sub-captions (by using the the sub-caption feature of the caption package).
% This package is incompatible with the subfigure and subfig packages.

%%%%%%%%%%%%%%%%%%%%%%%%%%%%%%%%%%%%%%%%%%%%%%%%%%%%%%%%%%%%%%%%%%%%%%%%%%%%%%%%%%%%%%%%%%%%%%%%%%%%%%%%%%%%%%%%%%%%%%%%%%%%%%%%%%%%%%%%%%%%%%%%%%%%%%%%%%%%%%

% \usepackage{imakeidx}
% \makeindex

%%%%%%%%%%%%%%%%%%%%%%%%%%%%%%%%%%%%%%%%%%%%%%%%%%%%%%%%%%%%%%%%%%%%%%%%%%%%%%%%%%%%%%%%%%%%%%%%%%%%%%%%%%%%%%%%%%%%%%%%%%%%%%%%%%%%%%%%%%%%%%%%%%%%%%%%%%%%%%

% \usepackage{bm}
% The bm package defines a command \bm which makes its argument bold. The argument may be any maths object from a single symbol to an expression. This is closely related to the specification of the \boldsymbol command in AMS-LATEX, but \bm is rather more careful in the way it does things.

%%%%%%%%%%%%%%%%%%%%%%%%%%%%%%%%%%%%%%%%%%%%%%%%%%%%%%%%%%%%%%%%%%%%%%%%%%%%%%%%%%%%%%%%%%%%%%%%%%%%%%%%%%%%%%%%%%%%%%%%%%%%%%%%%%%%%%%%%%%%%%%%%%%%%%%%%%%%%%

% \usepackage{showkeys}
% \usepackage[color]{showkeys}
% \definecolor{refkey}{gray}{0.5}
% \definecolor{labelkey}{gray}{0.2}

%%%%%%%%%%%%%%%%%%%%%%%%%%%%%%%%%%%%%%%%%%%%%%%%%%%%%%%%%%%%%%%%%%%%%%%%%%%%%%%%%%%%%%%%%%%%%%%%%%%%%%%%%%%%%%%%%%%%%%%%%%%%%%%%%%%%%%%%%%%%%%%%%%%%%%%%%%%%%%

\pdfsuppresswarningpagegroup=1
% Solve Multiple PDFs with page group included in a single page warning

%%%%%%%%%%%%%%%%%%%%%%%%%%%%%%%%%%%%%%%%%%%%%%%%%%%%%%%%%%%%%%%%%%%%%%%%%%%%%%%%%%%%%%%%%%%%%%%%%%%%%%%%%%%%%%%%%%%%%%%%%%%%%%%%%%%%%%%%%%%%%%%%%%%%%%%%%%%%%%

%\usepackage{xfrac}

% The package allows the user to typeset fractions in the form ‘n/d’.

%%%%%%%%%%%%%%%%%%%%%%%%%%%%%%%%%%%%%%%%%%%%%%%%%%%%%%%%%%%%%%%%%%%%%%%%%%%%%%%%%%%%%%%%%%%%%%%%%%%%%%%%%%%%%%%%%%%%%%%%%%%%%%%%%%%%%%%%%%%%%%%%%%%%%%%%%%%%%%

% \graphicspath{./img}
% The folder in which images are stored for this project.
% If this is enabled, the folder doesn't need to be specified in each
% call to \includegraphics, i.e
%   \includegraphics{picturename}
% rather than
%   \includegraphics{img/picturename}

%%%%%%%%%%%%%%%%%%%%%%%%%%%%%%%%%%%%%%%%%%%%%%%%%%%%%%%%%%%%%%%%%%%%%%%%%%%%%%%%%%%%%%%%%%%%%%%%%%%%%%%%%%%%%%%%%%%%%%%%%%%%%%%%%%%%%%%%%%%%%%%%%%%%%%%%%%%%%%

\usepackage{todonotes}
% The todonotes package allows you to insert to–do items in your document. At any point in the document a list of all the inserted to–do items can be listed with the \listoftodos command.
% obeyDraft option (\usepackage[obeyDraft]{todonotes}) will display todonotes if your document is in draft mode and hide them if it isn't in draft mode. However, be aware that the draft mode also only shows bounding boxes for images, so that might be an issue.
% With the option obeyFinal, todonotes will be disabled only in final mode.


% The todonotes package makes three commands available to the user: \todo[]{}, \missingfigure{} and \listoftodos. \todo[]{} and \missingfigure{} makes it possible to insert notes in your document about things that has to be done later (todonotes . . . ).
%\setuptodonotes{fancyline}
% The fancyline option inserts a curved arrow, pointing from the inserted note to the insertion point.
\setlength{\marginparwidth}{2cm}
% Set the length of the notes in the margin.

%%%%%%%%%%%%%%%%%%%%%%%%%%%%%%%%%%%%%%%%%%%%%%%%%%%%%%%%%%%%%%%%%%%%%%%%%%%%%%%%%%%%%%%%%%%%%%%%%%%%%%%%%%%%%%%%%%%%%%%%%%%%%%%%%%%%%%%%%%%%%%%%%%%%%%%%%%%%%%

% For more packages
% https://tex.stackexchange.com/questions/232007/best-practice-packages-for-a-standard-technical-article-with-math-and-figures
% https://tex.stackexchange.com/questions/553/what-packages-do-people-load-by-default-in-latex?noredirect=1&lq=1

% \usepackage{amsmath,amssymb,mathrsfs}
% \usepackage{amsthm}

% \usepackage{inputenc}
% %\usepackage{bm}
% \usepackage{avant}
% \usepackage[english]{babel}
% \usepackage{thmtools,thm-restate}
% \usepackage[nottoc]
% {tocbibind}
% \usepackage{graphicx}
% \usepackage{enumitem}
% \usepackage{tikz}

% \usepackage{bbm}
% \usepackage{esint}
% \usepackage{authblk}
% \usepackage{mathtools}
% \usepackage{dsfont}
% \usepackage{interval}
% \intervalconfig{soft open fences}

% \usepackage{nicematrix}
% \NiceMatrixOptions{cell-space-limits = 1pt}

% %\addtolength{\hoffset}{-2.2cm}

% %\addtolength{\voffset}{-2cm}
% %\addtolength{\textheight}{5.2cm}
% %\addtolength{\textwidth}{3.8cm}
% \usepackage[inner=2.4cm,outer=2.4cm,top=2.8cm,bottom=2.8cm]{geometry}
% \usepackage[all]{xy}
%  \usepackage{todonotes}
 \newcommand{\red}{\color{red}}
 \newcommand{\blue}{\color{blue}}
 \newcommand{\green}{\color[HTML]{3CB371}}
% \usepackage[bookmarks=true]{hyperref}
% \usepackage{xcolor}
\hypersetup{
    colorlinks,
    linkcolor={red!50!black},
    citecolor={blue!50!black},
    urlcolor={blue!80!black},
}


% \mathtoolsset{showonlyrefs}  % uncomment to tag only the referred equations (post-production)

%SYMBOLS

%standard notations
\DeclareMathOperator{\id}{id}
% \DeclareMathOperator{\rank}{rk}
\DeclareMathOperator{\diam}{diam}
\DeclareMathOperator{\supp}{supp}
\DeclareMathOperator{\area}{Area}
\newcommand{\scal}[2]{\ensuremath{\langle #1 , #2 \rangle}} % scalar product
% \newcommand{\norm}[1]{\left\lVert#1\right\rVert}
\newcommand{\normdot}{{\|\!\cdot\!\|}}
\newcommand*\diff{\mathop{}\!\mathrm{d}}
\newcommand{\normc}[1]{\left\lvert#1\right\rvert}
\newcommand{\normdotc}{{|\!\cdot\!|}}

\newcommand{\modu}[1]{\left |#1\right |}
\newcommand{\Leb}{\mathscr{L}}
\newcommand{\Borel}{\mathscr{B}}
\newcommand{\setR}{\mathbb{R}}
\newcommand{\N}{\mathbb{N}}
\newcommand{\R}{\mathbb{R}}
\newcommand{\Q}{\mathbb{Q}}
\newcommand{\Z}{\mathbb{Z}}
\newcommand{\p}{\mathtt p}
\newcommand{\bbH}{\mathbb{H}}

\newcommand{\prob}{\mathscr P}
\newcommand{\g}{\mathrm{g}}
\newcommand{\vol}{\mathrm{dvol}}
%projection
\newcommand{\de}{\ensuremath{\, \mathrm d}} % in the integrals
\newcommand{\dee}{\ensuremath{\mathrm d}}
\newcommand{\mres}{\mathbin{\vrule height 1.6ex depth 0pt width
0.13ex\vrule height 0.13ex depth 0pt width 1.3ex}}
\newcommand{\suchthat}{\ensuremath{\,:\,}} % such that inside, for example, the sets definition
\newcommand\restr[2]{{% we make the whole thing an ordinary symbol
  \left.\kern-\nulldelimiterspace % automatically resize the bar with \right
  #1 % the function
  \right|_{#2} % this is the delimiter
  }}

% weak convergences
\newcommand{\weakto}{\rightharpoonup}
\newcommand{\Weakto}[1]{{\stackrel
{#1}{\weakto}}}
\newcommand{\weakstarto}{\overset*\rightharpoonup}

%useful symbols
\newcommand{\thmsymbol}{\( \square \)}
\newcommand{\sfD}{\mathsf D}
\newcommand{\Dreg}{\sfD^+_{reg}}

%optimal transport notations
\newcommand{\cd}{\mathsf{CD}}
\newcommand{\T}{\mathrm T}
\newcommand{\bm}{\mathsf{BM}}
\newcommand{\BE}{\mathsf{BE}}
\newcommand{\RCD}{\mathsf{RCD}}
\newcommand{\MCP}{\mathsf{MCP}}
\newcommand{\Ric}{\mathrm{Ric}}
\newcommand{\SC}{\mathsf{SC}}
\newcommand{\Ch}{\mathsf{Ch}}
\newcommand{\X}{\mathsf{X}}
\newcommand{\I}{\mathcal{I}}
\newcommand{\M}{\mathcal{M}}
\newcommand{\J}{\mathcal{J}}
\newcommand{\Ks}{\mathsf{KS}}
\newcommand{\ks}{\mathsf{ks}}
\newcommand{\Lip}{\mathsf {Lip}}
\newcommand{\Lipglob}{\textbf {Lip}}
\newcommand{\Lipfunc}{\text {Lip}}
\DeclareMathOperator{\res}{restr}
\newcommand{\rest}[2]{{{\rm restr}_{#1}^{#2}}}
\newcommand{\di}
{\mathsf d} %standard mms distance notation
\newcommand{\m}{\mathfrak m} %standard mms measure notation
\DeclareMathOperator{\Ent}{Ent}
\DeclareMathOperator{\Geo}{Geo}
\DeclareMathOperator{\OptGeo}{OptGeo}
\DeclareMathOperator{\OptPlans}{OptPlans}
\newcommand{\Prob}{\mathscr{P}}
\newcommand{\ProbTwo}{
\mathscr{P}_2}
\newcommand{\Cost}{\ensuremath{\mathcal C}} % cost of a transport map or a transport plan
\newcommand{\act}{\ensuremath{\mathcal{A}_{2}}} % action of a geodesic on [0, 1]
\newcommand{\action}{\ensuremath{\mathcal{A}}} % Action of a vector field against a probability measure
\newcommand{\ppi}{{\mbox{\boldmath$\pi$}}}
\newcommand{\ssigma}{{\mbox{\boldmath$\sigma$}}}

\newcommand{\dis}{\mathcal D}
\newcommand{\A}{\mathcal A}
\newcommand{\sF}{sub-Fin\-sler }
\newcommand{\sr}{sub-Rie\-man\-nian }
\newcommand{\ariem}{almost-Rie\-man\-nian }
\newcommand{\Ltwo}{L^2\big([0,1];(\R^k,\normdot)\big)}
\DeclareMathOperator{\ann}{Ann}
\newcommand{\sinom}{\sin_\Omega}
\newcommand{\cosom}{\cos_\Omega}
\newcommand{\sinomp}{\sin_{\Omega^\circ}}
\newcommand{\cosomp}{\cos_{\Omega^\circ}}
\newcommand{\Sbb}{\mathbb{S}}
\newcommand{\cut}{{\rm Cut}}
\newcommand{\e}{{\rm e}}
\newcommand{\Mod}[1]{\ (\mathrm{mod}\ #1)}
\newcommand{\hei}{\mathbb{H}}
\renewcommand{\S}{\mathbb{S}}
\newcommand{\Splu}{\mathbb{S}_+}
\renewcommand{\epsilon}{\varepsilon}
\renewcommand{\phi}{\varphi}
\newcommand{\hes}{\mathrm{Hess}}
\newcommand{\G}{\mathbb{G}}
\newcommand{\Scal}{\mathcal{S}}
% \usepackage{titlesec}
% \titleformat{\section}
%   {\scshape\large\centering}
%   {\thesection}{0.5em}{}
% \titleformat{\subsection}
%   {\scshape\centering}
%   {\thesubsection}{0.4em}{}

\date{\today}
\title{\textbf{\uppercase{\large{Curvature-dimension condition of sub-Riemannian $\alpha$-Grushin half-spaces}}}}
\author{Samu\"el Borza\footnote{University of Vienna, Universit\"atsring 1, 1010 Vienna, Austria. \textit{E-mail}:  \href{mailto:samuel.borza@univie.ac.at}{samuel.borza@univie.ac.at}}, \ 
and Kenshiro Tashiro\footnote{Okinawa Institute of Science and Technology, 1919-1 Tancha, Onna-son, Kunigami-gun, Okinawa, Japan. 
\textit{E-mail}:  \href{mailto:kenshiro.tashiro@oist.jp}{kenshiro.tashiro@oist.jp}}}

% \title{\textbf{Curvature dimension condition for $\alpha$-Grushin model spaces}}
% % \title{\textbf{Many norms}}
% \date{}

% % \author{50 Cent\footnote{Institut f\"ur Angewandte Mathematik, Universit\"at Bonn. \textit{E-mail}:  \href{mailto:cent@iam.uni-bonn.de}{cent@iam.uni-bonn.de}}}
% \author{Samu\"el Borza\footnote{Scuola Internazionale Superiore di Studi Avanzati (SISSA, Trieste). \textit{E-mail}:  \href{mailto:sborza@sissa.it}{sborza@sissa.it}}, \
% and Kenshiro Tashiro\footnote{Okinawa Institute of Science and Technology. \textit{E-mail}:  \href{mailto:kenshiro-tashiro@oist.jp}{kenshiro.tashiro@oist.jp}}}


%\author{}
\newtheoremstyle{remark}% <name>
        {10pt}% <Space above>
        {10pt}% <Space below>
        {}% <Body font>
        {}% <Indent amount>
        {\itshape}% <Theorem head font>
        {.}% <Punctuation after theorem head>
        {.4em}% <Space after theorem head>
        {}% <Theorem head spec (can be left empty, meaning 'normal')>

\newtheoremstyle{proof}% <name>
        {10pt}% <Space above>
        {10pt}% <Space below>
        {}% <Body font>
        {}% <Indent amount>
        {\itshape}% <Theorem head font>
        {.}% <Punctuation after theorem head>
        {.4em}% <Space after theorem head>
        {}% <Theorem head spec (can be left empty, meaning 'normal')>

\newtheoremstyle{definition}% <name>
        {10pt}% <Space above>
        {10pt}% <Space below>
        {}% <Body font>
        {}% <Indent amount>
        {\bfseries}% <Theorem head font>
        {.}% <Punctuation after theorem head>
        {.4em}% <Space after theorem head>
        {}% <Theorem head spec (can be left empty, meaning 'normal')>





\newtheoremstyle{theorem}% <name>
        {10pt}% <Space above>
        {10pt}% <Space below>
        {\slshape}% <Body font>
        {}% <Indent amount>
        {\bfseries}% <Theorem head font>
        {.}% <Punctuation after theorem head>
        {.4em}% <Space after theorem head>
        {}% <Theorem head spec (can be left empty, meaning 'normal')>

\theoremstyle{theorem}
\newtheorem{theorem}{Theorem}[section]
\newtheorem{prop}[theorem]{Proposition}
\newtheorem{corollary}[theorem]{Corollary}
\newtheorem{lemma}[theorem]{Lemma}
\newtheorem{conj}[theorem]{Conjecture}

\theoremstyle{definition}
\newtheorem{definition}[theorem]{Definition}

\theoremstyle{remark}
\newtheorem{remark}[theorem]{Remark}
\newtheorem{example}[theorem]{Example}
\newtheorem{notation}[theorem]{Notation}


\theoremstyle{proof}
\newtheorem*{pro}{Proof}
\newenvironment{pr}{\begin{pro}%
 \renewcommand{\qedsymbol}{\thmsymbol}\pushQED{\qed}}%
 {\popQED\end{pro}}

\makeatletter
\renewcommand\xleftrightarrow[2][]{%
  \ext@arrow 9999{\longleftrightarrowfill@}{#1}{#2}}
\newcommand\longleftrightarrowfill@{%
  \arrowfill@\leftarrow\relbar\rightarrow}
\makeatother

\makeindex
\begin{document}

\maketitle

\providecommand{\keywords}[1]
{
	\textbf{\textit{Keywords---}} #1
}

\providecommand{\msc}[1]
{
	\textbf{\textit{MSC (2020)---}} #1
}

\begin{abstract}
	 We provide new examples of sub-Riemannian manifolds with boundary equipped with a smooth measure that satisfy the $\mathsf{RCD}(K , N)$ condition. They are constructed by equipping the half-plane, the hemisphere and the hyperbolic half-plane with a two-dimensional almost-Riemannian structure and a measure that vanishes on their boundary. The construction of these spaces is inspired from the geometry of the $\alpha$-Grushin plane.
\end{abstract}

\keywords{Sub-Riemannian geometry, RCD spaces, optimal transport}

\msc{53C17, 53C21, 49Q22}


% Sam: My idea of structure (suggestion, I am happy to change if you think that something else is better).\\
% \textbf{Section 1} introduction \\
% \textbf{Section 2} The $\alpha$-spaces and halspaces fully described as metric measure space (including the measure). One subsection for each model (three in total).\\
% \textbf{Section 3} Equivalence CD and Ric \\
% \textbf{Section 4} Computation of the Ricci curvatures\\
% What do you think?
% \\
% Ken: I like this structure. But what is the $\alpha$-space?
% In my mind there are three types:\\
% (i)Grushin spaces with $\alpha=(\alpha_1,\dots,\alpha_k)$,\\
%    (ii) Grushin spaces with $\alpha=(\alpha,0,\dots,0)$,\\
%     (iii) Grushin surfaces\\
% If we adopt (iii), then the tentative $\alpha$-Grushin model spaces would be nice.
% (ii), I think we can introduce $n$-dim. model space with a small effort.
% (i), It might be possible, but it is too technical for non-sub-Riemannian researchers. \\
% Sam: Sorry when I said space I just meant the three $\alpha$-models that we are looking at. I am not sure it's worth thinking too much about $n$-dimensional ones. Of course it's interesting and we could do it if you want.\\
% ken: Ok! I'm trying to calculate $n$-dimensional case. Let's decide after the computation is done\\
% Sam: sure! I keep going on
\tableofcontents
% \todo[inline]{Ken: I prefer $\alpha$-Grushin surfaces as a whole name. How do you think?\\ Sam : To me a surface is an $(n - 1)$-dimensional submanifold of an $n$-dim manifold \\ Sam: Note that $\mathsf{RCD}(K, 2)$ can only be satisfied if the measure is proportional to the Riemannian volume.\\
% Ken: Ok, then all ``$\alpha$-Grushin model spaces'' will be changed into ``$\alpha$-Grushin spaces''?}
\section{Introduction}
\label{section:Introduction}

% \todo[inline]{
%     1. Let Riemannian model space $g_K=dx^2 + f_K(x)^2dy^2$,
%     and Grushin metric is $g_\alpha = dx^2 + \frac{f_K(x)^2}{w(x)^{2\alpha}}dy^2$. Add the computation at the begin of section 5, introduce the model in the introduction and unify in section 3 every time we introduce the model.\\
%     33. Check the computation of $\infty$-Grushin (section 5)\\
%     9. Change notation tangent space. Sam: seems DONE by Kenshiro. Sam: I changed more mathrm Ken: I have checked, complete.\\ 12. Change $T$ to $T = 1$. I think also DONE.}
In the past few decades, the curvature-dimension condition $\mathsf{CD}(K, N)$ was specifically developed to generalise to the non-smooth setting of metric measure spaces
% {\red this bound}
the concept of a lower bound on the $N$-Bakry-Émery Ricci tensor $\mathrm{Ric}_{N,V} \geq K$ on weighted Riemannian manifolds, using the theory of optimal transport (see \cite{villani,lott--villani2009, sturm2006}).
The $\RCD(K,N)$ condition is a restriction of the $\cd(K,N)$ condition by requiring in addition that the space is infinitesimally Hilbertian, which excludes Finsler manifolds from being $\RCD$-spaces (see \cite{ambrosioICMs, Gigli12,Erbar2015} for example).
The conditions $\cd(K,N)$ and $\RCD(K,N)$ have not only established known meaningful geometric inequalities to non-smooth spaces, but also helped obtain new results, even in the smooth setting.

However, it has been found that the large class of sub-Riemannian manifolds equipped with a smooth positive measure does not satisfy any of the $\mathsf{CD}$ conditions. In the Heisenberg group, this was first proven in \cite{juillet2005} and then for any Carnot groups in \cite{ambrosio2020}. The same author then showed in \cite{juillet2020} that the same result holds for any sub-Riemannian manifold whose distribution has constant rank strictly smaller than its topological dimension. A no-$\mathsf{CD}$ theorem for two-dimensional rank-varying structures was then found in \cite{magnaboscorossi}. Finally, the most general result to date, established in \cite{rizzi2023failure}, states that no curvature-dimension condition can hold for any sub-Riemannian manifold equipped with a positive smooth measure. As a side note, we refer the reader interested in the study of curvature-dimension in the sub-Finsler setting to the works \cite{borzatashiro}, \cite{magnabosco2023failure}, \cite{borzatashiroSIU2024n1} and \cite{borzatashiroSIU2024n2}.

Surprisingly, it was discovered in \cite{panwei,Pan23} that it suffices to consider a sub-Riemannian structure on a manifold with boundary and equip it with a smooth measure that vanishes on the boundary points to construct an example of a $\mathsf{CD}$-space in sub-Riemannian geometry.
% {\blue Surprisingly, contrary to the expectation from researchers from Riemannian/sub-Riemannian geometry, Pan--Wei constructed an example of $\RCD(0,N)$ space which admits sub-Riemannian structure on its boundary points.}\todo{Let's discuss this sentence}
This example, which we revisit here, is also discussed in \cite{rizzi2023failure}. For $\alpha \geq 0$, the $\alpha$-Grushin plane $\mathbb{G}_\alpha$ is the sub-Riemannian structure on $\R^2$ induced by the vector fields
\[
X = \partial_x,~~~~Y^\alpha = \abs{x}^\alpha \partial_y.
\]
This defines a metric space $(\mathbb{G}_\alpha, \diff_{\mathbb{G}_\alpha})$ that admits the metric tensor
\[g_{\G_\alpha}=\diff x\otimes \diff x+\frac{1}{|x|^{2\alpha}}\diff y\otimes \diff y,\]
at non-singular points.
In \cite{panwei,rizzi2023failure}, the authors equip this metric space with the
the $\beta$-weighted measure $\mathfrak{m}_\beta = \abs{x}^{\beta - \alpha} \diff x \diff y$. Note that this weighted measure vanishes on the singular set $\{x = 0\}$ if $\beta>\alpha$, it coincides with the Lebesgue measure if $\beta=\alpha$,
and it fails to be locally finite if {\blue $\beta\leq \alpha-1$}. Although \cite{panwei} and \cite{rizzi2023failure} examined the same space, their techniques are different. In \cite{panwei}, Pan--Wei constructed a Riemannian manifold $M := [0,+\infty) \times_{f} \S^{n-1} \times_{g} \S^1$ with a doubly warped product metric such that
% {\red the contribution of the metric tensor over $\S^{n-1}$ is huge, and that over $\S^1$ is small for large $r \in [0,+\infty)$}
$f(r)\sim\sqrt{r}$ and $g(r)\sim r^{-2\alpha}$ as the first factor $r$ goes to $+\infty$.
It can be verified that $M$ has positive Ricci curvature.
By taking the asymptotic cone of the universal cover based at a specific point,
it can be argued that its (collapsed) limit is a Ricci limit space satisfying $\RCD(0,n+1)$.
Montgomery pointed out that this limit space is actually the  $\alpha$-Grushin half-plane.


In \cite{rizzi2023failure},
Rizzi--Stefani revisited this example by showing that the half-plane $\overline{\mathbb{G}}_\alpha^+ := \{ x \geq 0\}$ and the open half-plane $\mathbb{G}_\alpha^+ := \{ x > 0\}$ are both geodesically convex subsets of $\mathbb{G}_\alpha$ and that the weighted incomplete Riemannian manifold $(\mathbb{G}_\alpha^+, \diff_{\mathbb{G}_\alpha}, \mathfrak{m}_{\beta})$ satisfies the condition $\mathrm{Ric}_{N,V} \geq K$ if and only if
\begin{equation}
    \label{eq:KNagrushinhalfplane}
    K\leq 0,~~\text{and}~~-\alpha^2-\alpha+\beta\cdot \min\left(\alpha,1-\frac{\beta}{N-2}\right)\geq 0,
\end{equation}
where $V(x,y)=-\beta\log|x|$ is a singular potential.
Actually, \cite{rizzi2023failure} only considered the case $\alpha = 1$ but the general case $\alpha \geq 0$ is obtained in the same way.
%{\red Using a limiting argument, the $K$-convexity of the $N$-Rényi entropy functional on $\overline{\G}_\alpha^+$ can be inherited from $\G^+_\alpha$,  and thus }
They showed that the $\alpha$-Grushin half-plane $\overline{\G}_\alpha^+$ verifies $\cd(K,N)$ for the values of $K \in \R$ and $N \in (2, +\infty]$ such that \cref{eq:KNagrushinhalfplane} holds, using the geodesic convexity and the weighted Ricci curvature lower bound on the interior $\G_\alpha^+$. They also proved that $\overline{\G}_\alpha^+$ is infinitesimally Hilbertian, and thus the $\mathsf{RCD}(K, N)$ is also satisfied.
% {\red The proof of \cite[Theorem 1.11]{rizzi2023failure} which we generalise in \cref{theorem:equivalence} then shows that these values of $N$ and $K$ are exactly the ones for which the $\mathsf{RCD}(K, N)$ condition holds in the sub-Riemannian manifold (with boundary) $(\overline{\mathbb{G}}_\alpha^+, \diff_{\mathbb{G}_\alpha}, \mathfrak{m}_{\beta})$. }

The goal of this paper is to provide a few more examples of this phenomenon. We introduce a sub-Riemannian structure on the hemisphere and on the hyperbolic half-plane that depend on a parameter $\alpha \geq 0$, and by tweaking their unbounded ``Riemannian'' volume measure with a parameter $\beta \geq \alpha$, we obtain metric measure spaces that are $\mathsf{RCD}$-spaces. One motivation for studying these kinds of singular spaces is the singular Weyl's laws studied in \cite{boscain-prandi-seri2016, Chitour-Prandi-Rizzi24, honda2023} and the references therein.

In \cref{section:alphaGrushin}, we define the $\alpha$-Grushin hemisphere and its $\beta$-weighted measure. This metric measure space is an example of sub-Riemannian manifold where the $\mathsf{RCD}(K, N)$ condition holds for some $K > 0$.

\begin{theorem}
\label{thm:main1}
The metric measure space consisting of the sub-Riemannian $\alpha$-Grushin hemisphere equipped with its $\beta$-weighted measure satisfies the following properties.
\begin{enumerate}[label=\normalfont(\roman*)]
    \item For any $\alpha\geq 0$ and $K>0$, there is $\beta\geq\alpha$ such that the $\RCD(K,\infty)$ condition is satisfied.
    \item For any $\alpha\geq 1$ and $N\geq 2+4\alpha(\alpha+1)$, there is $\beta>\alpha$ such that $\RCD(0,N)$ holds.
    \item When $\RCD(0,N)$ is satisfied for some $N \in \linterval{2}{+\infty}$, there is $K>0$ such that $\RCD(K,N)$ holds.
\end{enumerate}
\end{theorem}
Similarly, we introduce in \cref{section:alphaGrushin2} the $\alpha$-Grushin hyperbolic half-plane and its $\beta$-weighted measure. This metric measure space is an example of sub-Riemannian manifold where the $\mathsf{RCD}(K, N)$ condition holds for some $K < 0$.


\begin{theorem}
\label{thm:main2}
    The metric measure space consisting of the sub-Riemannian $\alpha$-Grushin hyperbolic half-plane equipped with its $\beta$-weighted measure satisfies the following properties.
    \begin{enumerate}[label=\normalfont(\roman*)]
    \item For any $\alpha\geq 1$ and $N\geq 2+4\alpha(\alpha+1)$,
    there are $\beta>\alpha$ and $K<0$ such that $\RCD(K,N)$ holds.
    \item For any $\alpha,\beta\geq0$ and $N\in(2,+\infty)$, it does not satisfy $\RCD(0,N)$. Furthermore, it satisfies $\RCD(0,+\infty)$ if and only if $\alpha=1$ and $\beta\geq 2$.
    \item For any $\alpha\geq 1$,
        there are $\beta>\alpha$ and $N\in(-\infty,0)$ such that $\cd(0,N)$ holds.
    % \todo{2. It is possible to define $\RCD(K,N)$ for negative $N$, but since we do not know the equivalence of $\BE$ and $\cd$ for negative $N$, it does not give further information. Ask Chiara.}
\end{enumerate}
\end{theorem}
The negativity of $K$ is essential in the sense that the $\alpha$-Grushin hyperbolic half-plane satisfies the $\cd(0,N)$ condition only if $N$ is negative or $+\infty$.
An odd phenomenon occurs at $\alpha=1$, where the $1$-Grushin hyperbolic half-plane satisfies $\RCD(0,+\infty)$. It is unclear if this phenomenon is relevant to the validity of hyperbolicity or $\mathsf{CAT}(0)$ property in the $1$-Grushin hyperbolic plane.
\begin{remark}
In \cite{Pan23}, Pan constructed a sequence of doubly warped metric spaces which collapses to the $1$-Grushin hemisphere. We do not know if a similar construction holds for other $\alpha$-Grushin spaces, nor if the limit measure in such a construction would be equal to our $\beta$-weighted measure.
\end{remark}

As another generalisation of the $\alpha$-Grushin plane, we introduce a new sub-Riemannian manifold with infinite Hausdorff dimension, which we call the $\infty$-Grushin plane.
The description of this space is found in \cref{section:alphaGrushin3}.
By tweaking the Riemannian volume measure by multiplicative factors that depend on two parameters $\beta$ and $\gamma$,
the $\infty$-Grushin half-plane is shown to satisfy the $\RCD(0,+\infty)$.
% {\red We wished to also give an example of a true $\mathsf{RCD}(0, \infty)$ in sub-Riemannian geometry, and this is realised with the $\alpha$-Grushin half-plane equipped with a weighted measure that depends on two parameters $\beta$ and $\gamma$. The description of this space is found in \cref{section:alphaGrushin3}.}

\begin{theorem}
\label{thm:main3}
   The metric measure space consisting of the sub-Riemannian $\infty$-Grushin half-plane equipped with its $(\beta, \gamma)$-weighted measure satisfies the following properties.
\begin{enumerate}[label=\normalfont(\roman*)]
    \item There are $\beta,\gamma\geq 0$ such that $\RCD(0,+\infty)$ holds.
    \item For any $\beta,\gamma\geq 0$ and $K\in \R$, there is no $N\in(2,+\infty)$ such that $\RCD(K,N)$ holds.
\end{enumerate}
\end{theorem}


% {\blue
% \begin{remark}
%     In \cite{Pan23}, Pan constructed a sequence of doubly warped metric spaces which collapses to $1$-Grushin hemisphere.
%     We do not know if a similar construction holds for other $\alpha$-Grushin model spaces,
%     nor the limit measure in such a construction is equal to our weighted measure.
% \end{remark}}\todo{I don't think this remark is relevant here, maybe in another section}

  \begin{remark}\label{rmk:compact}
    The following consequence is to be noted.
\cref{thm:main3} implies that, for any $N\geq 2$, Gromov--Hausdorff limits of a sequence of metric measure spaces in the class
    \[\mathcal{X}_N:=\left\{(X,\di,\m)\mid \text{proper, }\RCD(0,\infty) \text{ and } \dim_H(X,\di)\leq N\right\}\]
    can have infinite Hausdorff dimension, i.e. $\overline{\mathcal{X}_N}\not\subset \bigcup_{n\in\N}\mathcal{X}_n$.
    Indeed, for $\epsilon>0$, the subsets $\R_{\geq \epsilon}\times \R$ of the $\infty$-Grushin half-plane are weighted Riemannian manifolds with boundary contained in $\mathcal{X}_2$,
    which converges to the $\infty$-Grushin half-plane as $\epsilon\to 0$, which has infinite Hausdorff dimension (see \cref{lemma:infinitehausdorffdimension}).
\end{remark}
%\todo{Sam: I don't like itemize in the remark, it doesn't look good. Before the first bullet was outside, we could go back to that. We could also have two remarks.\\
%K: I put the first remark after hyperbolic. I think this is the correct space. S: PErfect!}

The strategy of the proof of the previous results is fairly simple. After studying these spaces in detail in \cref{section:themodels}, we show in \cref{section:equivalence} that, under the right conditions of smoothness and geodesic convexity (see \cref{theorem:equivalence}), the validity of the $\mathsf{CD}(K, N)$ condition for a metric measure space whose interior is Riemannian is equivalent to the bound $\mathrm{Ric}_{N,V} \geq K$ on its interior. The values of $(K, N)$ for which the $\mathsf{CD}(K, N)$ and $\mathsf{RCD}(K, N)$ conditions holds in the $\alpha$-Grushin half-spaces are then just a matter of Ricci curvature computations, which we provide in \cref{section:riccicomputation}.

\section*{Acknowledgments}
This project has received funding from the European Research Council (ERC) under the European Union’s Horizon 2020 research and innovation program (grant agreement No. 945655).
The authors thank Shouhei Honda for leading the authors' attention to these problems, and telling us the consequences in \cref{rmk:compact}.

\section{Preliminaries}
\subsection{The \texorpdfstring{$\cd(K,N)$}{CD(K, N)} and \texorpdfstring{$\RCD(K,N)$}{RCD(K, N)}  conditions}
Firstly,
we start by recalling the curvature-dimension condition $\cd(K,N)$ and Riemannian curvature-dimension condition $\RCD(K,N)$.
A metric measure space is a triple $(\X,\diff,\m)$ where $(\X,\diff)$ is a complete, separable, locally compact and geodesic metric space, and $\m$ is a non-negative Radon measure on it.
We emphasize that the metric measure spaces will always be assumed to be essentially non-branching.
We denote by $C([0, 1], \X)$ the space of continuous curves from $[0, 1]$ to $\X$, and for $s \in [0, 1]$, we let $e_s : C([0, 1], \X)\ni \gamma\mapsto \gamma(s) \in \X$ be the evaluation map.
A geodesic $\gamma:[0,1]\to \X$ is a curve such that $\diff(\gamma(s), \gamma(t)) = |t - s| \diff(\gamma(0), \gamma(1))$ for all $s, t \in [0, 1]$, and we denote by $\Geo(\X)$ the space of all geodesics on $(\X,\diff)$.
Furthermore, the set of Borel probability measures on $\X$ is denoted by $\Prob(\X)$ and the set of those having finite second momentum by $\Prob_2(\X) \subseteq \Prob(\X)$. We endow the space $\Prob_2(\X)$ with the Wasserstein distance $W_2$, defined by
\begin{equation}\label{wasserstein}
W_2^2(\mu_0, \mu_1) := \inf_{\pi \in \mathsf{Adm}(\mu_0,\mu_1)}  \int \diff^2(x, y) \, \diff \pi(x, y),
\end{equation}
where $\mathsf{Adm}(\mu_0, \mu_1):=\{\pi\in \X\times \X\mid (\p_i)_\sharp \pi = \mu_i,\ \p_i : \text{the projection to the $i$-th factor}~(i=1,2)\}$.
The metric space $(\Prob_2(\X),W_2)$ is itself complete, separable and geodesic.
A probability measure $\pi\in \prob(\X\times \X)$ which attains the minimum values in \cref{wasserstein} is called an optimal transport plan.


\paragraph{The \texorpdfstring{$\cd(K,N)$}{CD(K, N)} condition for positive \texorpdfstring{$N$}{N}.}
In this context, the curvature-dimension condition is defined as follows.
For every $K \in \R$, $N\in [1,\infty)$ and $t\in [0,1]$, the \emph{distortion coefficients} are the functions

% \todo{5. wrong. Sam Check. DONE}
\begin{equation*}
\tau_{K,N}^{(t)}(\theta):=
%
\begin{cases}

%
\displaystyle  +\infty & \textrm{if}\  (N-1)\pi^{2}\leq K\theta^{2} ~\text{ and }~K\theta^2>0,\\
%
\displaystyle t^{1/N}\left(\frac{\sin(t\theta\sqrt{K/N})}{\sin(\theta\sqrt{K/N})}\right)^{1-1/N} & \textrm{if}\  0 < K\theta^{2} < (N-1)\pi^{2},\\
%
t & \textrm{if}\ %K \theta^{2}<0 \ \textrm{and}\ N=0, \ \textrm{or  if}\
K\theta^2 =0, ~~\text{or}~~K\theta^2<0~\text{ and }~N=1, \\
%
\displaystyle  t^{1/N}\left(\frac{\sinh(t\theta\sqrt{-K/N})}{\sinh(\theta\sqrt{-K/N})}\right)^{1-1/N} & \textrm{if}\ K\theta^2 < 0~\text{ and }~N>1.
%
\end{cases}
\end{equation*}

\begin{definition}[{$\cd(K,N)$ condition for $N\in[1,+\infty]$, \cite{sturm2006, lott--villani2009}}]\label{def.2.1}
% \todo{6. SAM $N=1$? $t\in[0,1] $ a.e.? $N=+\infty$ case? DONE.}
For $K\in\R$ and $N\in[1,+\infty]$, a metric measure space $(\X,\di,\m)$ is said to satisfy the $\cd(K,N)$ condition if for every pair of absolutely continuous measures $\mu_0=\rho_0\m,\mu_1= \rho_1 \m \in \Prob_2(\X)$, there exists an optimal transport plan $\pi\in \prob(\X\times \X)$ and an absolutely continuous $W_2$-geodesic $(\rho_t\m)_{t\in [0,1]}$ connecting them such that the following inequality holds for every $N' \geq N$ and every $t \in [0,1]$:
\begin{equation*}\label{eq:CDpositive}
    \int_\X \rho_t^{1-\frac 1{N'}} \diff \m \geq \int_{\X \times \X} \Big[ \tau^{(1-t)}_{K,N'} \big(\diff(x,y) \big) \rho_{0}(x)^{-\frac{1}{N'}} +    \tau^{(t)}_{K,N'} \big(\diff(x,y) \big) \rho_{1}(y)^{-\frac{1}{N'}} \Big]    \diff \pi( x,y),
\end{equation*}
if $N < +\infty$, and, if $N = +\infty$,
  % \[
  %   \mathrm{Ent}(\mu_t | \m) \leq (1 - t) \mathrm{Ent}(\mu_0|\m) + t \mathrm{Ent}(\mu_1 | \m) - \frac{K}{2} t(1 - t) W_2^2(\mu_0, \mu_1).
  % \]
 \[
    \int_\X\rho_t\log\rho_t \diff\m \leq (1 - t)\int_\X\rho_0\log\rho_0 \diff\m + t \int_\X\rho_1\log\rho_1 \diff\m - \frac{K}{2} t(1 - t) W_2^2(\mu_0, \mu_1).\]
\end{definition}
\noindent It is not difficult to see that if a metric measure space satisfies the $\mathsf{CD}(K, N)$ condition for some $K \in \R$ and $N \in [1, +\infty]$, then it also satisfies the $\mathsf{CD}(K', N')$ for any $K' \leq K$ and $N' \in [N,+\infty]$, see \cite[Prop. 1.6]{sturm2006}.

% \todo{Sam: add the $N = \infty$ case\\
% Ken: Is this same?}

\paragraph{The \texorpdfstring{$\cd(K,N)$}{CD(K, N)} condition for negative \texorpdfstring{$N$}{N}.}
The curvature-dimension condition $\cd(K,N)$ for negative $N$ is defined in a similar way.
For $K\in\R$ and $N<0$,
 we set the distortion coefficients
\begin{equation*}
\tilde{\tau}_{K,N}^{(t)}(\theta):=
%
\begin{cases}

%
\displaystyle  +\infty & \textrm{if}\  (N-1)\pi^{2}\geq K\theta^{2},\\
%
\displaystyle t^{1/N}\left(\frac{\sin(t\theta\sqrt{K/N})}{\sin(\theta\sqrt{K/N})}\right)^{1-1/N} & \textrm{if}\  0>K\theta^{2} > (N-1)\pi^{2},\\
%
t & \textrm{if}\ %K \theta^{2}<0 \ \textrm{and}\ N=0, \ \textrm{or  if}\
K\theta^2 =0, \\
%
\displaystyle  t^{1/N}\left(\frac{\sinh(t\theta\sqrt{-K/N})}{\sinh(\theta\sqrt{-K/N})}\right)^{1-1/N} & \textrm{if}\ K\theta^2 > 0.
%
\end{cases}
\end{equation*}


% \todo{8. Do we need negative $N$?}
\noindent The following definition was first introduced in \cite{Ohta-negativeCD}.
\begin{definition}[{$\cd(K,N)$ condition for $N\in(-\infty,0)$}]
% \todo{Sam: A small paragraph about the relevance of CD with $N$ negative + standard properties like consistency + monotonicity. DONE.}
For $K\in\R$ and $N\in(-\infty,0)$,
a metric measure space $(\X,\di,\m)$ is said to satisfy the $\cd(K,N)$ condition if for every pair of absolutely continuous measures $\mu_0=\rho_0\m,\mu_1= \rho_1 \m \in \Prob_2(\X)$, there are an optimal transport plan $\pi\in\prob(\X\times \X)$ and an absolutely continuous a $W_2$-geodesic $(\rho_t\m)_{t\in [0,1]}$ connecting them such that the following inequality holds for every $N'\in[N,0)$ and every $t \in [0,1]$:
\begin{equation*}\label{eq:CDnegative}
    \int_\X \rho_t^{1-\frac 1{N'}} \diff \m \leq \int_{\X \times \X} \Big[ \tilde{\tau}^{(1-t)}_{K,N'} \big(\diff(x,y) \big) \rho_{0}(x)^{-\frac{1}{N'}} +    \tilde{\tau}^{(t)}_{K,N'} \big(\diff(x,y) \big) \rho_{1}(y)^{-\frac{1}{N'}} \Big]    \diff \pi( x,y).
\end{equation*}
\end{definition}
\noindent The $\cd(K,N)$ condition for negative $N$ has been found to appear naturally when studying harmonic measures on the sphere in \cite{Milman2017}, where the author uses previous results from \cite{Milman2017a} to obtain new isoperimetric inequalities for these harmonic measures. This condition is further studied in works such as \cite{Magnabosco2023b, Magnabosco2023a, Magnabosco2023c}. We also have the following consistency property: if a metric measure space satisfies the $\mathsf{CD}(K, N)$ condition for some $K \in \R$ and $N < 0$, then it also satisfies the $\mathsf{CD}(K', N')$ condition for every $K' \leq K$ and $N' \in [N, 0)$. We also have that the $\mathsf{CD}(K, +\infty)$ condition implies the $\mathsf{CD}(K, N)$ condition for any $N < 0$, see \cite[Lemma 2.9]{Ohta-negativeCD}.

\paragraph{The \texorpdfstring{$\RCD(K,N)$}{RCD(K, N)} condition.}
In addition to the $\cd(K,N)$ condition,
we need to introduce infinitesimal Hilbertianity to define $\RCD(K,N)$ spaces, see \cite{Gigli12}.
On a proper metric measure space $(\X,\diff,\m)$,
the Cheeger energy $\mathrm{Ch}:L^2(\X,\m)\to \R$ is defined by
\[\mathrm{Ch}(f):=\inf\left\{\liminf_{n\to \infty}\frac{1}{2}\int_\X\Lip(f_n)\diff \m\mid f_n\in\Lip_b(\X,\diff)\cap L^2(\X,\m),\|f_n-f\|_{L^2}\to 0\right\},\]
where $\Lip_b(\X,\diff)$ is the space of bounded Lipschitz functions and for $f\in\Lip_b(\X,\diff)$ and $x\in \X$,
\[\Lip(f)(x):=\limsup_{y\to x}\frac{|f(y)-f(x)|}{\diff(x,y)}.\]
Then the space of functions $H^{1,2}(\X,\diff,\m):=\{f\in L^2(\X,\m)\mid \mathrm{Ch}(f)<+\infty\}$ defines a Banach space with the norm $\|f\|_{H^{1,2}}:=\left(\|f\|_{L^2}^2+2\mathrm{Ch}(f)^2\right)^{1/2}$.

There are different equivalent definitions of infinitesimal Hilbertianity, and the following is the one we choose here, see \cite{Gigli12} and \cite{Ambrosio2013} for more details.

\begin{definition}[Infinitesimal Hilbertianity and $\RCD(K,N)$ condition]
    A metric measure space $(\X,\diff,\m)$ is said to be infinitesimally Hilbertian if $(H^{1,2}(\X,\diff,\m),\|\cdot\|_{H^{1,2}})$ is a Hilbert space.

    Furthermore, given $K\in\R$ and $N\in[1,+\infty]$, the metric measure space $(\X,\diff,\m)$ satisfies the $\RCD(K,N)$ condition if it verifies the $\cd(K,N)$ condition and is infinitesimally Hilbertian.
\end{definition}

%{\blue }\todo{K: I think we can skip this sentence since we mention it in the next remark. Sam: I think we should keep it otherwise it's two definitions in a row which is pretty ugly.}

%\begin{definition}[{$\RCD(K,N)$ condition}]
%    For $K\in\R$ and $N\in[1,+\infty]$, a metric measure space $(\X,\diff,\m)$ satisfies the $\RCD(K,N)$ condition if it satisfies the $\cd(K,N)$ condition and is infinitesimally Hilbertian.
%\end{definition}

\begin{remark}
    We do not speak about the $\mathsf{RCD}$ condition for $N < 0$ for a few reasons. Firstly, the $\mathsf{CD}$ condition for negative $N$, contrary to when $N \geq 1$, allows for measures $\mathfrak{m}$ that are not locally finite. If that's the case, it is not known if the equivalences of the weak gradients in \cite[Sections 7 and 8]{Ambrosio2013} still hold. Secondly, one of the reasons the $\mathsf{RCD}(K, N)$ condition is  particularly successful is because it is equivalent to the (synthetic) Bakry-\'Emery condition $\mathsf{BE}(K, N)$ under the Sobolev-to-Lipschitz property, see \cite{Erbar2015} and \cite{Ambrosio2015}.
    %{\red  and with the entropic curvature dimension condition $\cd^e(K,N)$,} see \cite{Erbar2015, Ambrosio2015}
    %This equivalence is not yet known for $N < 0$. We refrain from using the $\RCD(K,N)$ condition for negative $N$ because of the lack of a unified treatment in the literature.
    In our specific setting, the measures are always nonnegative, and we only use the equivalence between $\cd(K,N)$ and the smooth Ricci lower bound $\mathrm{Ric}_{N,V} \geq K$.
    % \todo[inline]{Bakry--Emery condition is
    % \[|\nabla P_tu|^2 + \frac{1-e^{-2Kt}}{NK}({\rm L}P_t u)^2\le e^{-2Kt}P_t|\nabla u|^2,\]
    % where $P_t u $ is the solution to the heat equation with the initial value condition $P_0 u = u$. Sam: Ok. I meant the Bakry-Emery tensor.\\
    % K:Then, how about using the notation Ricci lower bound?} This probably mean that any reasonable definition of $\mathsf{RCD}(K, N)$ for $N < 0$ would be applicable in our case.}
    %This probably means that any reasonable definition of $\mathsf{RCD}(K, N)$ for $N < 0$ would be applicable in our case.
    Once the gaps in the $\mathsf{RCD}$ theory for negative effective dimension will have been filled, one should be able to replace $\mathsf{CD}(K, N)$ with $\mathsf{RCD}(K, N)$ when $N < 0$ in this paper, e.g. in (iii) of \cref{thm:main2}
\end{remark}


\subsection{Sub-Riemannian geometry}
\label{subsection:subR}

Before we introduce the $\alpha$-Grushin spaces, we recall some basic facts about sub-Riemannian and metric geometry. For a more comprehensive account of these topics, we refer the reader to \cite{ABB-srgeom} and \cite{burago-burago-sergei2001}, for example.
The $\alpha$-Grushin spaces considered in this work will all be two-dimensional almost-Riemannian manifolds and we recommend especially \cite[Chapter 9]{ABB-srgeom}.

A sub-Riemannian structure on an $n$-dimensional manifold $M$ is given by a set of $m$ globally defined vector field $\mathcal{F} := \{X_1, \dots, X_m\}$, also called the generating frame. The associated distribution is defined as the family, indexed with $x \in M$, of the vector subspaces
\[
\mathcal{D}_x := \mathrm{span} \{X_1(x), \dots, X_m(x)\} \subseteq T_xM.
\]
From this data, it is possible to introduce an inner product $g_x$ on $\mathcal{D}_x$ by applying the polarisation formula to
\begin{equation}
    \label{eq:subRmetric}
    g_x(v, v) := \inf \left\{ \sum_{i = 1}^m u_i^2 \mid \sum_{i = 1}^m u_i X_i(x) = v  \right\}.
\end{equation}

The rank of the sub-Riemannian structure at $x \in M$ is defined by $r(x) := \dim (\mathcal{D}_x)$. A two-dimensional almost-Riemannian manifold is a sub-Riemannian structure on a two-dimensional manifold $M$ such that the cardinal of $\mathcal{F}$ is two.
% \todo{10. DONE. meaning of $|\mathcal{F}|$. Maybe we have to mention a general definition of almost Riemannian manifolds}.
In a two dimensional sub-Riemannian structure, a point $x \in M$ such that $r(x) = 2$ (resp. $r(x) = 1$) is called a Riemannian point (resp. singular point). The singular set, i.e. the set of singular points, must necessarily be small (see \cite[Section 9.1.1]{ABB-srgeom}). In a neighborhood of Riemannian points,
% \todo{11. DONE. In a neighborhood of Riemannian points,? I feel it is a bit dangerous}
the inner product \cref{eq:subRmetric} is a well-defined metric tensor and we can introduce, at those points, a Riemannian volume density (using the same formula as in Riemannian geometry) which diverges when approaching a singular point.

An admissible (or horizontal) curve $\gamma : \interval{0}{1} \to M$ is an absolutely continuous path such that there exists a control $u \in \mathrm{L}^2(\interval{0}{1}, \R^m)$ satisfying
% \todo{12. $T=1$? DONE.}
\[
\dot{\gamma}(t) = \sum_{i = 1}^m u_i(t) X_i(\gamma(t)), \ \text{ for almost every } t \in \interval{0}{1}.
\]
The sub-Riemannian length of an admissible curve $\gamma : \interval{0}{1} \to M$ is then defined by
% \todo{ 13. scalar product was missing DONE}
\[
\mathrm{Length}(\gamma) := \int_0^1 \sqrt{g_{\gamma(t)}(\dot{\gamma}(t), \dot{\gamma}(t))} \mathrm{d}t,
\]
and the sub-Riemannian distance between two points $x, y \in M$ is
\begin{equation}\label{e:distance}
\mathrm{d}(x, y) := \inf \left\{ \mathrm{Length}(\gamma) \mid \gamma \text{ admissible and joins } x \text{ and } y \right\}.
\end{equation}
It is not always given that $\mathrm{d}$ really defines a distance function. If $\mathcal{F}$ satisfies the bracket-generating condition (see \cite[Definition 3.1]{ABB-srgeom}), for instance, then Rashevskii-Chow theorem \cite[Section 3.2]{ABB-srgeom} implies that there exists an admissible curve between every two points of $M$, and that $(M, \diff)$ is a metric space with the metric and manifold topology coinciding.

% Denoting by $\pi: \T^\ast(\mathds{H}) \to \mathds{H}$ is the canonical bundle projection, the Hamiltonian of the sub-Riemannian structure is the map $\mathcal{H} : \T^*(\mathds{H}) \times \mathds{R}^m \times \{0, -1\} \to \mathds{R}$ defined by
% \[
% \mathcal{H}(\lambda, u, \nu) := \sum_{i = 1}^m u_i \langle \lambda_0, X_k(\pi(p)) \rangle + \frac{\nu}{2} \sum_{i = 1}^m u_i^2,
% \]
% while the maximised Hamiltonian is given by
% \[
% H : \mathrm{T^*}M \to \mathds{R} : (p, \lambda_0) \mapsto H(p, \lambda_0) :=  \frac{1}{2}\sum_{k = 1}^m \langle \lambda_0, X_k(p) \rangle^2.
% \]

% The rank $r(x)$ of the sub-Riemannian structure at $x \in M$ is defined by $\dim \mathcal{D}_x$. A two-dimensional almost-Riemannian manifold is a sub-Riemannian structure on a two-dimensional manifold $M$ such that $|\mathcal{F}| = 2$. A point $x \in M$ such that $r(x) = 2$ (resp. $r(x) = 1$) is called a Riemannian point (resp. singular point). The singular set, i.e. the set of singular point, must necessarily be small (see \cite[Section 9.1.1]{ABB-srgeom}). At Riemannian points, we have a well-defined metric tensor and Riemannian volume measure which diverge when approaching a singular point.

Denoting by $\pi: T^\ast M \to M$ is the canonical bundle projection, the Hamiltonian is the map $H : T^*M \to \R$ defined by
\[
H(\lambda) :=  \frac{1}{2}\sum_{k = 1}^m \langle \lambda, X_k(\pi(\lambda)) \rangle^2 ~~~~\forall \lambda\in T_xM
\]
%{\red A (constant speed) length minimiser, or minimising geodesic, is a horizontal curve $\gamma : \interval{0}{1} \to M$ such that $\diff(\gamma(s), \gamma(t)) = |t - s| \diff(\gamma(0), \gamma(1))$ for all $s, t \in [0, 1]$.}
% \todo{15: Geodesics should be constant speed. DONE}
Pontryagin's Maximum Principle is very helpful in the search for
%{\red length-minimisers.}
geodesics.

\begin{theorem}[Pontryagin's Maximum Principle]
\label{theorem:PMP}
     If $\gamma$ is length minimiser parametrised by constant speed, then there exists a Lipschitz curve $\lambda(t) \in T^*_{\gamma(t)}M$ such that one and only one of the following is satisfied:
\begin{enumerate}[label=\normalfont(\roman*)]
\item $\dot{\lambda} = \overrightarrow{H}(\lambda)$, where $\overrightarrow{H}$ is the unique vector field in $T^*M$ such that $\sigma(\cdot, \overrightarrow{H}(\lambda)) = \mathrm{d}_\lambda H$ for all $\lambda \in T^*M$;
\item $\langle \lambda(t), X_i(\gamma(t)) \rangle = 0$ for all $i = 1, \dots, m$, and $\lambda(t) \neq 0$ for all $t \in \interval{0}{1}$.
\end{enumerate}
\end{theorem}
A curve $\lambda : \interval{0}{1} \to T^*M$ satisfying (i) (resp. (ii)) in the theorem above is called a normal (resp. abnormal) extremal. \cref{theorem:PMP} is thus stating that a (constant speed) minimising geodesic has a cotangent lift that is a normal or an abnormal extremal. Note that an extremal in a two-dimensional almost-Riemannian manifold is abnormal if and only if its projection is a constant curve that lies on the singular set (see \cite[
Theorem 9.2]{ABB-srgeom}).
% \todo[inline]{KEn: I'm feeling that the description of Popp's volume makes the confusion.On the regular points, it is just Riemannian volume, and is not defined on singular points. Also I feel that it is too chalenging to talk about Popp's volume on non-eqiuregular manifolds.\\ Sam: let's discuss that}
% {\red On neighborhoods where the rank of the sub-Riemannian structure is constant, it
% is possible to define an intrinsic smooth measure, called Popp’s measure. The general construction is somewhat involved and can be found, for instance, in \cite{popprizzi}. For the purposes of this work, it is sufficient to note that on a 2D almost-Riemannian manifold, the Popp measure coincides with the Riemannian volume in a neighborhood of Riemannian points.
% For the remainder of this work, we will adopt the notation $(\cdot)^{2 \alpha} := ((\cdot)^2)^\alpha$ for every $\alpha \geq 0$.}


% For the remainder of this work, we will adopt the notation $(\cdot)^{2 \alpha} := ((\cdot)^2)^\alpha$ for every $\alpha \geq 0$.

% % {\red A sub-Riemannian structure on an $n$-dimensional manifold $M$ is given by a set of $m$ globally defined vector field $\mathcal{F} := \{X_1, \dots, X_m\}$, also called the generating frame.The associated distribution is defined as the family, indexed with $x \in M$, of the vector subspaces
% % \[
% % \mathcal{D}_x := \mathrm{span} \{X_1(x), \dots, X_m(x)\} \subseteq \mathrm{T}_xM.
% % \]
% % From this data, it is possible to introduce an inner product $g_x$ on $\mathcal{D}_x$ by applying the polarisation formula to
% % \[
% % g_x(v, v) := \inf \left\{ \sum_{i = 1}^m u_i^2 \mid \sum_{i = 1}^m u_i X_i(x) = v  \right\}.
% % \]
% % An admissible (or horizontal) curve $\gamma : \interval{0}{T} \to M$ is an absolutely continuous path such that there exists a control $u \in \mathrm{L}^2(\interval{0}{T}, \R^m)$ satisfying
% % \[
% % \dot{\gamma}(t) = \sum_{i = 1}^m u_i(t) X_i(\gamma(t)), \ \text{ for almost every } t \in \interval{0}{T}.
% % \]
% % The sub-Riemannian length of an admissible curve $\gamma : \interval{0}{T} \to M$ is then defined by
% % \[
% % \mathrm{L}(\gamma) := \int_0^T \sqrt{\dot{\gamma}(t), \dot{\gamma}(t)} \mathrm{d}t,
% % \]
% % and the sub-Riemannian distance between two points $x, y \in M$ is
% % \[
% % \mathrm{d}(x, y) := \inf \left\{ \mathrm{L}(\gamma) \mid \gamma \text{ admissible and joins } x \text{ and } y \right\}.
% % \]}

% % A sub-Riemannian structure on an $n$-dimensional manifold $M$ \todo{{\blue of class  $C^k$}} is given by a pair $(E,f)$, where $E$ is a vector bundle over $M$ equipped with a Euclidean metric $\langle\cdot,\cdot\rangle$,
% % and $f:E\to \mathrm{T}M$ a morphism of vector bundles and is fiber-wise linear \todo{Sam: let's chat about regularity. We need $k \geq 2$ (or 3?) to have smoothness of $H$. Our structures are not $C^\infty$. If we just say ``morphism of vector bundle'' as in the book ABB, we can avoid stating the regularity everywhere Ken: Let's say morphism}. The associated distribution is defined by the image, $\mathcal{D}_x := f(E_x) \subseteq \mathrm{T}_xM$. The bundle projection will be denoted by $\pi_E : E \to M$.
% % From this data, it is possible to introduce an inner product $g_x$ on $\mathcal{D}_x$ by applying the polarisation formula to
% % \begin{equation}
% %     \label{eq:subRmetric}
% %     g_x(v, v) := \inf \left\{ \langle u, u \rangle \mid v=f(u) \right\} , \ \ v \in \mathcal{D}_x .
% % \end{equation}

% % The rank of the sub-Riemannian structure at $x \in M$ is defined by $r(x) := \dim (\mathcal{D}_x)$. A two-dimensional almost-Riemannian manifold is a sub-Riemannian structure $(E, f)$ on a manifold $M$ such that $M$ is two-dimensional and $E$ has bundle rank two. In a two dimensional sub-Riemannian structure, a point $x \in M$ such that $r(x) = 2$ (resp. $r(x) = 1$) is called a Riemannian point (resp. singular point). The singular set, i.e. the set of singular points, must necessarily be small (see \cite[Section 9.1.1]{ABB-srgeom}). At Riemannian points, the inner product \cref{eq:subRmetric} is a well-defined metric tensor and we can introduce, at those points, a Riemannian volume density (using the same formula as in Riemannian geometry) which diverges when approaching a singular point.
% % \todo{Ken: $\mathrm{L}$ is for the length or space of functions?}
%  An admissible (or horizontal) curve $\gamma : \interval{0}{T} \to M$ is an absolutely continuous path such that there exists a control $u \in \mathrm{L}^2(\interval{0}{T}, E)$ with $\pi_E(u(t)) = \gamma(t)$ satisfying
% \[
% \dot{\gamma}(t) = f(u(t)), \ \text{ for almost every } t \in \interval{0}{T}.
% \]

% % {\red An admissible (or horizontal) curve $\gamma : \interval{0}{T} \to M$ is an absolutely continuous path such that there exists a curve $u \in \mathrm{L}^{+\infty}(\interval{0}{T}, \R^m)$ $(m=\mathrm{rank}(\mathrm{U}))$ satisfying
% % \[
% % \dot{\gamma}(t) = f(\gamma(t),u(t)), \ \text{ for almost every } t \in \interval{0}{T}.
% % \]}
% The sub-Riemannian length of an admissible curve $\gamma : \interval{0}{T} \to M$ is then defined by
% \[
% \mathrm{L}(\gamma) := \int_0^T \sqrt{g_{\gamma(t)}(\dot{\gamma}(t), \dot{\gamma}(t))} \ \mathrm{d}t,
% \]
% and the sub-Riemannian distance between two points $x, y \in M$ is
% \[
% \mathrm{d}(x, y) := \inf \left\{ \mathrm{L}(\gamma) \mid \gamma \text{ is admissible and joins } x \text{ and } y \right\}.
% \]
% It is not always given that $\mathrm{d}$ really defines a distance function. If $(E, f)$ satisfies the bracket-generating condition (see \cite[Definition 3.1]{ABB-srgeom}), for instance, then Rashevskii-Chow theorem \cite[Section 3.2]{ABB-srgeom} implies that there exists an admissible curve between every two points of $M$, and that $(M, \diff)$ is a metric space with the metric and manifold topology coinciding.

% Denoting by $\pi: \T^\ast(\mathds{H}) \to \mathds{H}$ is the canonical bundle projection, the Hamiltonian of the sub-Riemannian structure is the map $\mathcal{H} : \T^*(\mathds{H}) \times \mathds{R}^m \times \{0, -1\} \to \mathds{R}$ defined by
% \[
% \mathcal{H}(\lambda, u, \nu) := \sum_{i = 1}^m u_i \langle \lambda_0, X_k(\pi(p)) \rangle + \frac{\nu}{2} \sum_{i = 1}^m u_i^2,
% \]
% while the maximised Hamiltonian is given by
% \[
% H : \mathrm{T^*}M \to \mathds{R} : (p, \lambda_0) \mapsto H(p, \lambda_0) :=  \frac{1}{2}\sum_{k = 1}^m \langle \lambda_0, X_k(p) \rangle^2.
% \]

% {\red The rank $r(x)$ of the sub-Riemannian structure at $x \in M$ is defined by $\dim \mathcal{D}_x$. A two-dimensional almost-Riemannian manifold is a sub-Riemannian structure on a two-dimensional manifold $M$ such that $|\mathcal{F}| = 2$. A point $x \in M$ such that $r(x) = 2$ (resp. $r(x) = 1$) is called a Riemannian point (resp. singular point). The singular set, i.e. the set of singular point, must necessarily be small (see \cite[Section 9.1.1]{ABB-srgeom}). At Riemannian points, we have a well-defined metric tensor and Riemannian volume measure which diverge when approaching a singular point.}

% {\red Denoting by $\pi: \T^\ast(\mathds{H}) \to \mathds{H}$ is the canonical bundle projection, the Hamiltonian is the map $H : \mathrm{T^*}M \to \mathds{R}$ defined by
% \[
% H(\lambda) :=  \frac{1}{2}\sum_{k = 1}^m \langle \lambda_0, X_k(\pi(\lambda)) \rangle^2.
% \]}

% Denoting by $\pi: \T^\astM \to M$ is the cotangent bundle projection, the Hamiltonian is the map $H : \mathrm{T^*}M \to \R$ defined by,
% for $\lambda\in \mathrm{T}^*M$,\todo{Ken: $\langle\cdot,\cdot,\rangle$ is already used \\ Sam: to me that's not a problem, every one knows Ken: it is okay for sub-Riemannian geometers, but it is much kind to change it for other readers}
% \[
%  H(\lambda) :=  \max_{\substack{u \in E \\ \pi_E(u) = \pi(\lambda)}}\left(\langle \lambda, f(u)\rangle-\frac{1}{2}\|u\|^2\right)
% .\]

% A length minimiser, or minimising geodesic, is a horizontal curve $\gamma : \interval{0}{T} \to M$ such that $\diff(\gamma(0), \gamma(T)) = \mathrm{L}(\gamma)$. Pontryagin's Maximum Principle is very helpful in the search for length-minimisers. For this reason, we introduce the Hamitonian $\mathcal{H}^\nu : \T^*M \oplus E \to \mathds{R}$ is defined as
% \[
% \mathcal{H}^\nu(\lambda, u) = \langle \lambda, f(u)\rangle+\frac{1}{2}\|u\|^2,
% \]
% \[
% \mathcal{H} : (\T^*M \oplus U) \times \{0, - 1\} \to \R : (\lambda, u, \nu) \mapsto \langle \lambda, f(u)\rangle+\frac{\nu}{2}\|u\|^2,
% \]
% where $\nu \in \mathds{R}$ and $\pi: \T^\astM \to M$ is the cotangent bundle projection. The Hamiltonian $\mathcal{H}^\nu$ is differentiable in $\lambda$ and the Hamiltonian vector field $\overrightarrow{\mathcal{H}}^\nu(\cdot, u) : \T^*M \to \T^*M$ is then the unique vector field on $\T^*M$ satisfying
% \[
% \sigma(\cdot, \overrightarrow{\mathcal{H}}^\nu(\lambda, u)) = \diff_{\lambda} \mathcal{H}^\nu(\cdot, u),
% \]
% where $\sigma$ denotes the canonical symplectic form of $\T^*M$.\todo{Sam: I think something is wrong here. I am not sure how to to fix $u$ and let $\lambda$ varies. I could make $\lambda$ vary in the same fiber of a fixed $u$ but I don't think that's giving us the right thing.}
% \todo{Ken: Reference. maybe not on ABB?}
% \todo[inline]{Sam: I think that the coordinate-free PMP for vector-bundle formalism is contained in \cite{agrachev2005}. The Hamitonian $\mathcal{H}^\nu : \T^*M \oplus E \to \mathds{R}$ is defined as
% \[
% \mathcal{H}^\nu(\lambda, u) = \langle \lambda, f(u)\rangle+\frac{\nu}{2}\|u\|^2.
% \]
% Let's write $\mathfrak{p} : \T^*M \oplus E \to \T^*M$ the projection on the first component and $\sigma$ the symplectic form on $\T^*M$. Now I am just guessing from \cite{agrachev2005} which only has Hamilton's equation for the normal case:
% }
% \begin{theorem}[{Pontryagin's Maximum Principle}]
%     \label{PMP}
%     Let $\gamma$ is length minimiser parametrised by constant speed controlled by $u \in \mathrm{L}^{+\infty}(\interval{0}{T}, E)$. Then there exists an absolutely continuous curve $\lambda : \interval{0}{1} \to \T^*(\mathds{H})$ satisfying $\pi(\lambda(t)) = \gamma(t)$, and a number $\nu \in \{0, -1\}$, such that for almost all $t \in \interval{0}{1}$
%     \begin{enumerate}[label=\normalfont(\roman*)]
%     \item $\lambda(t) \neq 0$ if $\nu = 0$;
%     \item $\sigma_{}(\diff \mathfrak{p}[\dot{\lambda}(t), \cdot], \diff \mathfrak{p}[\dot{\lambda}(t), \cdot]) = \diff_{(\lambda(t), u(t))} \mathcal{H}^\nu : \T_{\lambda(t)}(\T^*M) \oplus \T_{u(t)}(E) \to \R$;
%     \todo[inline]{Sam: I am too tired and don't understand Agrachev. Maybe we will have a look at page 12 of \cite{agrachev2005}. That's where the equation is written, but I struggle to interpret it.}
%     \item $\displaystyle \mathcal{H}^\nu(\lambda(t), u(t)) = \max_{\substack{u \in E \\ \pi_E(u) = \gamma(t)}} \mathcal{H}^\nu(\lambda(t), u)$.
% \end{enumerate}
% \end{theorem}
% {\red \begin{theorem}[{Pontryagin's Maximum Principle}]
%     \label{PMP}
%     Let $\gamma$ is length minimiser parametrised by constant speed controlled by $u \in \mathrm{L}^{+\infty}(\interval{0}{T}, E)$. Then there exists an absolutely continuous curve $\lambda : \interval{0}{1} \to \T^*(\mathds{H})$ satisfying $\pi(\lambda(t)) = \gamma(t)$, and a number $\nu \in \{0, -1\}$, such that for almost all $t \in \interval{0}{1}$
%     \begin{enumerate}[label=\normalfont(\roman*)]
%     \item $\lambda(t) \neq 0$ if $\nu = 0$;
%     \item $\dot{\lambda}(t) = \overrightarrow{\mathcal{H}}^\nu(\lambda(t), u(t))$;
%     \item $\displaystyle \mathcal{H}^\nu(\lambda(t), u(t)) = \max_{\substack{u \in E \\ \pi_E(u) = \gamma(t)}} \mathcal{H}^\nu(\lambda(t), u)$.
%     \end{enumerate}
% \end{theorem}}

% {\red \begin{theorem}[Pontryagin's Maximum Principle]
% \label{theorem:PMP}
%      If $\gamma$ is length minimiser parametrised by constant speed, then there exists a Lipschitz curve $\lambda(t) \in \mathrm{T}^*_{\gamma(t)}M$ such that one and only one of the following is satisfied:
% \begin{enumerate}[label=\normalfont(\roman*)]
% \item $\dot{\lambda} = \overrightarrow{H}(\lambda)$, where $\overrightarrow{H}$ is the unique vector field in $\mathrm{T}^*M$ such that $\sigma(\cdot, \overrightarrow{H}(\lambda)) = \mathrm{d}_\lambda H$ for all $\lambda \in \mathrm{T}^*M$;
% \item {\red $\mathcal{D}_{\gamma(t)}\subset \ker\lambda$ for all $t\in[0,T]$.}\todo{Is it okay?\\ I don't think that ok. I think we cannot introduce abnormal extremals without the non-maximised hamiltonian}
%for all $i = 1, \dots, m$, and $\lambda(t) \neq 0$ for all $t \in \interval{0}{T}$.
% \end{enumerate}
% \end{theorem}}
% A curve $\lambda : \interval{0}{T} \to \mathrm{T}^*M$ satisfying the conclusions of \cref{PMP} with $\nu = -1$ (resp. $\nu = 0$) is called a normal (resp. abnormal) extremal. \cref{PMP} is thus stating that a (constant speed) minimising geodesic has a cotangent lift that is a normal or an abnormal extremal. Note that an extremal in a two-dimensional almost-Riemannian manifold is abnormal if and only if its projection is a constant curve that lies on the singular set (see \cite[Theorem 9.2]{ABB-srgeom}).
% \todo[inline]{KEn: I'm feeling that the description of Popp's volume makes the confusion.On the regular points, it is just Riemannian volume, and is not defined on singular points. Also I feel that it is too chalenging to talk about Popp's volume on non-eqiuregular manifolds.\\ Sam: let's discuss that}
% {\red On neighborhoods where the rank of the sub-Riemannian structure is constant, it
% is possible to define an intrinsic smooth measure, called Popp’s measure. The general construction is somewhat involved and can be found, for instance, in \cite{popprizzi}. For the purposes of this work, it is sufficient to note that on a 2D almost-Riemannian manifold, the Popp measure coincides with the Riemannian volume in a neighborhood of Riemannian points.
% For the remainder of this work, we will adopt the notation $(\cdot)^{2 \alpha} := ((\cdot)^2)^\alpha$ for every $\alpha \geq 0$.}

% \begin{remark}
%     Given $m$ global vector fields $X_1, \dots, X_m : M \to \mathrm{T}M$ on a manifold $M$, we can induce on $M$ a sub-Riemannian structure in the following way. We set $E = M \times \mathds{R}^m$ the trivial bundle of rank $m$, $f : E \to \mathrm{T}M : (p, (u_1, \dots, u_m)) \mapsto \sum_{k = 1}^{m} u_k X_k(p)$ and finally the metric on $E$ is the Euclidean one. In this way, we induce an inner product on $\mathcal{D}_p = \mathrm{span}\{X_1(p), ..., X_m(p)\}$ by the polarisation formula applied to the norm
% \begin{equation*}
% 	\label{eq:freeinnerproduct}
% 	\|u\|^2_{\mathcal{D}_p} := \min \left\{ \sum_{k = 1}^{m} u^2_i \mid \sum_{k = 1}^{m} u_i X_k(p) = u \right\}.
% \end{equation*}
% The family $(X_1, \dots, X_m)$ is said to be a generating family of the sub-Riemannian structure $(E, f)$ on $M$. A \textit{free} sub-Riemannian structure is one that is induced from a generating family. Every sub-Riemannian structure is \textit{equivalent} to a free one (see \cite[Section 3.1.4]{ABB-srgeom}).
% {\blue For a manifold $M$ equipped with a free sub-Riemannian structure,
% the maximised Hamiltonian $H:\mathrm{T}^*M\to \R$ is deduced to the so called sub-Riemannian Hamiltonian
% \[H(\lambda)=\frac{1}{2}\sum_{i=1}^m\langle \lambda,X_i(\pi(\lambda))\rangle^2,\]
% and a normal extremal $\lambda(t)=(x(t),p(t))$ is characterised by the solution to the Hamilton's equation
% \[\dot{x}=\frac{\partial H}{\partial p},~~\dot{p}=-\frac{\partial H}{\partial x}\]}
% \end{remark}

For the remainder of this work, we will adopt the notation $(x)^{2 \alpha} := (x^2)^\alpha  \in \R_{\geq 0}$ for every $\alpha \geq 0$ and $x\in\R$.

\section{Geometry of \texorpdfstring{$\alpha$}{alpha}-Grushin half-spaces}
\label{section:themodels}

\subsection{The \texorpdfstring{$\alpha$}{alpha}-Grushin sphere and hemisphere}
\label{section:alphaGrushin}

On the two-dimensional Riemannian sphere $(\S^2, g_{\S^2})$, we introduce the following coordinate chart, the validity of which can be found in \cite{boscain-prandi-seri2016}.
% \todo{16. Since the equator is vertical, $N$ and $S$ should be $E$ and $W$, the east and west pole? Answer: $N, S$ do not belong to $\gamma$.}
Fix a large circle $\gamma:\R/2\pi\Z\to \S^2$,
and let $N$ and $S$ be the north pole and south pole of $\S^2$ respectively, with respect to $\gamma$. For $p\in \S^2\setminus\{N,S\}$,
we define $x := x(p) \in \ointerval{-\pi/2}{\pi/2}$ as the signed (spherical) distance $\diff_{\S^2}(p,\mathrm{Im}(\gamma))$,
where the sign is positive (resp. negative) if $p$ belongs to the hemisphere containing $N$ (resp. $S$).
Furthermore,
define the number $y=y(p)\in \R/2\pi\Z$ so that $\gamma(y)$ is the perpendicular foot from $p$ to $\mathrm{Im}(\gamma)$. The map $\S^2 \to \ointerval{-\pi/2}{\pi/2}\times \R/2\pi\Z  : p \mapsto (x(p), y(p))$ is a well-defined coordinate chart which can be naturally extended to $N$ and $S$ under the identification
$(\frac{\pi}{2},y_1) \sim (\frac{\pi}{2},y_2)$ and $(-\frac{\pi}{2},y_1) \sim (-\frac{\pi}{2},y_2)$ for any $y_1,y_2\in \R/2\pi\Z$. Note that $N$ (resp. $S$) corresponds to the equivalence class of $(\frac{\pi}{2},y)$ (resp. $(-\frac{\pi}{2},y)$). It is easily shown that in these coordinates, the spherical Riemannian metric tensor $g_{\S^2}$ possesses the warped product structure
\[
g_{\S^2} = \diff x \otimes \diff x + \cos^2(x) \diff y \otimes \diff y.
\]
\begin{remark}
% \todo{17. Do we use spherical coordinates? We choose it like this because it is exactly the same construction for the hyperbolic case, and justifies it. DONE.}
The coordinate system $(x, y)$ is, up to rotations, the same as the spherical coordinates $(\phi,\theta) \in  (-\pi/2, \pi/2)\times \R / 2\pi\Z$
% \todo{20. This should be open? DONE.}
which parametrises the standard sphere $\S^2 \setminus \{N, S\}\subseteq \R^3$ by
\[(\theta,\phi)\mapsto (\cos(\theta)\cos(\phi),
\sin(\theta)\cos(\phi),\sin(\phi)).\]
\end{remark}
With this in mind, we can introduce the $\alpha$-Grushin sphere and hemisphere.
\begin{definition}
% \todo{21. We should say somewhere that $\alpha$ can be non-integer, since the PMP holds, but it is no longer almost Riemannian manifolds}
\label{def:alphasphere}
    For $\alpha \geq 0$, the $\alpha$-Grushin sphere $\S_\alpha$ is the sub-Riemannian structure on $\S^2$ induced from the vector field $X$ and $Y^\alpha$ given by
\[
X:=\partial_{x},~~Y^\alpha:=\frac{\abs{\sin(x)}^{\alpha}}{\cos(x)}\partial_{y}.
\]
The $\alpha$-Grushin hemisphere $\overline{\S}^+_\alpha$ (resp. open hemisphere $\S^+_\alpha$) is the subset of $\S_\alpha$ defined by
\[
\overline{\S}^+_\alpha := \{ p \in \S^2 \mid x \in [0,\pi/2] \} \text{ (resp. $\S^+_\alpha := \{ p \in \S^2 \mid x \in (0,\pi/2] \})$}.
\]

\begin{remark}
    When $\alpha$ is non-integer, the vector fields are not smooth and they are not bracket-generating. However, any pair of points can still be joined with a horizontal curve and Pontryagin's Maximum Principle stated in \cref{theorem:PMP} can be applied. Therefore, the Carnot-Carathéodory metric \cref{e:distance} can be constructed as in smooth sub-Riemannian geometry. This remark remains valid for the other model spaces introduced in this work.
\end{remark}

\begin{remark}
     Note that, strictly speaking, the structure introduced in \cref{def:alphasphere} does not fall into the definition of sub-Riemannian structure laid out in \cref{subsection:subR}. Indeed, the vector fields $X$ and $Y^\alpha$ are not global vector fields: they are not defined at the poles, i.e. at $x=\pm\pi/2$. To be completely rigorous, one should therefore check that it is still a sub-Riemannian manifold but according to the general definition of \cite[Definition 3.2]{ABB-srgeom}. There, a sub-Riemannian manifold is given by couple $(E, f)$ where $E$ is a Euclidean vector bundle and $f : E \to T M$ is a morphism of vector bundles.
     In our setting, $E = T\Sbb^2$ and $f$ is a morphism that we now construct explicitly.

    Denote by $p:T\Sbb^2\to \Sbb^2$ the bundle projection. Letting $\mathcal{U}_1 := \mathbb{S}^2 \setminus \{N, S\}$, the coordinate chart
    $\phi_1 := (x, y) : \mathcal{U}_1 \to \mathbb{R}^2$
    induces a chart on $T\mathbb{S}^2$, and we define
    $f_{1} : p^{-1}(\mathcal{U}_1) \to T \mathbb{S}^2$
    as the map, linear on fibers, that satisfies
    \[
    f_{1}(\partial_x) = \partial_x, \qquad f_{1}(\partial_y) = \abs{\sin(x)}^{\alpha} \partial_y.
    \]
    Letting $\mathcal{U}_2 := \{ (s, t, z) \in \mathbb{S}^2\subset \R^3 \mid z > 0 \}$ the map
    \[
    \phi_2 : \mathcal{U}_2 \to \mathbb{R}^2 :  (s, t, \sqrt{1 - s^2 - t^2}) \mapsto (s, t)
    \]
    is another coordinate chart on $\mathbb{S}^2$, which also induces a chart on $T \mathbb{S}^2$. We set $f_{2} : p^{-1}(\mathcal{U}_2) \to T \mathbb{S}^2$
    the map, linear on fibers, satisfying
    \[
    f_2(\partial_s) = \frac{1}{s^2+t^2} \left[ \left(s^2+(1-s^2-t^2)^{\alpha/2}t^2\right)  \partial_s + st(1-(1-s^2-t^2)^{\alpha/2}) \partial_t \right]
    \]
    and
    \[
    f_2(\partial_t) = \frac{1}{s^2+t^2} \left[ st(1-(1-s^2-t^2)^{\alpha/2}) \partial_s + \left((1-s^2-t^2)^{\alpha/2}s^2 + t^2 \right) \partial_t \right].
    \]
    This ensures $f_1=f_2$ on $p^{-1}(\mathcal{U}_1\cap\mathcal{U}_2)$ via the coordinate transformation $s=\cos(x)\cos(y)$, $t=\cos(x)\sin(y)$. The apparent singularity at $(s,t)=(0,0)$ is removable since
\[
f_2(\partial_s)\to\partial_s,
\quad
f_2(\partial_t)\to\partial_t
\quad\text{as }(s,t)\to(0,0),
\]
so $f_2$ extends smoothly over the origin, giving the identity on the fiber at $N$. The maps $f_1$ and $f_2$ are patched together to define a smooth, globally defined bundle morphism \(f\colon E\to T\Sbb^2\), and the pair $(E,f)$ defines a sub‐Riemannian manifold in the sense of \cite[Definition 3.2]{ABB-srgeom}.


%     {\red and the generating family is the pushforward by $f$ of an orthonormal frame in $E$.
%     In our setting, $E$ is the tangent bundle $T\Sbb^2$ with its (local) orthonormal frame $\{\partial_x,\frac{1}{\cos(x)}\partial_y\}$, and the morphism $f$ is defined on $T\Sbb^2\setminus\{T_N\Sbb^2\cup T_S\Sbb^2\}$ by
%     \[f_{(x,y)}=\begin{pmatrix}
%         1 & 0\\
%         0 & |\sin(x)|^{\alpha}
%     \end{pmatrix}.
%     \]
%     % f(u\partial_x+v\frac{1}{\cos(x)}\partial_y)=u\partial_x + |x|^\alpha v\frac{1}{\cos(x)}\partial_y.\]
% The morphism $f$ can then be smoothly\todo{how smooth} extended to $N$ and $S$ by transforming the coordinates $(x,y)$ into $(s,t)=(\cos(x)\cos(y), \cos(x)\sin(y))$:
%     \[f_{(s,t)}=\frac{1}{s^2+t^2}\begin{pmatrix}
%         s^2+(1-s^2-t^2)^{\alpha/2}t^2 & st(1-(1-s^2-t^2)^{\alpha/2})\\
%         st(1-(1-s^2-t^2)^{\alpha/2}) & (1-s^2-t^2)^{\alpha/2}s^2 + t^2
%     \end{pmatrix}.\]
%     }
    %{\red  Our simpler definition is called a \textit{free} sub-Riemannian structure and, ultimately, every sub-Riemannian structure is equivalent to a free one (see \cite[Section 3.1.4]{ABB-srgeom}). In our case, the bundle $E$ is just $T \mathbb{S}^2$
    % \todo{22. $\Sbb_2\to \Sbb^2$: DONE}
    %while we have defined the morphism $f$ outside the poles. $f$ can then be smoothly extended to $N$ by taking another chart, for example transforming the coordinates $(x,y)$ into the canonical upper hemisphere coordinates $(s,t)\mapsto (s,t,\sqrt{1-s^2-t^2})$.
    % \todo{23. something is wrong in this arguments? Explain why this is right in the response to reviewer.}
    %It can be verified that everything stays smooth and that the morphism $f$ becomes the identity map at the poles, i.e. at $(s,t)=(0,0)$. }
    % \todo[inline]{Sam: add comment comment 21 $\alpha$ non integer.}
\end{remark}

% {\blue At the point $x=\pm\frac{\pi}{2}$, we cannot extend $X$ and $Y^\alpha$, however, the Riemannian metric is smoothly extended. Indeed, after transforming the coordinates $(x,y)$ into the canonical upper hemisphere coordinates $(s,t)\mapsto (s,t,\sqrt{1-s^2-t^2})$, we can  verify the smoothness at $(s,t)=(0,0)$.}
\end{definition}
The $0$-Grushin sphere is simply the two-dimensional Riemannian sphere $\S^2$. When $\alpha > 0$, the $\alpha$-Grushin sphere is a two-dimensional almost-Riemannian structure with $\{x = 0\}$ being its set of singular points. The Grushin sphere studied in \cite{boscain-prandi-seri2016,Pan23} corresponds to the $1$-Grushin sphere. At non-singular points, this sub-Riemannian structure admits the Riemannian metric
\begin{equation}
    \label{eq:SalphaRiemannianMetric}
    g_{\S_{\alpha}} = \diff x \otimes \diff x + \frac{\cos^2(x)}{\sin^{2\alpha}(x)} \diff y \otimes \diff y.
\end{equation}
A simple computation shows that the
% \todo{24. unweighted should be removed? DONE}
 Riemannian volume induced from \cref{eq:SalphaRiemannianMetric}, is given by
\[
\vol_{\S_\alpha}=\cos(x)\abs{\sin(x)}^{-\alpha}\diff x\diff y.
\]
% \todo{25. Should we use $dx\wedge dy?$ No because we say measure (=density) and not form, so no negative sign which there could be with wedge. DONE.}
We introduce the following weighted measure.
\begin{definition}
\label{def:MeasurealphaGrushinSphere}
For $\beta\geq \alpha$,
we consider the weighted measure given by
\begin{equation*}
    \label{eq:MeasurealphaGrushinSphere}
    \m^\beta_{\S_\alpha}:=\abs{\sin(x)}^\beta\vol_{\S_\alpha}=\cos(x)\abs{\sin(x)}^{\beta-\alpha}\diff x \diff y=e^{-V_{\S_\alpha}}\vol_{\S_\alpha},
\end{equation*}
where $V_{\S_\alpha}(x,y):=-\beta\log\abs{\sin(x)}$.
\end{definition}

% \todo[inline]{Sam: strictly speaking, is $(\S_\alpha, \diff_{\S_\alpha}, \mathfrak{m}^\beta_{\S_\alpha})$ a metric measure space even for $\beta < \alpha$ or only when $\beta \geq \alpha$?\\
% Ken: You are right, if $\beta-\alpha\in(-\infty,-1]$, then the measure is not locally finite. If $\beta-\alpha\in(-1,0)$, then it is locally finite (and Radon?) but possibly some regularity might fail. It would be safe to assume $\beta\geq \alpha$.}
\noindent The $\alpha$-Grushin hemisphere is a geodesically convex subset of $\S_\alpha$ and a
% {\red uniquely}
geodesic space when seen as a length subspace of $\S_\alpha$. This is made clear by the next result.

\begin{prop}
    \label{prop:geodesicspacealphaGrushinSphere}
     There is a
     % {\red unique}\todo{Sam: do we need to prove uniqueness of geodesics for our result? Ken: I think no}
     minimising geodesic contained within $\overline{\S}^+_\alpha$ that joins any two given points in $\overline{\S}^+_\alpha$. Furthermore, the $\alpha$-Grushin open hemisphere $\S_\alpha^+$ is a geodesically convex subset of $\overline{\S}^+_\alpha$ and has the structure of an incomplete weighted Riemannian manifold when equipped with the restriction of the measure $\mathfrak{m}^\beta_{\S_\alpha}$.
    % {\red In addition,
    % $\gamma$ is written by the image of the exponential map defined in \ref{expo}. Sam: Maybe we do not really use this information anymore?\\
    % Ken: Agree, maybe not. What we need is (strongly) geodesically convexity of $M$, and it can be proved without exponential map.}
\end{prop}

\begin{proof}

% {\red Note that if a path on $\Sbb^2$ is contained in both of $\Splu$ and $\S_-$,
% then by the reflection along the singular equator $\{x=0\}$,
% we can convert the path on $\S^2$ into that on $\Splu$ with the same length.
% Moreover,
% if a path has a corner on the singular equator,
% then a visual shortcut yields a truly shorter path.
% Therefore every length minimizer touches the singular equator at most once,
% and it fails to minimize after it passes through the equator.
% It implies that every length minimizer in $(\Splu,g_\alpha)$ is the restriction of one on $(\S^2,g_\alpha)$.

% Next we compute length minimizers on $(\S^2,g_\alpha)$ restricted to $\Splu$.
% Note that we always have $\sin(x),\cos(x)\geq 0$.
% Let $H:T^\ast \Splu$ be the sub-Riemannian Hamiltonian,
% that is,
% for the coordinates $\lambda=(x,y,u,v)\in T^\ast \Splu$,
% \[H(\lambda):=\frac{1}{2}\left[\langle \lambda,X\rangle^2+\langle\lambda,Y\rangle^2\right]=\frac{1}{2}\left[u^2+\frac{\sin^{2\alpha}(x)}{\cos^2(x)}v^2\right].\]
% Since we have no abnormal minimizers,
% every length minimizer $\gamma(t)=(x(t),y(t))$ is the solution to the Hamilton's equation
%     \[\begin{cases}
%     \dot{x}=u,\\
%     \dot{y}=\frac{\sin^{2\alpha}(x)}{\cos^2(x)}v,\\
%     \dot{u}=-[2\alpha\cos^2(x)+2\sin^2(x)]\cdot \frac{\sin^{2\alpha-1}(x)}{\cos^3(x)}\cdot v^2,\\
%     \dot{v}=0.
% \end{cases}\]
% From the equation,
% $v\equiv v_0$ is a constant function.
% If $v\equiv 0$,
% then $\gamma(t)=(x(t),y(t))=(x_0+u_0t,y_0)$ i.e. the great circle through the north pole.
% If $v_0\not\equiv 0$,
% then $y(t)$ is strictly increasing/decreasing,
% and $x(t)$ is concave.
% Since every length minimizer is the solution to the Hamilton's equation \ref{eq:Hamilton},
% the above monotonicity/convexity arguments imply that the length minimizer is uniquely deteremined.

% Therefore,
% for the initial point $z=(x_0,y_0)$,
% we can define the exponential map $\exp_{z}^t(u,v)=(x(t),y(t))$ defined on an appropriate open subset $O\subset T_{z}^\ast\Splu$\footnote{If $z$ is on the singular equator,
% the tangent space $T_z^\ast \Splu$ is a closed half plane $\R_{\geq 0}\times \R$} such that $\exp_z^1$ is diffeomorphism onto its image.}
% \todo[inline]{Sam: I try to reorganise the proof. Let me know what you think.}
% {\blue

We start by noting that between every two points on the $\alpha$-Grushin sphere, there is indeed a horizontal path controlled by the vector fields $X$ and $Y^{\alpha}$ joining them. This means that the induced sub-Riemannian distance $\diff_{\S_\alpha}$ is well-defined and that $(\S_\alpha, \diff_{\S_\alpha})$ is a locally compact metric space. Even though the bracket generating condition is not verified when $\alpha \notin \N$, it is easy to see that the metric topology still coincides with the original topology of $\S^2$ by the monotonicity property of $\mathrm{d}_{\S_\alpha}$ with respect to $\alpha \geq 0$. Furthermore, any metric ball is compact and thus the metric space $(\S_\alpha, \diff_{\S_\alpha})$ is complete.

The sub-Riemannian Hamiltonian $H : T^*(\S_\alpha) \to \R$ can be written in the canonical coordinates $(x, y, u,v )$ induced from $(x, y)$  as
\[
H(\lambda):=\frac{1}{2}\left[\langle \lambda,X\rangle^2+\langle\lambda,Y^\alpha\rangle^2\right]=\frac{1}{2}\left[u^2+\frac{\sin^{2\alpha}(x)}{\cos^2(x)}v^2\right].
\]
A normal extremal $\lambda : \interval{0}{T} \to T^*(\S_\alpha) : t \mapsto (x(t), y(t), u(t), v(t))$ satisfies the following Hamiltonian system of equations
\begin{equation}
    \label{eq:alphaGrushinsphereHamilton}
    \begin{cases}
    \dot{x}=u,\\
    \displaystyle \dot{y}=\frac{\sin^{2\alpha}(x)}{\cos^2(x)}v,\\
    \displaystyle \dot{u}= -  v^2 \  \sin^{2(\alpha - 1)}(x) \tan(x) (\alpha + \tan^2(x)),\\
    \dot{v}=0.
\end{cases}
\end{equation}
% {\red By standard results in the theory of ordinary differential equations \cite{abraham-marsden1978}, we easily get that the Hamiltonian is a complete vector field on $\mathrm{T}^*(\S_\alpha)$.\todo{Sam: need to find the right necessary condition.}
% The sub-Riemannian Hopf-Rinow theorem (see \cite[Theorem 11.38]{ABB-srgeom}) thus implies that the metric space $(\S_\alpha, \diff_{\S_\alpha})$ is complete. Therefore,}
Here note that the extremal reaches to the undefined point $x=\pm \pi/2$ only if $v\equiv 0$ (this follows from the non-integrability of $\tan(x)$ near $x=\pm\pi/2$). In this case, a (Euclidean) large circle passing through the north pole becomes a length minimizing geodesic. By completeness, there is a sub-Riemannian geodesic between every two points of $\S_\alpha$ by \cite[Theorem 2.5.23]{burago-burago-sergei2001}. These are obtained from Hamilton's equation \cref{eq:alphaGrushinsphereHamilton} since there are no non-trivial abnormal geodesics.

If a horizontal path of $\Sbb_\alpha$ is contained in both $\overline{\S}^+_\alpha$ and $\overline{\S}^-_\alpha := \{ p \in \S^2 \mid x \in [-\pi/2,0] \}$,
then a reflection $(x, y) \mapsto (-x,y)$ of the part of path that is in $\overline{\S}^-_\alpha$ produces a curve contained in $\overline{\S}^+_\alpha$ with the same length. This shows that a geodesic between points in the $\alpha$-Grushin hemisphere $\overline{\S}^+_\alpha$ is contained within $\overline{\S}^+_\alpha$. Length-minimisers are smooth since they satisfy Hamilton's equation \cref{eq:alphaGrushinsphereHamilton}. Thus, a constant-speed minimising geodesic $\gamma(t) = (x(t), y(t))$ that touches the singular equator at a point other than its endpoints must do so tangentially, and \cref{eq:alphaGrushinsphereHamilton} implies that $x(t)$ vanishes for all $t$. Consequently, $y(t)$ also vanishes, and $\gamma$ becomes a constant curve. In particular, a minimising geodesic between points of $\S_\alpha^+$ is also contained within $\S_\alpha^+$.

The fact that $\S_\alpha^+$ is also an incomplete Riemannian manifold follows easily since it doesn't contain any singular points of $\S_\alpha$.
% }
\end{proof}

% {\red Sam: I am thinking to combine the following with the previous result into one.

% The open hemisphere $\S_\alpha^+$ is also a geodesically convex subset. Actually, it is also a (non-complete) Riemannian manifold since it doesn't contain any singular points of $\S_\alpha$.

% \begin{prop}
%     The $\alpha$-Grushin open hemisphere $\S_\alpha^+$ is a geodesically convex subset of $\overline{\S}^+_\alpha$ and has the structure of a non-complete weighted Riemannian manifold when equipped with the restriction of the measure $\mathfrak{m}^\beta_{\S_\alpha}$.
% \end{prop}

% \begin{proof}
%     {\red I need to add but before I will check again if we really need strong convexity}
% \end{proof}

% }}
% \todo[inline]{Sam: should we use $\theta$ and $\phi$ as symbol for the coordinates of the sphere? Since they represent angles?\\
% Ken: I prefer $x,y$ since we use Einstein summation convention in Ricci computation.
% How do you think to put the green remark?\\
% Sam: Ok then, let's use $x, y$. The remark is helpful!}

% {\red \paragraph{\textbf{The $\alpha$-Grushin sphere}}Let $\S^2$ be the $2$-dimensional unit sphere.
% We endow the (singular) coordinates system $\{(x,y)\}$ as follows (for its validity, see \cite{boscain-prandi-seri2016}).
% Fix a large circle $\gamma:\R/2\pi\Z\to \S^2$,
% and let $N$ and $S$ be the north pole and south pole of $\S^2$ respectively, with respect to $\gamma$.
% For $p\in \S^2\setminus\{N,S\}$,
% let $x\in \R/2\pi\Z$ be the unique number such that $\gamma(x)$ is the perpendicular foot from $p$ to $\mathrm{Im}(\gamma)$,
% and $y$ be the (spherical) distance $\diff_{\S^2}(p,\gamma(x))$.
% It gives a parametrization of $\S^2\setminus\{N,S\}$ by $(x,y)\in \R/2\pi\Z\times (-\pi/2,\pi/2)$.
% Furthermore, it naturally extends to $N,S$ under the identification
% $(\frac{\pi}{2},y)\sim ( \frac{\pi}{2},y_2)$ and $(-\frac{\pi}{2},y)\sim(-\frac{\pi}{2},y_2)$ for any $y,y_2\in \R/2\pi\Z$.
% We call the coordinates $\{(x,y)\}=\R/2\pi\Z\times[-\pi/2,\pi/2]/\sim$ the suspension coordinates.
% \todo[inline]{Sam: why suspension?\\
% Ken: for me it looks like a topological suspention of the circle, but no preference\\ Sam: ok, I am not sure it's necessary to say it, we never mention it later?: Maybe no, so let's remove!}
% Note that,
% under this coordinates,
% the metric tensor $\g$ of the spherical metric is
% \[\g=\diff x\otimes \diff x+\cos^2(x)\diff y\otimes \diff y.\]
% % \todo[inline]{Sam: why not use the usual spherical coordinates? $g = \diff \theta \oplus \diff \theta + \sin^2(\theta) \diff \phi \oplus \diff \phi$. What's the advantage of introducing a general great circle $\gamma$?\\
% % Ken: I think it induces a different metric different from the standard spherical metric.}

% For $\alpha\geq 0$,
% define two vector fields
% \[X=\partial_1,~~Y=\frac{|\sin(x)|^{\alpha}}{\cos(x)}\partial_2,\]
% and define an almost Riemannian metric $g_\alpha$ so that $X$ and $Y$ are orthonormal.
% We call $(\S^2,g_\alpha)$ the $\alpha$-Grushin sphere.
% It is a generalization of the Euclidean sphere ($\alpha=0$),
% and the Grushin sphere ($\alpha=1$), see \cite{boscain-prandi-seri2016,Pan23}.
% }


\subsection{The \texorpdfstring{$\alpha$}{alpha}-Grushin hyperbolic plane and half-plane}\label{section:alphaGrushin2}

On the two-dimensional hyperbolic plane $(\bbH^2, g_{\bbH^2})$, we consider the following coordinate chart, called Lobachevsky's coordinates (see \cite[Section 33.1]{Martin-book}, for example).
%\todo{maybe we add a reference for that Ken: Note tha this coordinates is called the Lobachevsky coordinates, but not introduced by him. I think we should find a reference other than Russian, but it looks nothing}
We fix an infinite minimising geodesic ray $\gamma : \R \to \bbH^2$, and for $p \in \bbH^2$, we let $x := x(p) \in \R$ be signed hyperbolic distance $\diff_{\bbH^2}(p, \mathrm{Im}(\gamma))$,
where the signature is positive (resp. negative) if $p$ belongs to the left hand side (resp. right hand side) of $\gamma$. Furthermore, let $y=y(p)\in \R$ be the unique number such that $\gamma(y)$ is the perpendicular foot from $p$ to $\mathrm{Im}(\gamma)$. The map $\bbH^2 \to \R \times \R : p \mapsto (x(p), y(p))$ defines a global coordinate chart, and a short computation shows that the hyperbolic Riemannian metric $g_{\bbH^2}$ has the warped product structure
\[
g_{\bbH^2} = \diff x\otimes \diff x+\cosh^2(x)\diff y\otimes \diff y.
\]

\begin{definition}\label{def:hyperbolic}
    For $\alpha \geq 0$, the $\alpha$-Grushin hyperbolic plane $\bbH_\alpha$ is the sub-Riemannian structure on $\bbH^2$ induced from the vector field $X$ and $Y^\alpha$ given by
\[
X:=\partial_{x},~~Y^\alpha:=\frac{\abs{\sinh(x)}^{\alpha}}{\cosh(x)}\partial_{y}.
\]
The $\alpha$-Grushin hyperbolic half-plane $\overline{\bbH}^+_\alpha$ (resp. open half-plane $\bbH^+_\alpha$) is the subset of $\bbH_\alpha$ defined by
\[
\overline{\bbH}^+_\alpha := \{ p \in \bbH^2 \mid x \geq 0 \} \text{ (resp. $\bbH^+_\alpha := \{ p \in \bbH^2 \mid x > 0 \}$)}.
\]
\end{definition}
\noindent The $0$-Grushin hyperbolic plane is simply the two-dimensional Riemannian hyperbolic plane $\bbH^2$. When $\alpha > 0$, the $\alpha$-Grushin hyperbolic plane is a two-dimensional almost-Riemannian structure with $\{x = 0\}$ being its set of singular points.
% {\blue
To the best of our knowledge, this definition, although very natural, is new.
% } \todo{Sam: what do you think? Is that true? Ken: At least Pan and Honda didn't know, and it a.s. means that nobody knows.}.
At non-singular points, this sub-Riemannian structure admits the Riemannian metric
\begin{equation}
    \label{eq:HalphaRiemannianMetric}
    g_{\bbH_{\alpha}} = \diff x \otimes \diff x + \frac{\cosh^2(x)}{\sinh^{2\alpha}(x)} \diff y \otimes \diff y.
\end{equation}
A simple computation shows that the Riemannian volume induced from \cref{eq:HalphaRiemannianMetric} is given by
\[
\vol_{\bbH_\alpha}=\cosh(x)\abs{\sinh(x)}^{-\alpha}\diff x\diff y.
\]
We introduce the following weighted measure.
\begin{definition}\label{def:hyperbolicmeasure}
For $\beta\geq \alpha$,
we consider the weighted measure given by
\begin{equation*}
    \label{eq:MeasurealphaGrushinHyperbolic}
    \m^\beta_{\bbH_\alpha}:=\abs{\sinh(x)}^\beta\vol_{\bbH_\alpha}=\cosh(x)\abs{\sinh(x)}^{\beta-\alpha}\diff x \diff y=e^{-V_{\bbH_\alpha}}\vol_{\bbH_\alpha},
\end{equation*}
where $V_{\bbH_\alpha}(x,y):=-\beta\log\abs{\sinh(x)}$.
\end{definition}

\noindent As for the previous section, the $\alpha$-Grushin hyperbolic half-plane is a geodesically convex subset of $\bbH_\alpha$ and is a
% {\red uniquely}
geodesic space when seen as a length subspace of $\bbH_\alpha$. The open half-plane $\bbH_\alpha^+$ is also a geodesically convex subset and it is an incomplete Riemannian manifold since it doesn't contain any singular points of $\bbH_\alpha$.

\begin{prop}
    \label{prop:geodesicspacealphaGrushinHyperbolic}
     There is a
     % {\red unique}
     minimising geodesic contained within $\overline{\bbH}^+_\alpha$ that joins any two given points in $\overline{\bbH}^+_\alpha$. Furthermore, The $\alpha$-Grushin hyperbolic open half-plane $\bbH_\alpha^+$ is a geodesically convex subset of $\overline{\bbH}_\alpha^+$ and has the structure of an incomplete weighted
     % \todo{26. typo. DONE.}
     Riemannian manifold when equipped with the restriction of the measure $\mathfrak{m}^\beta_{\bbH_\alpha}$.
    % {\red In addition,
    % $\gamma$ is written by the image of the exponential map defined in \ref{expo}. Sam: Maybe we do not really use this information anymore?\\
    % Ken: Agree, maybe not. What we need is (strongly) geodesically convexity of $M$, and it can be proved without exponential map.}
\end{prop}

\begin{proof}
% {\red Sam: We add the proof or we explain how it is the same as the previous section?\\
% Ken: I think just mentioning the prvious proposition is enough}
Here, the Hamiltonian and the corresponding Hamilton's equation are given in Lobachevsky's coordinates, by
\begin{equation*}
    \label{eq:alphaGrushinHyperbolicHamilton}
    H(\lambda)=\frac{1}{2}\left[u^2+\frac{\sinh^{2\alpha}(x)}{\cosh^2(x)}v^2\right], \text{ and } \begin{cases}
    \dot{x}=u,\\
    \displaystyle \dot{y}=\frac{\sinh^{2\alpha}(x)}{\cosh^2(x)}v,\\
    \displaystyle \dot{u}= - 2 v^2 \ \sinh^{2(\alpha - 1)}(x) \tanh(x) (\alpha - \tanh^2(x)),\\
    \dot{v}=0.
\end{cases}
\end{equation*}
The rest of the proof follows exactly the arguments of \cref{prop:geodesicspacealphaGrushinSphere}.
\end{proof}

% {\red The open half-plane $N_\alpha^+$ is also a geodesically convex subset and it is a (non-complete) Riemannian manifold since it doesn't contain any singular points of $\bbH_\alpha$.

% \begin{prop}
%     The $\alpha$-Grushin hyperbolic open half-plane $N_\alpha^+$ is a strongly geodesically convex subset of $\bbH_\alpha^+$ and has the structure of a non-complete weighted Riemannian manifold when equipped with the restriction of the measure $\mathfrak{m}^\beta_{\bbH_\alpha}$.
% \end{prop}

% \begin{proof}
%     {\red Sam: Same, we add the proof or we explain how it is the same as the previous section?}
% \end{proof}
% }
% {\red The unweighted Riemannian (Popp's) measure is
% \[\vol=\cosh(x)|\sinh(x)|^{-\alpha}\diff x\diff y.\]
% For $\beta\geq 0$,
% we consider the weighted measure
% \[\m_\beta:=|\sinh(x)|^\beta\vol=\cosh(x)|\sinh(x)|^{\beta-\alpha}\diff x\diff y=e^{-f}\vol,\]
% where $f(x,y):=-\beta\log|\sinh(x)|$.}

% {\red Let us adopt the parametrization of the hyperbolic plane $\bbH^2$ by Lobachevsky's Cartesian-style coordinates $\{(x,y)\}$ defined as follows.
% First, fix an infinite geodesic ray $\gamma:\R\to \bbH^2$.
% For $p\in \bbH^2$,
% let $x\in\R$ be the unique number such that $\gamma(x)$ is the perpendicular foot from $p$ to $\gamma(\R)$,
% and $y$ the hyperbolic distance $\diff_{\bbH^2}(p,\gamma(x))$.
% Then $\{(x,y)\}$ gives a global coordinates of $\bbH^2$.
% In this coordinates,
% the metric tensor of the hyperbolic metric is
% \[\g=\diff x\otimes \diff x+\cosh^2(x)\diff y\otimes \diff y.\]

% % \todo[inline]{Sam: Same question here, why not use the simple usual hyperbolic coordinates?\\
% % The same reason to sphere. In addition, if we use other well-known coordinates like half-plane, then I checked that it is very difficult to $\alpha$-Grushinize it.}

% For $\alpha\geq 0$,
% consider two vector fields
% \[X=\partial_1,~~~~Y=\frac{|\sinh(x)|^\alpha}{\cosh(x)}\partial_2.\]
% This degenerated vector fields defines the almost Riemannian metric $g_\alpha$ on $\bbH^2$.
% We call $(\bbH^2,g_\alpha)$ the $\alpha$-Grushin hyperbolic plane.
% It is a generalization of the hyperbolic plane ($\alpha=0$).}
% \todo[inline]{K: I'm feeling that $\infty $-Grushin plane in the appendix is better. S: both fine to me. The paper is short, maybe it's better not to have an appendix? I am not sure. K:Ok!!}

\subsection{The \texorpdfstring{$\infty$}{infty}-Grushin plane and half-plane}\label{section:alphaGrushin3}

The geometry of the so-called $\alpha$-Grushin plane, where $\alpha \geq 0$, has been studied in \cite{li2012}, \cite{borza2022}, and \cite{borza2023}. The $\alpha$-Grushin half-plane and the validity of the $\cd$ condition in this space is studied in \cite{rizzi2023failure}, and we recalled some details in \cref{section:Introduction} . Instead, we introduce a model of a Grushin plane with infinite Hausdorff dimension. The global chart $(x, y)$ simply denotes the cartesian coordinates in this section.

\begin{definition}\label{def:infty}
    The $\infty$-Grushin plane $\mathbb{G}_\infty$ is the sub-Riemannian structure on $\R^2$ induced from the vector field $X$ and $Y$ given by
\[
X:=\partial_{x},~~Y:= e^{-1/\abs{x}}\partial_{y}.
\]
The $\infty$-Grushin half-plane $\overline{\mathbb{G}}^+_\infty$ (resp. open half-plane $\mathbb{G}^+_\infty$) is the subset of $\mathbb{G}_\infty$ defined by
\[
\overline{\mathbb{G}}^+_\infty := \{ p \in \R^2 \mid x \geq 0 \} \text{ (resp. $\mathbb{G}^+_\infty := \{ p \in \R^2 \mid x > 0 \}$)}.
\]
\end{definition}

\begin{lemma}
\label{lemma:infinitehausdorffdimension}
    The $\infty$-Grushin plane and half-plane have infinite Hausdorff dimension.
\end{lemma}
% \todo{27. There should be a short proof: since, locally, $\di_\infty\geq \di_\alpha$ for any $\alpha$, $\dim \G_\infty\geq \dim \G_\alpha$ for any $\alpha$. Yes, let's change\\ Ken: Changed. DONE.}
\begin{proof}
We will show that the Hausdorff dimension of $\mathcal{S}:=\{(0,y)\mid y\in\R\}\subseteq \G_\infty$ is $+\infty$. Let us denote by $\di_\alpha$ (resp. $\di_\infty$) the induced distance on $\G_\alpha$ (resp. $\G_\infty$).
It is well-known that the Hausdorff dimension of $\Scal\subseteq \G_\alpha$ is $\alpha+1$, see e.g. \cite{Franchi1983}. Since the inequality $|x|^\alpha\geq e^{-1/|x|}$ holds for sufficiently small $|x|$,
we have the inequality $\di_\alpha\geq \di_\infty$ in a small neighbourhood of an arbitrary point in $\Scal$. This implies that $\dim_H(\Scal,\di_\alpha)\leq \dim_H(\Scal,\di_\infty)$ and concludes the lemma.
% {\red Let $\mathcal{S}:=\{(0,y)\mid y\in\R\}$ be the singular line of $\mathbb{G}_\infty$.
% We will show that the Hausdorff dimension of $\mathcal{S}$ is infinite.
% We denote by $\diff_{\mathbb{G}_\infty}$ the sub-Riemannian distance of $\mathbb{G}_\infty$ and by $\sfD_\infty$ the sub-Finsler distance on $\R^2$ derived from the $\ell^{\infty}$-norm with respect to the frame $\mathcal{F} = \{X,Y\}$.
% Trivially, we have that $\sfD_{\infty}\leq \diff_{\mathbb{G}_\infty}\leq 2\sfD_\infty$. In particular, the distances $\sfD_{\infty}$ and $\diff_{\mathbb{G}_\infty}$ are bi-Lipschitz equivalent and it is enough to investigate the Hausdorff dimension of $\mathcal{S}$ with respect to $\sfD_{\infty}$. For $y> 0$,
% let $v_y>0$ be the positive number such that
% \[\sfD_{\infty}(\e,(v_y,y/2))=\min_{v\in\R}\{\sfD_{\infty}(\e,(v,y/2))\},\]
% where $\e=(0,0)$ is the origin.
% Since the reflection along a horizontal line is an isometry,
% it holds that
% \[\sfD_{\infty}(\e,(0,y))=2 \ \sfD_{\infty}(\e,(v_y,y/2)).\]
% Let us consider the curve $c(t)=(c_1(t),c_2(t))$ defined through the system of differential equations
% \[\dot{c}_1(t)=1,~~\dot{c}_2(t)=e^{-1/c_1(t)}~~~~\text{with the initial value }c_1(0)=c_2(0)=0.\]
% Then, $(c_1,c_2)$ is a unit speed length minimiser of $(\R^2,\sfD_\infty)$.
% This implies that $v_y$ is the unique positive number such that
% \[\int_0^{v_y}e^{-1/s}ds=\frac{y}{2},\]
% and it also holds that
% $\sfD_{\infty}(\e,(0,y))=2v_y$.
% To prove that the Hausdorff dimension of $\mathcal{S}$ is greater than $N$ for any $N \in \rinterval{1}{+\infty}$,
% we will verify that for any $v>0$ small enough,
% $\Leb(B(v))=o(v^N)$ as $v \to 0$,
% where $B(v)\subseteq (\mathcal{S},\sfD_{\infty})$ is the ball of radius $v$ centered at $\e$. %\todo{with respect to $\sfD_{\infty}$ no? Ken: It is same by bi-Lipschitz equivalence}
%    For $v>0$,
%    let $y_v$ be the positive number such that
%    \[\int_0^{v}e^{-1/s}ds=\frac{y_v}{2}.\]
%    From the above argument,
%    the interval $[-y_v,y_v]$ is the closed ball of radius $2v$ in $(\mathcal{S},\sfD_\infty)$ centered at $\e$.
%    % Since $\sfD_{\infty}\leq \diff_{\mathbb{G}_\infty}$, it contains the ball of radius $v$ in $(\mathcal{S},\diff_\infty)$.
%    From the definition of $y_v$,
%    its differential is $\frac{\diff}{\diff v}y_v=e^{-1/v}.$
%    Therefore,
%    the growth of the $v$-ball is so small that $\mathcal{L}(B(v))=o(v^N)$ as $v \to 0$ for any positive number $N$.
%    We conclude that $(\mathcal{S},\diff_{\mathbb{G}_\infty})$ has infinite Hausdorff dimension.}
\end{proof}

The $\infty$-Grushin plane is a two-dimensional almost-Riemannian structure with two-dimensional almost-Riemannian structure with $\{x = 0\}$ being its set of singular points.
% {\blue
To the best of our knowledge, this definition of $\infty$-Grushin plane is also new.
%} \todo{Sam: what do you think? Is that true?K: I think so. If somebody pointed out, let's add it.}.
At non-singular points, this sub-Riemannian structure admits the Riemannian metric
\begin{equation}
    \label{eq:GinftyRiemannianMetric}
    g_{\mathbb{G}_{\infty}} = \diff x \otimes \diff x + e^{2/\abs{x}} \diff y \otimes \diff y.
\end{equation}
A simple computation shows that the Riemannian volume induced from \cref{eq:GinftyRiemannianMetric} is given by
\[
\vol_{\mathbb{G}_\infty}=e^{1/\abs{x}}\diff x\diff y.
\]
We introduce the following weighted measure.
\begin{definition}\label{def:infty_measure}
For
% {\red $\beta, \gamma \geq 0$}{\blue
$\beta\geq 0$ and $\gamma>0$
% }
,
we consider the weighted
% {\blue
(Radon)
% }
measure given by
\begin{equation*}
    \label{eq:MeasurealphaGrushininfty}
    \m^{\beta, \gamma}_{\mathbb{G}_{\infty}}:=\abs{x}^\beta e^{-\gamma/x^2}\vol_{\mathbb{G}_{\infty}}=\abs{x}^\beta e^{-\gamma/x^2 + 1/\abs{x}}\diff x\diff y=e^{-V}\vol_{\mathbb{G}_{\infty}},
\end{equation*}
where $V_{\G_\infty}(x,y):=\frac{\gamma}{x^2}-\beta\log\abs{x}$.
\end{definition}
% \todo[inline]{Sam: We should explain why this is a well-defined Radon measure. Is that the case for every $\beta, \gamma \geq 0$?\\
% Ken: The following argument is correct? If $\beta\geq 0$ and $\gamma>0$, then the measure does not diverges (locally finite). Other continuity property holds since the density function is continuous.\\
% Sam: Should this be made rigourous in a remark?}

The next result is analogous to the corresponding one in the previous two sections. The $\infty$-Grushin half-plane is a geodesically convex subset of $\mathbb{G}_\infty$ and a geodesic space when seen as a length subspace of $\mathbb{G}_\infty$. Similarly, the open half-plane $\mathbb{G}_\infty^+$ is also a geodesically convex subset and it is an incomplete Riemannian manifold since it doesn't contain any singular points of $\mathbb{G}_\infty$.

\begin{prop}
     There is a minimising geodesic contained within $\overline{\mathbb{G}}^+_\infty$ that joins any two given points in $\overline{\mathbb{G}}^+_\infty$. Furthermore, the $\infty$-Grushin open half-plane $\mathbb{G}_\infty^+$ is a geodesically convex subset of $\overline{\mathbb{G}}_\infty^+$ and has the structure of an incomplete weighted Riemannian manifold when equipped with the restriction of the measure $\mathfrak{m}^{\beta, \gamma}_{\mathbb{G}_\infty}$.
    % {\red In addition,
    % $\gamma$ is written by the image of the exponential map defined in \ref{expo}. Sam: Maybe we do not really use this information anymore?\\
    % Ken: Agree, maybe not. What we need is (strongly) geodesically convexity of $M$, and it can be proved without exponential map.}
\end{prop}

\begin{proof}
The argument is again analogous to the proofs of \cref{prop:geodesicspacealphaGrushinSphere} and \cref{prop:geodesicspacealphaGrushinHyperbolic}, with the Hamiltonian and Hamilton's equation given in cartesian coordinates by
\begin{equation*}
    \label{eq:alphaGrushinInfHamilton}
    H(\lambda)=\frac{1}{2}\left[u+e^{-2/\abs{x}}v^2\right], \text{ and } \begin{cases}
    \dot{x}=u,\\
    \displaystyle \dot{y}=e^{-2/\abs{x}}v,\\
    \displaystyle \dot{u}= - v^2 \frac{e^{-2/\abs{x}}}{\abs{x}^2} x,\\
    \dot{v}=0.
\end{cases}
\end{equation*}
\end{proof}

Hereafter, we will collectively refer to the $\alpha$-Grushin plane (resp. half-plane), the $\infty$-Grushin plane (resp. half-plane), the $\alpha$-Grushin sphere (resp. hemisphere), and the $\alpha$-Grushin hyperbolic plane (resp. half-plane) as the $\alpha$-Grushin spaces (resp. half-spaces).

% {\red Similarly, the open half-plane $P_\infty^+$ is also a geodesically convex subset and it is a (non-complete) Riemannian manifold since it doesn't contain any singular points of $\mathbb{G}_\infty$.

% \begin{prop}
%     The $\infty$-Grushin open half-plane $P_\infty^+$ is a strongly geodesically convex subset of $\mathbb{G}_\infty^+$ and has the structure of a non-complete weighted Riemannian manifold when equipped with the restriction of the measure $\mathfrak{m}^{\beta, \gamma}_{\mathbb{G}_\infty}$.
% \end{prop}

% \begin{proof}
%     {\red Sam: Same, I add later}
% \end{proof}
% }
% On $\R^2$ with the cartesian coordinates $\{(x,y)\}$,
% define vector fields
% \[X=\partial_1,~~~~Y=e^{-\frac{1}{|x|}}\partial_2.\]
% This degenerated vector fields define the almost Riemannian metric $\g_\infty$ on $\R^2$.
% We call the pair $(\R^2,\g_\infty)$ the $\infty$-Grushin plane.
% Denote by $\diff_\infty$ the induced distance function on $\R^2$.

% \begin{lemma}
%     The Hausdorff dimension of the $\infty$-Grushin plane is $\infty$.
% \end{lemma}

% \begin{proof}
%         Let $\mathcal{S}:=\{(0,y)\mid y\in\R\}$ be the singular line.
%         We will show that the Hausdorff dimension of $(\mathcal{S},\diff_\infty)$ is $\infty$.

%     Denote by $\sfD_\infty$ the sub-Finsler distance on $\R^2$ derived from the $\ell^{\infty}$ norm with respect to the frame $\{X,Y\}$.
% Trivially we have $\sfD_{\infty}\leq \diff_\infty\leq 2\sfD_\infty$,
% thus $\sfD_{\infty}$ determines the bi-Lipschitz equivalence class of $\diff_\infty$.
% For $y> 0$,
% let $v_y>0$ be the positive number such that

% \[\sfD_{\infty}(\e,(v_y,y/2))=\min_{v\in\R}\{\sfD_{\infty}(\e,(v,y/2))\},\]
% where $\e=(0,0)$ is the zero point.
% Since the reflection along a horizontal line is an isometry,
% it holds
% \[\sfD_{\infty}(\e,(0,y))=2\sfD_{\infty}(\e,(v_y,y/2)).\]
% Let us consider a curve $c(t)=(c_1(t),c_2(t))$ defined by
% \[\dot{c}_1(t)=1,~~\dot{c}_2(t)=e^{-\frac{1}{c_1(t)}}~~~~\text{with the initial value }c_1(0)=c_2(0)=0.\]
% Then $(c_1,c_2)$ is a unit speed length minimizer of $(\R^2,\sfD_\infty)$.
% It implies that $v_y$ is the unique positive number such that
% \[\int_0^{v_y}e^{-\frac{1}{s}}ds=y/2,\]
% and it holds
% $\sfD_{\infty}(\e,(0,y))=2v_y$.


% To prove $\dim_{\mathrm{H}}(\mathcal{S})\geq N$,
% we will check that for any $v>0$ small enough,
% $\Leb(B(v))=O(v^N)$,
% where $B(v)\subset (\mathcal{S},\diff_\infty)$ is the ball of radius $v$ centered at $\e$.

%    For $v>0$,
%    let $y_v$ be the positive number such that
%    \[\int_0^{v}e^{-\frac{1}{s}}ds=\frac{y_v}{2}.\]
%    From the above argument,
%    the interval $[-y_v,y_v]$ is the closed ball of radius $2v$ in $(\mathcal{S},\sfD_\infty)$ centered at $\e$.
%    Since $\sfD_{\infty}\leq \diff_\infty$,
%    it contains the ball of radius $v$ in $(\mathcal{S},\diff_\infty)$.
%    From the definition of $y_v$,
%    its differential is $\frac{\diff}{\diff v}y_v=e^{-\frac{1}{v}}.$
%    Therefore,
%    the growth of the $v$-ball is so small that $\mathcal{L}(B(v))=o(v^N)$ for any positive number $N$.
%    It implies that $(\mathcal{S},\diff_\infty)$ has infinite Hausdorff dimension.
% \end{proof}

% \subsection{Geodesics of \texorpdfstring{$\alpha$}{alpha}-Grushin model halfspaces}

% \todo[inline]{I change $(x,y)$ into $(x,y)$}
% {\red Let us start from length minimizers
% We study length minimizers on $\alpha$-Grushin hemisphere.
% First we start from an easy lemma.
%  \begin{lemma}
%     For given two points in $\Splu$,
%     there is a unique length minimizer $\gamma$ joining them in $(\Splu,g_\alpha)$.
%     In addition,
%     $\gamma$ is written by the image of the exponential map defined in \ref{expo}.
% \end{lemma}
% \begin{proof}

% Note that if a path on $\Sbb^2$ is contained in both of $\Splu$ and $\S_-$,
% then by the reflection along the singular equator $\{x=0\}$,
% we can convert the path on $\S^2$ into that on $\Splu$ with the same length.
% Moreover,
% if a path has a corner on the singular equator,
% then a visual shortcut yields a truly shorter path.
% Therefore every length minimizer touches the singular equator at most once,
% and it fails to minimize after it passes through the equator.
% It implies that every length minimizer is the restriction of one on $(\S^2,g_\alpha)$.

% Next we compute length minimizers on $(\S^2,g_\alpha)$ restricted to $\Splu$.
% Note that we always have $\sin(x),\cos(x)\geq 0$.
% Let $H:T^\ast \Splu$ be the sub-Riemannian Hamiltonian,
% that is,
% for the coordinates $\lambda=(x,y,u,v)\in T^\ast \Splu$,
% \[H(\lambda):=\frac{1}{2}\left[\langle \lambda,X\rangle^2+\langle\lambda,Y\rangle^2\right]=\frac{1}{2}\left[u^2+\frac{\sin^{2\alpha}(x)}{\cos^2(x)}v^2\right].\]
% Since we have no abnormal minimizers,
% every length minimizer $\gamma(t)=(x(t),y(t))$ is the solution to the Hamilton's equation
%     \[\begin{cases}
%     \dot{x}=u,\\
%     \dot{y}=\frac{\sin^{2\alpha}(x)}{\cos^2(x)}v,\\
%     \dot{u}=-[2\alpha\cos^2(x)+2\sin^2(x)]\cdot \frac{\sin^{2\alpha-1}(x)}{\cos^3(x)}\cdot v^2,\\
%     \dot{v}=0.
% \end{cases}\]
% From the equation,
% $v\equiv v_0$ is a constant function.
% If $v\equiv 0$,
% then $\gamma(t)=(x(t),y(t))=(x_0+u_0t,y_0)$ i.e. the great circle through the north pole.
% If $v_0\not\equiv 0$,
% then $y(t)$ is strictly increasing/decreasing,
% and $x(t)$ is concave.
% Since every length minimizer is the solution to the Hamilton's equation \ref{eq:Hamilton},
% the above monotonicity/convexity arguments imply that the length minimizer is uniquely deteremined.

% Therefore,
% for the initial point $z=(x_0,y_0)$,
% we can define the exponential map $\exp_{z}^t(u,v)=(x(t),y(t))$ defined on an appropriate open subset $O\subset T_{z}^\ast\Splu$\footnote{If $z$ is on the singular equator,
% the tangent space $T_z^\ast \Splu$ is a closed half plane $\R_{\geq 0}\times \R$} such that $\exp_z^1$ is diffeomorphism onto its image.


% \end{proof}}

\section{Equivalence between \texorpdfstring{$\cd(K,N)$}{CD(K,N)} and \texorpdfstring{$\Ric_N\geq K$}{RicN>=K} for some almost Riemannian manifolds}
\label{section:equivalence}

% \todo[inline]{Here? The intro could contain the $\mathsf{R}$ from $\mathsf{RCD}$ with the reference to Le Donne. Since we have $\mathsf{CD}$ we can note that we also have $\mathsf{BM}$ (usual sub-Riemannian never has $\mathsf{BM}$), and maybe other properties like Levy-Gromov
% }

% \todo[inline]{ken: To state $\cd(0,N)$ for negative $N$, we need to mention $iKRW$ distance for stability. Do we put the third statement?\\
% Sam: let's think about it. If we can adapt the theorem from section 3 to include $N$ negative then let's do it.}

% % \todo[inline]{Copy-Paste of Luca--Giorgio's paper}
% % \begin{equation}\label{eq:grushin_rm}
% % \g=\diff x\otimes\diff x+\frac1{x^2}\,\diff y\otimes\diff y,
% % \quad
% % x\ne0.
% % \end{equation}

% % \begin{equation}
% % \label{eq:grushin_wmeas_p}
% % \m_p=e^{-V}\vol_{\g}, \quad V(x)=-(p+1)\log |x|,
% % \quad
% % x\ne0.
% % \end{equation}
% % We call the metric-measure space $\mathbb{G}_p=(\R^2,\diff,\m_p)$ the \emph{($p$-)weighted Grushin plane}.


% % \begin{lemma}\label{lem:bericci}
% % Let $p\in \R$ and $N>2$. The $N$-Bakry--Émery Ricci tensor of the Grushin metric \eqref{eq:grushin_rm}, with weighted measure $\m_p = |x|^p\diff x \diff y$, for all $x\neq 0$ is
% % \begin{equation}
% % \mathrm{Ric}_{N,V} = \frac{p-1}{x^2}\g -\frac{(p+1)^2}{N-2}\frac{\diff x\otimes \diff x}{x^2},
% % \end{equation}
% % with the convention that $1/\infty=0$.
% % \end{lemma}

% % \begin{theorem}\label{res:halfgrushinCD}
% % Let $p\geq 1$. The weighted Grushin half-plane $\mathbb{G}_p^+ $ satisfies the $\cd(0,N)$ condition if and only if $N\geq N_p$, where $N_p \in (2,\infty]$ is given by
% % \begin{equation}
% % \label{eq:halfgrushin_Np}
% % N_p = \frac{(p+1)^2}{p-1}+2,
% % \end{equation}
% % with the convention that $N_1 = \infty$.
% % Furthermore, $\mathbb{G}_p^+ $ is infinitesimally Hilbertian, and it is thus an $\RCD(0,N)$ space for $N\geq N_p$.
% % \end{theorem}

% % \begin{proof}
% % The statement is a consequence of the geodesic convexity of $\mathbb{G}_p^+$ and the computation of the $N$-Bakry--Émery curvature in \ref{lem:bericci}.
% % Since the proof uses quite standard arguments, we simply sketch its main steps.

% % The interior of $\mathbb{G}_p^+$, i.e., the open half-plane, can be regarded as a (non-complete) weighted Riemannian manifold with metric $g$ as in \eqref{eq:grushin_rm} and weighted volume as in~\eqref{eq:grushin_wmeas_p}.
% % Let $\mu_0,\mu_1\in\prob_2(\mathbb{G}_p^+)$, $\mu_0,\mu_1<<\m_p$, with bounded support contained in the Riemannian region $\{x>\epsilon\}$, for some $\epsilon\geq 0$.

% % Let $(\mu_s)_{s\in[0,1]}$ be a $W_2$-geodesic joining $\mu_0$ and $\mu_1$.
% % By a well-known representation theorem (see \cite[Cor.\ 7.22]{villani}), there exists $\nu\in \prob(\mathrm{Geo(\mathbb{G}_p^+)})$, supported on the set $\Gamma=(e_0\times e_1)^{-1}(\supp\mu_0\times \supp\mu_1)$, such that $\mu_s = (e_s)_\sharp \nu$ for all $s\in[0,1]$.
% % Since the set $\{x\geq \epsilon\}$ is a geodesically convex subset of the full Grushin plane $\mathbb{G}_p$ (by the same argument of \cite[Prop.\ 5]{Rizzi-Gluing2018}), any $\gamma \in \Gamma$ is contained for all times in the region $\{x>0\}$.
% % Therefore, $\Gamma$ is a set of Riemannian geodesics contained in the weighted Riemannian structure $(\{x>0\},g,e^{-V}\mathrm{vol}_{g})$.
% % By \ref{lem:bericci}, we have $\mathrm{Ric}_{N,V} \geq 0$ for all $N\geq N_p$, where $N_p$ is as in~\eqref{eq:halfgrushin_Np}.
% % At this point, a standard argument shows that the Rényi entropy is convex along Wasserstein geodesics joining $\mu_0$ with $\mu_1$, see the proof of \cite[Thm. \ 1.7]{sturm2006} for example.

% % The extension to $\mu_0,\mu_1\in\prob_2(\mathbb G_p^+)$, with $\mu_0,\mu_1<< \m_p$ and compact support possibly touching the singular region $\{x=0\}$, is achieved via a standard approximation argument.
% % More precisely, one reduces to the previous case and  exploits the stability of optimal transport~\cite[Thm. \ 28.9]{villani} and the lower semi-continuity of the Rényi entropy~\cite[Thm. \ 29.20]{villani}.

% % Finally, the extension to general $\mu_0,\mu_1\in\prob_2(\mathbb G_p^+)$ follows the routine argument outlined in~\cite[Rem.\ 2.12]{Bacher-Sturm2010}, which works when $\mu_s = (e_s)_\sharp \nu$, $s\in[0,1]$, and $\nu$ is concentrated on a set of non-branching geodesics. This proves the `if' part of the statement.

% % The `only if' part is also standard. The $\cd(0,N)$ condition for $N>2$ implies that, on the Riemannian region $\{x>0\}$, $\mathrm{Ric}_{N,V} \geq 0$, but this is false for $N<N_p$.

% % The fact that $\mathbb{G}_{p}^+$ is infinitesimally Hilbertian follows from  \ref{rmk:inf_hilb_local}, by noting that $\m_p$ is positive and smooth out of the closed set $\{x=0\}$, which has zero measure. An alternative proof follows from the observation that $\mathbb{G}_p^+$ is a Ricci limit, see~\cite{panwei}.
% % \end{proof}

% \textbf{Sam: I tried very hard to get a general theorem. I think that the theorem below should (I hope) be correct, but let's check it. There are a two technical gaps that you may be able to see.}

% \begin{theorem}
% 	Let $(\X, \di, \mathfrak{m})$ be a (complete, locally compact, geodesic) metric measure space and $M$ be a Borel subset of $\X$ such that
% 	\begin{enumerate}[label=\normalfont(\roman*)]
% 	\item $M$ is a geodesically convex subset of $(\X, \di)$,
% 	\item $(M, \di|_M, \mathfrak{m}|_M)$ is a weighted Riemannian manifold $(M, \di_g, e^{-V} \mathrm{vol}_g)$,
% 	\item $\mathfrak{m}(\X \setminus M) = 0$ and $\lim_{x \in M \to \partial M} e^{-V(x)} = 0$.
% 	\end{enumerate}
%  Then, the metric measure space $(\X, \di, \mathfrak{m})$ satisfies the $\mathsf{RCD}(K, N)$ condition if and only if $\mathrm{dim}M \leq N$ and $\mathrm{Ric}_{V}^N \geq K$ on $(M, g)$.

% \end{theorem}

% \begin{proof}
% 	First of all, an $\mathsf{RCD}(K, N)$ condition on $(\X, \di, \mathsf{m})$ directly implies that $\mathrm{dim}M \leq N$ and $\mathrm{Ric}_{V}^N \geq K$ on $(M, g)$ by \cite[Part (e) of the proof of Theorem 1.7]{sturm2006}. We therefore focus now on the other implication.{\red Should we cite \cite[Thm\, 17.36]{villani} for the coherence?}

% 	Consider $\mu_0, \mu_1 \in \mathscr{P}_2^{\mathrm{ac}}(\X, \di, \mathfrak{m})$ with continuous densities $\rho_0$ and $\rho_1$ respectively. We need to argue that there exists a $W_2$-geodesic $(\mu_s)_{s \in \interval{0}{1}}$ joining $\mu_0$ to $\mu_1$ with $\mu_s = (e_s)_\sharp\nu$ for some $\nu \in \mathscr{P}(\mathrm{Geo}(\X, \di))$ such that the $N'$-Rényi entropy is $K$-convex along $(\mu_s)_{s \in \interval{0}{1}}$ for all $N' \geq N$.

% 	\paragraph{Step 1. The family of sets $(M_\epsilon)_{\epsilon \geq 0}$.} For all $\epsilon \geq 0$, we define $M_\epsilon := \left\{x \in M \mid \di_g(x, \partial M) > \epsilon \right\}$ and we note that the closure $\overline{M}_\epsilon$ of $M_\epsilon$ in $\X$ is $\left\{x \in \X \mid \di(x, \partial M) \geq \epsilon \right\}$. Conditions (i) and (ii) implies that for $\epsilon \geq 0$ small enough, the set $\overline{M}_\epsilon$ is a geodesically convex subset of $\X$. {\red Sam: is that true? is that obvious? How to prove it?}
% 	% {\green I think we need the following assumption instead of geodesically convexity:
%  % \begin{itemize}
%  %     \item[(i')]For every $x\in \partial M$,
%  %     there is a local chart $(x,\dots,x_n)\in \R_{\geq 0}\times \R^{n-1}$ around $x=(0,\dots,0)$ such that,
%  %     an orthonormal frame $x,\dots,X_n$ is written by
%  %     \[x=\partial_1,~~X_i=\sum_{j=2}^nf_{ij}(x,\dots,x_n)\partial_j,\]
%  %     where $f_{ij}$ is a positive smooth function such that $f_{ij}(\cdot,y,\dots,x_n)$ is strictly increasing + some uniform conditions (maybe we do not need uniform condition if we restrict ourselves to a given $\mu_i^\epsilon$ with the supports are in a compact subsets, not global $M_\epsilon$).
%  % \end{itemize}}
% 	\paragraph{Step 2. The support of $\mu_0$ and $\mu_1$ are compact and contained in $M_\epsilon$ for some $\epsilon > 0$.}
%  If that is the case for some $\epsilon > 0$ small enough, then geodesic convexity of $\overline{M}_{\epsilon}$ and the Riemannian structure that $M$ possesses implies, by \cite[Theorem 13]{mccann}, that there exists a unique $W_2$-geodesic $(\mu_s)_{s \in \interval{0}{1}}$ joining $\mu_0$ to $\mu_1$. By the Representation Theorem \cite[Corollary 7.20]{villani}, we also know that there exists $\nu \in \mathscr{P}(\mathrm{Geo}(\X, \di))$ such that $\mu_s = (e_s)_\sharp \nu$ for all $s \in \interval{0}{1}$ and $\mathrm{supp}(\nu) \subseteq \Gamma := (e_0 \times e_1)^{-1}(\mathrm{supp}(\mu_0) \times \mathrm{supp}(\mu_1))$. A geodesic in $\gamma \in \Gamma$ joins points in $\mathrm{supp}(\mu_0)$ with points in $\mathrm{supp}(\mu_1)$, and since $\overline{M}_\epsilon$ is geodesically convex we infer that $\gamma$ in contained in $M$. Since $\mathrm{dim}M \leq N$ and $\mathrm{Ric}_{V}^N \geq K$ on $(M, g)$, \cite[Parts (a), (b), (c), and (d) of the proof of Theorem 1.7]{sturm2006} implies that the $N'$-Rényi entropy is $K$-convex along $(\mu_s)_{s \in \interval{0}{1}}$ for all $N' \geq N$.

% 	\paragraph{Step 3. The support of $\mu_0$ and $\mu_1$ are compact and contained in $\X$.} In this case, we consider $i \in \{0, 1\}$ and $\epsilon > 0$, and we approximate $\rho_i$ with continuous densities $\rho_i^\epsilon$ such that $\mu_i^\epsilon := \rho_i^\epsilon \mathfrak{m}$ has support contained in $M_\epsilon$ and such that $\mu_i^\epsilon \rightharpoonup \mu_i$ as $\epsilon \to 0$. {\red (Sam: Need to figure out the right convergence and to argue that we can do this approximation. Here it should be important that $M$ is open for the topology of $\X$ and that (iii) holds)} Then, by stability of optimal transport (see \cite[Theorem 28.9]{villani}), the $W_2$-geodesic $(\mu_s^\epsilon)_{s \in \interval{0}{1}}$ joining $\mu_0^\epsilon$ to $\mu_1^\epsilon$ converges, as $\epsilon \to 0,$ to a $W_2$-geodesic $(\mu_s)_{s \in \interval{0}{1}}$ joining $\mu_0$ to $\mu_1$. The same behaviour happens for the measure $\nu^\epsilon \in \mathscr{P}(\mathrm{Geo}(\X, \di))$ from the Representation Theorem. By the previous step, the $N'$-Rényi entropy is $K$-convex along $(\mu_s^\epsilon)_{s \in \interval{0}{1}}$ for all $N' \geq N$. The lower semi-continuity of the Rényi entropy (see \cite[Theorem 29.20]{villani}) allows us to conclude that the $N'$-Rényi entropy is also $K$-convex along $(\mu)_{s \in \interval{0}{1}}$ for all $N' \geq N$.

% 	\paragraph{Step 4. The general case.} The case where the support of $\mu_0$ and $\mu_1$ are not necessarily compact is obtained by the previous step and by using an exhaustion by compact sets.

%  \todo[inline]{I wrote down so that the convergence and how-to-apply Thm 28.9 is clear. I'm wondering if we need to consider compact support}
% 	{\blue

% \begin{theorem}
% 	Let $(\X, \di, \mathfrak{m})$ be a (complete, separable) metric measure space and
% 	$(M_k)_{k\in\N}$ be a sequence of open subsets of $\X$ such that
% 	\begin{enumerate}[label=\normalfont(\roman*)]
% 	\item $M_k$ is a geodesically convex subset of $(\X, \di)$,
% 	\item $(M_k, \di|_{M_k}, \mathfrak{m}|_{M_k})$ possesses a weighted Riemannian structure,
%  \item $\bar M_k\subset M_{k+1}$,
%  \item $\bar M_k$ converges to $\X$ as $k\to\infty$ in the pointed measured Gromov--Hausdorff topology, with the inclusion maps $\iota_k$ being approximation maps.
% 	\end{enumerate}
%  \todo[inline]{Sam: Do you think this your (i), (iii) and (iv) is equivalent to the previous (i)?}
% 	Then, the metric measure space $(\X, \di, \mathfrak{m})$ satisfies the $\mathsf{RCD}(K, N)$ condition if and only if $\mathrm{dim}(M_k) \leq N$ and $\mathrm{Ric}_{V}^N \geq K$ on $M_k$ for all $k\in\N$.
% \end{theorem}

% 	\paragraph{Step 1. The support of $\mu_0$ and $\mu_1$ are contained in $M_k$ for some $k\in\N$.}
%  If that is the case, then the assumptions (i, ii) on $M_k$ and results in \cite{Fathi-Figalli2010,Figalli2007,Figalli-Gigli2011} implies that there exists a unique $W_2$-geodesic $(\mu_s)_{s \in \interval{0}{1}}$ joining $\mu_0$ to $\mu_1$, induced from the transportation map $T_s(x)=\exp_x(s\nabla \phi(x))$ for some $\mu_0$-almost everywhere differentiable function $\phi$.\todo{Ken:I'm not satisfied with this argument:( What we need is there is such $\phi$, but I'm not sure what is the best reference}
%  Note that $\supp(\mu_s)\subset M_{k}$ holds by the assumption (i). Combined with the assumptions $\mathrm{dim}(M_k) \leq N$ and $\mathrm{Ric}_{V}^N \geq K$ on $M_{k}$, \cite[Thm\, 17.36]{villani} implies that the $N$-Rényi entropy is $K$-convex along $(\mu_s)_{s \in \interval{0}{1}}$.

% 	\paragraph{Step 2. The support of $\mu_0$ and $\mu_1$ are contained in $\X$.} By the assumption (iv) and \cite[Thm\, 28.12]{villani},
%  there are probability measures $\mu_i^k$ $(i=0,1)$ such that the supports are contained in $M_k$ and that $(\iota_k)_\#\mu_i^k \to\mu_i$ in $W_2$-distance.
%  Let $\nu_k$ be the optimal transport plan joining $\mu_0^k$ to $\mu_1^k$ in $M_k$.
%  Then, by \cite[Thm\, 28.9 and Exer\, 28.15]{villani},
%  $\nu_k$ weakly converges to an optimal transport plan $\nu$ in $\X$ joining $\mu_0$ to $\mu_1$,
%  and the $W_2$-geodesics $(\mu_s^k)_{s \in \interval{0}{1}}$ uniformly converges to a $W_2$-geodesic $(\mu_s)_{s \in \interval{0}{1}}$ in $\X$ joining $\mu_0$ to $\mu_1$.
%  By the previous step and (iii), the $N$-Rényi entropy is $K$-convex along $(\mu_s^k)_{s \in \interval{0}{1}}$.
%  The same arguments with \cite[Thm\, 29.24]{villani}, i.e. the proof of Eq. (29.40) from Eq. (29.39), allows us to conclude that the $N$-Rényi entropy is also $K$-convex along $(\mu_s)_{s \in \interval{0}{1}}$.
%  }
% \end{proof}

% \todo[inline]{Ken: Under a stronger condition such $M_k$ exist, but I'm not sure if strongly $\Leftrightarrow$ geod. cvx?}
% {\green

% \begin{theorem}
% 	Let $K \in \R$, $N \in \interval{1}{+\infty}$, $(\X, \di, \mathfrak{m})$ be a (complete, separable) metric measure space and $M$ be an open subset of $\X$ such that
% 	\begin{enumerate}[label=\normalfont(\roman*)]
% 	\item $M$ is a strongly geodesically convex subset of $(\X, \di)$, i.e. for any closed subset $K$ in $M$, $\overline{\mathrm{Cvxhul}}(K)\subset M$,
% 	\item $(M, \di|_M, \mathfrak{m}|_M)$ possesses a weighted Riemannian manifold structure,
% 	\item $\mathfrak{m}(\X \setminus M) = 0$ and $\X\setminus M\subset \partial M$.
% 	\end{enumerate}
% 	Then, the metric measure space $(\X, \di, \mathfrak{m})$ satisfies the $\mathsf{RCD}(K, N)$ condition if and only if $\mathrm{dim}M \leq N$ and $\mathrm{Ric}_{V}^N \geq K$ on $M$.
% \end{theorem}

% \begin{proof}
% 	First of all, an $\mathsf{RCD}(K, N)$ condition on $(\X, \di, \mathsf{m})$ directly implies that $\mathrm{dim}M \leq N$ and $\mathrm{Ric}_{V}^N \geq K$ on $M$ by \cite[Thm\, 17.36]{villani}. We therefore focus now on the other implication.

% 	Consider $\mu_0, \mu_1 \in \mathscr{P}_2^{\mathrm{ac}}(\X, \di, \mathfrak{m})$ with continuous densities $\rho_0$ and $\rho_1$ respectively. We need to argue that there exists a $W_2$-geodesic $(\mu_s)_{s \in \interval{0}{1}}$ joining $\mu_0$ to $\mu_1$ with $\mu_s = (e_s)_\sharp\nu$ for some $\nu \in \mathscr{P}(\mathrm{Geo}(\X, \di))$ such that the $N$-Rényi entropy is $K$-convex along $(\mu_s)_{s \in \interval{0}{1}}$.

%  \todo[inline]{Sam: Shouldn't step 1 be ``support of $\mu_0$ and $\mu_1$ are contained in $M_k$.'' (and not $M$)?\\
%  Ken: If we start from $M_k$, then at some point we need to retake $k'$ so that the optimal transport plan $\nu$ is in $M_{k'}$? Maybe I'm misunderstanding something. Anyway, I have no preference!\\
%  Sam: Ok I need to understand better that part. }


% 	\paragraph{Step 1. The compact support of $\mu_0$ and $\mu_1$ are contained in $M$.}
%  Let $U:=\overline{\mathrm{Cvxhul}}(\supp(\mu_0)\cup\supp(\mu_1))$.
%  By the assumption (ii) and (iii),
%  there is a unique optimal transport plan $\nu\in\Prob(\Geo(\X,\di))$ joining $\mu_0$ and $\mu_1$.
%  Note that $\supp(\nu)\subset U$ holds by the assumption (i).
%  Furthermore, the results in \cite{Fathi-Figalli2010} imply that there exists a unique $W_2$-geodesic $(\mu_s)_{s \in \interval{0}{1}}$ joining $\mu_0$ to $\mu_1$, induced from the transportation map $T_s(x)=\exp_x(s\nabla \phi(x))$ for some $\mu_0$-almost everywhere differentiable function $\phi$.
%   Combined with the assumptions $\mathrm{dim}M \leq N$ and $\mathrm{Ric}_{V}^N \geq K$ on $ U\subset M$, \cite[Thm\, 17.36]{villani} implies that the $N$-Rényi entropy is $K$-convex along $(\mu_s)_{s \in \interval{0}{1}}$.

%  \todo[inline]{Sam: Is it possible that $M_k \subseteq M$ even without the assumption of strong geodesic convexity? so only with geodesic convexity. If yes and if my previous remark is also true, then we don't need \textbf{strong} geodesic convexity.\\
%  Sketch of idea of proof: Suppose that there is $p \in M_k$ with $p \in \partial M$. This means that we have a sequence of points $x_i, y_i$ such that $\diff(x_i, \partial M)$ and $\diff(y_i, \partial M) \geq 1/k$. Furthermore, on each unique geodesic $\gamma_i$ joining $x_i$ to $y_i$, there is a point $p_i$ such that $p_i$ tends to $p$ as $i \to +\infty$. Up to a subsequence, $x_i$ converges to $x$ and $y_i$ converges to $y$. Note that $\diff(x, \partial M)$ and $\diff(y, \partial M) \geq 1/k$ by continuity (this is crucial!). So $x, y \in M$ and $\notin \partial M$. I now want to say that the sequence of geodesics $\gamma_i$ converges to a geodesic $\gamma$ that contains $p$. I am a bit tired, correct me if I am wrong but this is because of Arzela-Ascoli and the fact that $M$ is uniquely geodesic. This is in contradiction with the geodesic convexity of $M$.\\
%  Sam again: I see the flaw in my idea. I don't know if $x_i$ (resp. $y_i$) really converges to some $x$ (resp. to $y$). The sequence could diverge, and then we have a problem. Could we take $M_k:=\overline{\mathrm{Cvxhul}}\{x\in M\mid k \geq \di(x,\X\setminus M)\geq \frac{1}{k}\}$ to save the argument?\\
%  Ken: I think the following slight modification should work!!!(Changed at 8/31)\\
%  Fix $x_0\in\X\setminus M$,
% and let $M_k:=\{x\in M\mid \frac{1}{k}\leq \di(x,\X\setminus M),~\di(x,x_0)\leq k\}$.
%  Since $\X$ is proper, $M_k$ is compact.
%  Choose $x_i,y_i,p_i,\gamma_i$ and $p$ as above (I think $\gamma_i$ need not be unique).
%  Furthermore, let $\eta_i$ be a concatenation of geodesics joining $x\to x_i\to p_i\to p\to p_i\to y_i\to y$.
%  Then the length of $\eta_i$ converges to $\di(x,y)$ as $i\to+\infty$.
%  Since $\X$ is proper, by the Ascoli--Arzela, $\eta_i$ sub-converges to a path $\eta$ (see also \cite[Theorem 2.5.14]{burago-burago-sergei2001}). Since a length functional is lower semicontinuous under pointwise convergence of paths (see \cite[Proposition 2.3.4]{burago-burago-sergei2001}), $\eta$ is a geodesic from $x$ to $y$ passing through $p$.\\
%  Ken: Now I'm convinced that the convexhull may not be a geodesic metric space, so we need to change the later arguments little bit.\\
%  Sam: Perfect! I am happy about it. I'll go over it again tomorrow or Monday.\\
%  Sam: The geodesic hull might not be geodesically convex. I am going through the argument again to see if it can work.}
% % Ken: Your diverging $x_i,y_i$ inspires the following example. Let's consider a metric on $\R^2$ induced from an o.n.f. $X=(1+y^2)\partial_x,Y=\partial_y$.
% % Then the half plane $\{x>0\}$ is geodesically convex, however, the convexhull of $M_{\epsilon_1,\epsilon_2}=\{u\in\R^2\mid \epsilon_2\geq \di(u,\{0\}\times \R)\geq \epsilon_1\}$ touches to the $y$-axis for any $\epsilon_1,\epsilon_2>0$.
% % Indeed, choose $y$ so that $\epsilon_1<\frac{1}{1+y^2}<\epsilon_2$. Then for any $\delta>0$,
% % $(\delta,y),(\delta,-y)$ are in $M_{\epsilon_1,\epsilon_2}$ and its geodesic is a straight segment.\\
% % However, in this example, we can take $M_k=\{(x,y)\mid x\geq \frac{1}{k}\}$, then it is desired one.}


% 	\paragraph{Step 2. The support of $\mu_0$ and $\mu_1$ are contained in $\X$.}
%  Let $M_k\subset M$ be the compact metric spaces given above.
%  The assumption (iii) implies that $M_k\to \X$ in the pointed measured Gromov--Hausdorff convergence with the inclusion map $\iota_k:M_k\to \X$ being approximation map.

%  Note that for any $k\in\N$,
%  there is $k'\in \N$ such that $\mathrm{Cvxhul}(M_k)\subset M_{k'}$.
%  In the same argument with \cite[Theorem 28.12]{villani},
%  $(\iota_{k'})_{\#}|_{\prob(M_k)}$ is also an approximation map.
%  Thus there are probability measures $\mu_i^k$ $(i=0,1)$ such that the supports are contained in $M_k$ and that $(\iota_{k'})_\#\mu_i^k \to\mu_i$ in $W_2$-distance.
%  Let $(\mu_s^k)_{s \in \interval{0}{1}}$ be the Wasserstein geodesic joining $\mu_0^k$ and $\mu_1^k$.
%  By the previous step, the $N$-Rényi entropy is $K$-convex along $(\mu_s^k)_{s \in \interval{0}{1}}$.
%  Furthermore,
% $(\mu_s^k)_{s \in \interval{0}{1}}$ uniformly converges to a $W_2$-geodesic $(\mu_s)_{s \in \interval{0}{1}}$ in $\X$ joining $\mu_0$ to $\mu_1$.
% The $N$-Rényi entropy is also $K$-convex along $(\mu_s)_{s \in \interval{0}{1}}$ since it is lower semicontinuous.

% \end{proof}
% }
The $\alpha$-Grushin spaces and half-spaces introduced in \cref{section:themodels} are metric measure spaces $(\X, \di, \mathfrak{m})$ with a weighted Riemannian manifold $(M, \diff_g, e^{-V}\vol_g)$ as their interior. For $N \in (-\infty,0)\cup[n,+\infty]$
% \todo{28. Here the domain should be $(-\infty,0)\cup (n,+\infty)$, with the convention that if $N=n$ then $V\equiv $ const... ?}
, we recall that the $N$-Ricci tensor of an $n$-dimensional weighted Riemannian $(M, \diff_g, e^{-V} \vol_g)$ is defined by
% \todo{29. DONE $\Ric_N\to \Ric_{N,V}$}
\begin{equation}
    \label{eq:NRicci}
    \Ric_{N,V}:= \begin{cases}
        \Ric & \text{ if } N=n, \\
        \Ric + \hes(V) & \text{ if } N=+\infty,\\
        \Ric+\hes(V)-\dfrac{\diff V\otimes\diff V}{N-n} & \text{ otherwise},
    \end{cases}
\end{equation}
with the convention that $V$ must be constant when $N=n$.

The following theorem provides sufficient conditions under which the differential condition $\Ric_{N,V} \geq K$ on $M$ is equivalent to $(\X, \di, \mathfrak{m})$ satisfying the $\mathsf{CD}(K, N)$ condition. The proof generalises the sketch found in \cite[Section 3.5]{rizzi2023failure}, which is specific to the $\alpha$-Grushin half-plane. This theorem is, to some extent, related to the conjecture stated in \cite{han2020}.
\begin{theorem}
    \label{theorem:equivalence}
    Let $K \in \R$, $N \in \ointerval{-\infty}{0} \cup \interval{n}{+\infty}$, $(\X, \di, \mathfrak{m})$ be a metric measure space and $M$ be an open subset of $\X$ such that
	\begin{enumerate}[label=\normalfont(\roman*)]
	\item $M$ is a geodesically convex
    % \todo{30. what is geodesically convex. DONE. }
    subset of $(\X, \di)$, i.e. for every $x,y\in M$, there is a geodesic joining $x$ and $y$ and any such curve is contained in $M$,
	\item $(M, \di|_M, \mathfrak{m}|_M)$ possesses a weighted $n$-dimensional Riemannian manifold structure,
	\item $\mathfrak{m}(\X \setminus M) = 0$.
 %{\red and $\X\setminus M\subset \partial M$}\todo{Sam: I still think we can easily remove it. See in the proof}.
	\end{enumerate}
	Then, the metric measure space $(\X, \di, \mathfrak{m})$ satisfies the $\cd(K,N)$ condition if and only if $\mathrm{Ric}_{N,V} \geq K$ on $M$.
\end{theorem}
% \todo[inline]{Ken: I think Remark 1.9 is not good. Since, to prove their $BE$-condition, we need to check Sobolev-to-Lipschitz property, see \cite{honda2018}.
% I suggest to prove the equivalence the $Ric^N\geq K$ and $\cd$, and use \cite{ledonne-lucic-pasqualetto2023}. Check the arguments in the end of this section.\\ Sam: Yes I agree. We might combine it in the main theorem? or separate like your wrote?\\ Sam again: so you mean that in general it might not be true, we are only sure it's true for $\alpha$-Grushin models.\\
% Ken: Yes, I guess $RCD$ holds in the general setting, however, it requires further discussions.}
\begin{proof}
   We provide the proof only for $N \in \interval{n}{+\infty}$. The details for $N \in \ointerval{-\infty}{0}$ are exactly the same, but one needs to replace the relevant key results used in the proof by the analogous ones when $N$ is negative. In particular, the Gromov--Hausdorff convergence which we use below must be replaced by the pointed $\mathsf{iKRW}$ convergence described in \cite{chiara-mattia-sosa2023}.

   Firstly, we note that the $\cd(K,N)$ condition on $(\X, \di, \mathsf{m})$ directly implies that $\mathrm{Ric}_{N,V} \geq K$ on $M$ by \cite[Theorem 17.36]{villani} (see also \cite[Part (e) of the proof of Theorem 1.7]{sturm2006}). We therefore focus now on the other implication.

	For a closed metric subspace $A\subseteq \X$, we shall denote by $W_2^A$ the Wasserstein distance on $A$. Consider $\mu_0, \mu_1 \in \mathscr{P}_2^{\mathrm{ac}}(\X, \di, \mathfrak{m})$ with continuous densities $\rho_0$ and $\rho_1$ respectively. We need to argue that there exists a $W_2^\X$-geodesic $(\mu_s)_{s \in \interval{0}{1}}$ joining $\mu_0$ to $\mu_1$ with $\mu_s = (e_s)_\sharp\nu$ for some $\nu \in \mathscr{P}(\mathrm{Geo}(\X, \di))$ such that the inequality in \cref{def.2.1} is satisfied.
    % the $N'$-Rényi entropy is $K$-convex
    % \todo{31. refer the relevant inequality in $\cd(K,N)$? DONE.}
    % along $(\mu_s)_{s \in \interval{0}{1}}$ for all $N' \geq N$.
    For all $k \in \N$ and given an arbitrary $x_0 \in \X$, we define the sets
\[
M_k:=\mathrm{cl}(\mathrm{geo}(N_k)), \text{ and } \ \ N_k := \{x\in M\mid k \geq \di(x_0,x), \text{ and }\di(x,\X\setminus M)\geq 1/k\},
\]
where $\mathrm{geo}(A)$ is the geodesic hull of a subset $A \subseteq \X$, i.e. the union of all geodesics starting at $x \in A$ and ending at $y \in A$. The set $N_k$ is clearly closed and bounded. Since $(\X, \diff, \mathfrak{m})$ is a complete metric measure space, the Heine--Borel property (see \cite[Theorem 2.5.28]{burago-burago-sergei2001}) implies that $N_k$ is also compact in $\X$.
% that $\mathrm{geo}(\{x\in M\mid k \geq \di(x,\X\setminus M)\geq 1/k\})$ is also bounded. Since $(\X, \diff, \mathfrak{m})$ is a (complete) metric measure space, .

    \paragraph{Step 1. The closure of the geodesic hull of a compact subset of $M$ is in $M$.}

    Let $K$ be a compact set of $\X$ contained in $M$. We claim that $\mathrm{cl}(\mathrm{geo}(K))$ is contained within $M$. Since $M$ is assumed to be geodesically convex in $\X$, we note that their closure in $\X$ satisfies $\mathrm{cl}(\mathrm{geo}(K)) \subseteq \mathrm{cl}(M)$. Because $M$ is open in $\X$, we also have that $\diff(x, \partial M) \geq \diff(K, \partial M) > 0$ for all $x \in K$. Suppose by contradiction that there exists $p \in \mathrm{cl}(\mathrm{geo}(K))$ such that $p \in \partial M$. This means that for all $n \in \N$, there exist $x_n, y_n \in K$, a geodesic $\gamma_n$ joining $x_n$ to $y_n$ contained in $M$, and a point $p_n \in M$ lying on $\gamma_n$ such that $p_n \to p$ as $n \to +\infty$. By compactness, we have, after extracting a converging subsequence, that $x_n \to x$ and $y_n \to y$ for some $x, y \in K$. In particular, the geodesics $\gamma_n$ have uniformly bounded lengths. Since $(\X, \diff)$ is a complete locally compact length space, the Arzela-Ascoli theorem from \cite[Theorem 2.5.14]{burago-burago-sergei2001} implies that the sequence of curves $\gamma_n$ contains a uniformly converging subsequence. This limit, which we denote by $\gamma$, is a curve of $\X$ with endpoints $x$ and $y$ and passing through $p$. Actually, the curve $\gamma$ is a minimising geodesics between $x$ and $y$ by \cite[Proposition 2.5.17]{burago-burago-sergei2001}. By geodesic convexity of $M$ assumed at (ii), the curve $\gamma$ is contained in $M$ while the point $p \notin M$, leading to a contradiction.

    % \paragraph{Step 2. The support of $\mu_0$ and $\mu_1$ are compact and contained in $M$.}
    % Let $U:=\overline{\mathrm{Cvxhul}}(\supp(\mu_0)\cup\supp(\mu_1))$.
    % By the assumption (ii) and (iii), there is a unique optimal transport plan $\nu\in\Prob(\Geo(\X,\di))$ joining $\mu_0$ and $\mu_1$.
    % Note that $\supp(\nu)\subset U$ holds by the assumption (i). Furthermore, the results in \cite{Fathi-Figalli2010} imply that there exists a unique $W_2$-geodesic $(\mu_s)_{s \in \interval{0}{1}}$ joining $\mu_0$ to $\mu_1$, induced from the transportation map $T_s(x)=\exp_x(s\nabla \phi(x))$ for some $\mu_0$-almost everywhere differentiable function $\phi$. Combined with the assumptions $\mathrm{dim}(M) \leq N$ and $\mathrm{Ric}_{V}^N \geq K$ on $ U\subset M$, \cite[Thm\, 17.36]{villani} implies that the $N$-Rényi entropy is $K$-convex along $(\mu_s)_{s \in \interval{0}{1}}$.

    \paragraph{Step 2. The support of $\mu_0$ and $\mu_1$ are compact and contained in $N_k$.}
    Assume that $\supp(\mu_0)$ and $\supp(\mu_1)$ are compact and that $\supp(\mu_0) \cup \supp(\mu_1) \subseteq N_k$ for some $k \in \N$. The previous step applied to the compact $N_k$ implies that each pair of points $x \in \supp(\mu_0)$ and $y \in \supp(\mu_1)$ are connected by a geodesic contained in $M$. By \cite[Theorem 13]{mccann}, there exists a unique $W_2^{M_k}$-geodesic $(\mu_s)_{s \in \interval{0}{1}}$ joining $\mu_0$ to $\mu_1$. By \cite[Proposition 2.10]{lott--villani2009}
    % \todo{Cor. 7.22? Sam: Maybe just chapter 7? Ken: The optimal transference plan $\Pi$ in the book is your $\nu?$\\Sam: I think yes, but it's better to check}
    , we also know that there exists $\nu \in \mathscr{P}(\mathrm{Geo}(\X, \di))$ such that $\mu_s = (e_s)_\sharp \nu$ for all $s \in \interval{0}{1}$ and $\mathrm{supp}(\nu) \subseteq \Gamma := (e_0 \times e_1)^{-1}(\mathrm{supp}(\mu_0) \times \mathrm{supp}(\mu_1))$. A geodesic in $\gamma \in \Gamma$ joins points in $\mathrm{supp}(\mu_0)$ with points in $\mathrm{supp}(\mu_1)$ and is contained in $M$ thanks to Step 1. Since $\mathrm{dim}(M) \leq N$ and $\mathrm{Ric}_{V}^N \geq K$ on $(M, g)$, \cite[Theorem\, 17.36]{villani} (or also \cite[Parts (a), (b), (c), and (d) of the proof of Theorem 1.7]{sturm2006}) implies that the inequality in \cref{def.2.1} for all $N' \geq N$.
    % $N'$-Rényi entropy is $K$-convex along $(\mu_s)_{s \in \interval{0}{1}}$ for all $N' \geq N$.
% \todo[inline]{ken: If we do not use pmGH convergence, then we may use \cite[Theorem 5.20, 29.20]{villani}, and arguments in Corollary 29.23}
    \paragraph{Step 3. The support of $\mu_0$ and $\mu_1$ are compact and contained in $\X$.}
    We can assume without loss of generality that $\X \setminus M \subseteq \partial M$. Indeed, when it comes to the $\mathsf{CD}(K, N)$ condition or the Gromov--Hausdorff convergence that we are going to discuss in this step, only the support of the measure matters (see \cite[Remark 3.1]{gigli2015}).
    % \todo{Sam: Remard 3.1 in Gigli for GH and find another reference for CD}
    The definition of $N_k$ and assumption (iii) imply that $N_k \to \X$ as $k \to +\infty$ in the pointed measured Gromov--Hausdorff convergence (see \cite[Definition 3.9]{gigli2015}), when taking the inclusion map $\iota_k:N_k\to \X$ as approximation maps.
 %    {\blue From this construction, there are sequences of probability measures $(\mu_0^k),(\mu_1^k)$ with compact in $N_k$ which converges to $\mu_0,\mu_1$ in $W_2^\X$.
 %    This is argued in
 % \cite[Proposition 28.7]{villani} where it is shown that the inclusion map $(\iota_k)_\#$ gives an approximation map $(\prob(N_k),W_2^{N_k})\to (\prob(\X),W_2^\X)$
 % % ,under the assumption that the support of $\mu_0^k,\mu_1^k,\mu_0,\mu_1$ are in a compact subset $K\subset \X$, which we can assume without loss of generality
 % }
Although $N_k$ is not a geodesically convex subspace, we have shown in Step 1 that $M_k := \mathrm{cl}(\mathrm{geo}(N_k)) \subseteq M$.
Then, an argument similar to the proof of \cite[Theorem 28.13]{villani} shows that the metric space $(\prob(N_k),W_2^{N_k})$ converges to $(\prob(\X),W_2^\X)$ in the geodesic local Gromov--Hausdorff topology, via the inclusion map $(\iota_k)_\#$. Hence there are sequences of probability measures $(\mu_i^k)_{k\in\N}$ $(i=0,1)$, supported on $N_k$, such that their pushforwards $(\iota_k)_\#\mu_i^k$ converges to $\mu_i$ in $(\prob_2(\X),W_2^\X)$.
    % \todo[inline]{Sam: why? I don't see why $\X \setminus M \subseteq \partial M$ is important.\\
    % Ken: If we don't have it, then it admits the example Half-Grushin $\cup\{(x,0)\mid x\leq 0\}$. Maybe we can ignore $\{(x,0)\mid x\leq 0\}$ by considering noly on the support of reference measure. I wanted to skip this arguments.}
    %{\red Note that for any $k\in\N$, there is $k'\in \N$ such that $\mathrm{geo}(N_k)\subseteq N_{k'}$. By the proof of \cite[Theorem 28.13]{villani}, the $(\iota_{k'})_{\#}|_{\prob(M_k)}$ is also an approximation map and thus there are probability measures $\mu_i^k$ $(i=0,1)$ such that the supports are contained in $N_k$ and that $(\iota_{k'})_\#\mu_i^k \to\mu_i$ in $W_2$-distance.}\todo{Sam: I am still unsure about this step, let's discuss it again.}
    % \todo[inline]{Sam: We need to introduce $\mu_0^k$ and $\mu_1^k$ and argue why they exist.}
    Let $(\mu_s^k)_{s \in \interval{0}{1}}$ be the Wasserstein geodesic joining $\mu_0^k$ and $\mu_1^k$ and $\nu_k$ be the optimal transport plan joining $\mu_0^k$ to $\mu_1^k$ in $M_k$. By Step 3, the inequality in \cref{def.2.1} holds
    % $N'$-Rényi entropy is $K$-convex
    along $(\mu_s^k)_{s \in \interval{0}{1}}$ for all $N' \geq N$. Furthermore, by \cite[Theorem 28.9 and Exercise 28.15]{villani}, $\nu_k$ weakly converges to an optimal transport plan $\nu$ in $\X$ joining $\mu_0$ to $\mu_1$, and the $W_2^{M_k}$-geodesics $(\mu_s^k)_{s \in \interval{0}{1}}$ uniformly converges to a $W_2^{\X}$-geodesic $(\mu_s)_{s \in \interval{0}{1}}$ in $\X$ joining $\mu_0$ to $\mu_1$. By the same argument as in \cite[Theorem\, 29.24 and 29.21]{villani}, we conclude that the inequality in \cref{def.2.1} remains valid
    % $N'$-Rényi entropy is also $K$-convex
    along $(\mu_s)_{s \in \interval{0}{1}}$ for all $N' \geq N$, by lower semicontinuity.

    \paragraph{Step 4. The general case.} The case where the support of $\mu_0$ and $\mu_1$ are not necessarily compact is obtained by the previous step and by using an exhaustion by compact sets, see \cite[Appendix E]{lott--villani2009}.
    % \todo{Sam: find a right reference, maybe \cite[Theorem 29.20(i,iii)]{villani} {\red find better reference}.}
\end{proof}

As for the $\mathsf{RCD}(K, N)$ condition (for $K \in \R$ and $N \in \interval{1}{+\infty}$), it is enough for our purposes to use the result of \cite[Theorem 1.2]{ledonne-lucic-pasqualetto2023}, which shows that
sub-Riemannian manifolds equipped with a non-negative Radon measure are infinitesimally Hilbertian. Although the results in \cite{ledonne-lucic-pasqualetto2023} assume the bracket-generating condition and our $\alpha$-Grushin half-spaces do not when $\alpha \notin \N$, their Finsler approximation techniques can still be applied.
It is sufficient to regard the $\alpha$-Grushin half-spaces as $\alpha$-Grushin (full) spaces equipped with a measure supported on the half-space, for the following theorem to follow directly.
\begin{theorem}
% \todo[inline]{32. Instead of using LLP23, can we use the stability of infinitesimal Hilbertianity within $\cd$?\\
% Ken: Probably not. Since our $M_k$ may not me geodesically convex, and in particular, it may not be geodesic $\cd(K,N)$ space. To prove the stability, we need the $1$-st order convergence i.e. the convergence of the gradient flow in $L^2$ or the null $(1,2)$-capacity condition, which are unclear to us. Therefore, we can not say that $M_k$ is $\mathsf{CD}/\mathsf{RCD}$.}
     For $K \in \R$ and $N \in \interval{1}{+\infty}$, the $\cd(K,N)$ condition is equivalent to the $\RCD(K,N)$ condition for the $\alpha$-Grushin half-spaces.
\end{theorem}
% \begin{remark}
%     Although the results in \cite{ledonne-lucic-pasqualetto2023} assume the bracket generating condition,
%     their Finsler approximation techniques can be applied to $\alpha$-Grushin model spaces with $\alpha\notin \Z$.

% \end{remark}

% \todo[inline]{Copy-Paste. We do not need to keep it. Sam: Yes, the discussion of ``infinitesimally Hilbertianity'' can be done in the equivance theorem (the $\mathsf{R}$ of $\mathsf{RCD}$)}


% \newcommand{\sob}{\mathrm W}
% \newcommand{\sobh}{\mathrm{HW}}
% \newcommand{\slope}{{\mathsf D}}
% \newcommand{\be}{\mathsf{BE}}

% \begin{definition}[Bakry--Émery inequality] We say that a sub-Riemannian manifold $(M,\diff)$ endowed with a positive smooth measure $\m$ satisfies the \emph{Bakry--Émery $\be(K,\infty)$ inequality}, for $K\in \R$, if
% \begin{equation}
% \label{eq:be}
% \frac12\,\Delta(\|\nabla u\|^2)
% \ge
% \g(\nabla u, \nabla \Delta u)
% +
% K\|\nabla u\|^2
% \quad
% \text{for all}\ u \in C^\infty(M),
% \end{equation}
% where $\Delta$ is the corresponding sub-Laplacian, and $\nabla$ the sub-Riemannian gradient.
% \end{definition}

% \begin{theorem}[Infinitesimal Hilbertianity]\label{res:R}
% Let $q\in(1,\infty)$.
% Let $(M,\diff)$ be a complete sub-Riemannian manifold equipped with a positive smooth measure $\m$.
% The following hold.
% \begin{enumerate}[label=(\roman*),ref=\roman*,itemsep=.5ex]

% \item \label{item:sob=sob}
% $\sob^{1,q}(M,\diff,\m)=\sobh^{1,q}(M,\m)$, with
% $|\slope u|_{w,q}=\|\nabla u \|$ $\m$-a.e.\ on $M$ for all $u\in\sob^{1,q}(M,\diff,\m)$.
% In particular, taking $q=2$, $(M,\diff,\m)$ is infinitesimally Hilbertian.

% \item\label{item:bochner}
% If $(M,\diff,\m)$ satisfies the $\mathsf{CD}(K,\infty)$ condition for some $K\in\R$, then the Bakry--Émery $\be(K,\infty)$ inequality \eqref{eq:be} holds on $M$.
% \end{enumerate}
% \end{theorem}

% Note that \ref{res:R} holds for less regular measures, see \ref{rmk:moregeneralmeasures}.
% \begin{remark}\label{rmk:moregeneralmeasures}
% The above proofs work for more general measures $\m$. Namely, we can assume that, locally on any bounded coordinate neighborhood $\Omega \subset \R^n$, $\m = \theta \mathscr{L}^n$ with $\theta \in \sob^{1,1}(\Omega,\mathscr{L}^n) \cap \mathcal{L}^\infty(\Omega,\mathscr{L}^n)$. In this case, the positivity of $\m$ corresponds to the requirement that $\theta$ is locally essentially bounded from below away from zero, in charts.
% \end{remark}

% \begin{remark}[The case of a.e.\ smooth measures]\label{rmk:inf_hilb_local}
% \ref{res:R} can be adapted also to the case of a Borel and locally finite measure $\m$ which is smooth and positive only on $\overline{\Omega}$, where $\Omega \subset M$ is an open set with $\m(\partial\Omega)=0$. In this case, we obtain $\sobh^{1,q}(\Omega,\m)= \sob^{1,q}(\overline{\Omega},\diff,\m)$, with $|\slope u|_{w,q}=\|\nabla u \|$  $\m$-a.e.\ on $\Omega$ for all $u\in \sob^{1,q}(\overline{\Omega},\diff,\m)$.
% In particular, if $\m$ is smooth and positive out of a closed set $\mathcal{Z}$, with $\m(\mathcal{Z})=0$, an elementary approximation argument proves that $(M,\diff,\m)$ is infinitesimally Hilbertian and, if $(M,\diff,\m)$ satisfies the $\cd(K,\infty)$ condition for $K\in \R$, then the Bakry-Émery $\be(K,\infty)$ inequality \eqref{eq:be} holds on $M\setminus \mathcal{Z}$. This is the case, for example, of the Grushin planes and half-planes with weighted measures of Section \ref{sec:weigthedgrush}. The proof follows the same argument of the one of \ref{res:R}, exploiting the locality of the $q$-upper gradient, see for example~\cite[Sec.\ 8.2]{AmbrosioGigliSavare13} and~\cite[Prop.\ 2.6]{Gigli12}, and similar properties for the distributional derivative.
% \end{remark}



\section{Generalised Ricci curvature of \texorpdfstring{$\alpha$}{alpha}-Grushin half-spaces}
\label{section:riccicomputation}

The previous section shows that establishing the validity of the $\mathsf{CD}(K, N)$ condition in the $\alpha$-Grushin half-spaces introduced in \cref{section:themodels} is equivalent to a computation of the Bakry--\'Emery Ricci curvature. The following simple computation will be handy.

\begin{lemma}
\label{e:Rictensor}
    Let $(M, g)$ be a $2$-dimensional Riemannian manifold, $\mathcal{U} \subseteq M$ an open set, and $N \in (-\infty, 0) \cup (2, +\infty]$. Assume that $(x, y) : \mathcal{U} \to \mathbb{R}^2$ is a chart and that we are given two smooth functions $f : \mathcal{U} \to \mathbb{R}\setminus \{0\}$ and $V : \mathcal{U} \to \mathbb{R}$ that only depends on $x$. If, in this coordinate,
    \[
    g = \diff x \otimes \diff x + \frac{1}{f(x)^2} \diff y \otimes \diff y,
    \]
    then it holds that
\begin{equation*}
        \begin{aligned}
            \Ric_{N,V} = \left[ \left(\frac{f'}{f}\right)' - \left(\frac{f'}{f}\right)^2 + V'' - \frac{(V')^2}{N - 2}\right] \diff x \otimes \diff x
            + \frac{1}{f^2}\left[\left(\frac{f'}{f}\right)' - \left(\frac{f'}{f}\right)^2 - \frac{f'}{f} V'\right] \diff y \otimes \diff y,
        \end{aligned}
\end{equation*}
where all the derivatives are understood with respect to the variable $x$.
% \begin{equation}\label{e:Rictensor}
%     \begin{aligned}
%        \Ric_{N,V}=
%        \left[\frac{\partial_x^2f\cdot f-2(\partial_xf)^2}{f^2}+\partial_x^2 V - \frac{(\partial_x V)^2}{N-2}\right]dx\otimes dx\\
%        \qquad\qquad+ \left[\frac{\partial_xf\cdot f-2(\partial_x f)^2}{f^4}-\frac{\partial_x f}{f^3}\partial_x V\right]dy\otimes dy.
%     \end{aligned}
% \end{equation}
\end{lemma}

\begin{proof}
The two vector fields $X = \partial_x$ and $Y = f(x) \partial_y$ defined on $\mathcal{U}$ form a family of $g$-orthonormal fields. The only non-zero bracket relation is
\[
[X, Y] = \frac{f'}{f} Y.
\]
Using the Koszul formula, we easily obtain
\[
\nabla_X X = \nabla_X Y = 0, \qquad \nabla_Y X = - \frac{f'}{f} Y, \qquad \nabla_Y Y = \frac{f'}{f} X.
\]
It follows that the only non-zero entry of the Riemann curvature tensor is
\[
\mathrm{R}(X, Y, X, Y) = \left(\frac{f'}{f}\right)' - \left(\frac{f'}{f}\right)^2.
\]
Using the symmetries, we obtain
\[
\mathrm{Ric} = \left[\left(\frac{f'}{f}\right)' - \left(\frac{f'}{f}\right)^2\right] \diff x \otimes \diff x + \frac{1}{f^2}\left[\left(\frac{f'}{f}\right)' - \left(\frac{f'}{f}\right)^2\right] \diff y \otimes \diff y.
\]
We also find that $\nabla V =X(V)X$, and that the only non-zero entries of the Hessian are
\[\hes(V)(X,X) = V'',~~\hes(V)(Y,Y)=-\frac{f'}{f}V'.\]
Finally, one has the non-zero entry
\[\diff V\otimes \diff V(X,X)=(V')^2,\]
and the proof is complete with \cref{eq:NRicci}.
\end{proof}

% {\blue Given the metric tensor
% \[g=dx\otimes dx + f(x)^{-2}dy\otimes dy,\]
% and the weighted measure
% \[\m=e^{-V(x)}dx dy,\]
% we have that
% \begin{equation}\label{e:Rictensor}
%     \begin{aligned}
%        \Ric_{N,V}=\left[\frac{\partial_x^2f\cdot f-2(\partial_xf)^2}{f^2}+\partial_x^2 V - \frac{(\partial_x V)^2}{N-2}\right]dx\otimes dx\\
%        \qquad\qquad+ \left[\frac{\partial_xf\cdot f-2(\partial_x f)^2}{f^4}-\frac{\partial_x f}{f^3}\partial_x V\right]dy\otimes dy.
%     \end{aligned}
% \end{equation}
% }
% \todo{This computation should be helpful to shorten the proof}
% \todo[inline]{Sam: I think we don't need to recall the formulas from Riemannian geometry, i.e. Ricci tensor in coordinates, Christoffel symbols etc. Just a thought.\\
% Yes, I wrote them so that I can check it easily. I think the computation on the paper should be as least as possible.}
% {\red Adopting Einstein's summation convention with $(x_1, x_2) := (x, y)$, we recall that in coordinates,
% \begin{equation}
%     \label{eq:pureRicci}
%     \Ric =\left[\partial_a\Gamma_{ij}^a-\partial_j\Gamma_{ai}^a+\Gamma_{ab}^a\Gamma_{ij}^b-\Gamma_{ib}^a\Gamma_{aj}^b\right]\diff x_i\otimes \diff x_j,
% \end{equation}
% where the Christoffel symbols $\Gamma_{ij}^k$ are obtained by
% \begin{equation}
%     \label{eq:christoffelsymbols}
%     \Gamma_{ij}^k:=\frac{1}{2}\left[\partial_i g_{ja}+\partial_j g_{ia}-\partial_a g_{ij}\right]g^{ka}.
% \end{equation}
% Furthermore, the Hessian of a function $V : M \to \R$ can be computed as
% \begin{equation}
% \label{eq:Hessian}
% \hes(V):=\left[\partial_i\partial_j-\Gamma_{ij}^k\partial_k\right]V\diff x_i\otimes \diff x_j.
% \end{equation}
% }
\paragraph{The $\alpha$-Grushin hemisphere.}
We make use of the coordinate chart and the notation described in \cref{section:alphaGrushin}.

% {\red \begin{prop}
% \label{prop:NRiccialphaGrushin}
% The $N$-Ricci curvature of the $\alpha$-Grushin open hemisphere $(\Sbb_\alpha^+, \diff_{\S_\alpha}, \mathfrak{m}^\beta_{\S_\alpha})$ is given by
% \begin{align*}
%     \Ric_{N,V}=&-\frac{1}{\sin^2(x)}\left[3\alpha-1+(\alpha-1)^2\cos^2(x) - \beta + \frac{\beta^2}{N-2}\cos^2(x) \right]\diff x\otimes \diff x \\
%     &-\frac{\cos^2(x)}{\sin^{2\alpha+2}(x)}\left[3\alpha-1+(\alpha-1)^2\cos^2(x) - \beta \left(1+(\alpha-1)\cos^2(x)\right)\right]\diff y\otimes \diff y.
% \end{align*}
% \end{prop}

% \begin{proof}
%     Using \cref{eq:christoffelsymbols}, we get that the Christoffel symbols of this Riemannian space are given by
% \begin{align*}
% &\Gamma_{ij}^1=\begin{cases}
%     0 & (i,j)=(1,1),(1,2),(2,1),\\
%     \displaystyle \frac{\cos(x)\left[1+(\alpha-1)\cos^2(x)\right]}{\sin^{2\alpha}(x)\sin(x)} & (i,j)=(2,2),\end{cases}\\
% &\Gamma_{ij}^2=\begin{cases}
%     0 & (i,j)=(1,1),(2,2),\\
%     \displaystyle -\frac{1+(\alpha-1)\cos^2(x)}{\sin(x)\cos(x)} & (i,j)=(1,2),(2,1).
% \end{cases}
% \end{align*}
% The expression for $\Ric_N$ follows from \cref{eq:pureRicci}, \cref{eq:Hessian}, and \cref{eq:NRicci} together with $V(x, y) =-\beta\log\abs{\sin(x)}$ from \cref{def:MeasurealphaGrushinSphere}.
% \end{proof}

% A necessary and sufficient condition to have a lower bound on $\Ric_N$ can now be found.
% }
\begin{prop}
    Given $K \in \R$, $N \in \ointerval{-\infty}{0} \cup \linterval{2}{+\infty}$, $\alpha \geq 0$ and $\beta \geq \alpha$, the $N$-Ricci curvature of the $\alpha$-Grushin open hemisphere $(\Sbb_\alpha^+, \diff_{\S_\alpha}, \mathfrak{m}^\beta_{\S_\alpha})$ is
\begin{align*}
    \Ric_{N,V}=&-\frac{1}{\sin^2(x)}\left[3\alpha-1+(\alpha-1)^2\cos^2(x) - \beta + \frac{\beta^2}{N-2}\cos^2(x) \right]\diff x\otimes \diff x \\
    &-\frac{\cos^2(x)}{\sin^{2\alpha+2}(x)}\left[3\alpha-1+(\alpha-1)^2\cos^2(x) - \beta \left(1+(\alpha-1)\cos^2(x)\right)\right]\diff y\otimes \diff y.
\end{align*}
Furthermore, it holds $\Ric_{N,V} \geq K$ if and only if
    \[
    \beta-\alpha^2-\alpha+\min\left(-K+(\alpha-1)^2,~-\frac{\beta^2}{N-2},~\beta(\alpha-1)\right)\geq 0.
    \]
\end{prop}

\begin{proof}
 The first part is done by \cref{e:Rictensor} with $f(x)=\frac{\abs{\sin(x)}^\alpha}{\cos(x)}$ and $V(x)=-\beta\log\abs{\sin(x)}$. For the second part, we see by comparing the coefficients of the tensors $\diff x\otimes \diff x$ and $\diff y\otimes \diff y$ respectively that $\Ric_{N,V}\geq K$ holds if and only if
\begin{equation*}\label{ineq:Ricbound1}
-3\alpha+1+\beta-K+\left[K-(\alpha-1)^2-\frac{\beta^2}{N-2}\right]\cos^2(x)\geq 0\end{equation*}
and
\begin{equation*}\label{ineq:Ricbound2}
    -3\alpha+1+\beta-K+\left[K-(\alpha-1)^2+\beta(\alpha-1)\right]\cos^2(x)\geq 0.
\end{equation*}
Since $\cos(x)$ takes its value in $[0,1)$ on $\S_\alpha^+$,
the above two inequalities are equivalent to the following inequality:
\[-3\alpha+1+\beta-K+\min\left(0,K-(\alpha-1)^2-\frac{\beta^2}{N-2},K-(\alpha-1)^2+\beta(\alpha-1)\right)\geq 0,\]
which is trivially equivalent to the inequality in the statement.
\end{proof}
% \todo[inline]{Sam: This is just a suggestion. I suggest that in section 2 we fully define the $\alpha$-spaces up to them being metric measure spaces. In this section, only computation of Ricci curvature. I feel it's better for the reader. (but if you think otherwise, let's change)\\
% Ken: I agree, it should be in section 2!!}

% {\red For the $\alpha$-Grushin sphere $(\S^2,\g_\alpha)$,
% the unweighted Riemannian (Popp's) measure is
% \[\vol=\cos(x)|\sin(x)|^{-\alpha}\diff x\diff y.\]
% For $\beta\geq 0$,
% we consider the weighted measure given by
% \[\m_\beta:=|\sin(x)|^\beta\vol=\cos(x)|\sin(x)|^{\beta-\alpha}\diff x \diff y=e^{-V}\vol,\]
% where $V(x,y):=-\beta\log|\sin(x)|$.}


% {\red \begin{theorem}
%     For the $\alpha$-Grushin hemisphere with $\beta$-weighted measure $(\Splu,\g_\alpha,\m_\beta)$,
%     $\Ric_N\geq K$ holds if and only if the following inequality holds:
% \[\beta-\alpha^2-\alpha+\min\left(-K+(\alpha-1)^2,~-\frac{\beta^2}{N-2},~\beta(\alpha-1)\right)\geq 0.\]
% \end{theorem}

% \begin{remark}
% Suppose that $\Ric_N\geq K$ holds for some $\beta>0$ and $N\in[1,+\infty)$.
% Then $-\frac{\beta^2}{N-2}< 0$,
% and as a consequence,
% there is $\bar K>0$ such that $\Ric_N\geq \bar K$ holds.
%     \end{remark}
% \begin{proof}

% The metric tensor $\g_\alpha$ is
% \[\g_\alpha=\diff x\otimes \diff x+\frac{\cos^2(x)}{\sin^{2\alpha}(x)}g_{\S^{n-1}}.\]
% The Christoffel symbol $\Gamma_{ij}^k:=\frac{1}{2}\left[\partial_i g_{ja}+\partial_j g_{ia}-\partial_a g_{ij}\right]g^{ka}$ is
% \begin{align*}
%     &\Gamma_{ij}^1=\begin{cases}
%     0 & (i,j)=(1,1),(1,2),(2,1),\\
%     \frac{\cos(x)\left[1+(\alpha-1)\cos^2(x)\right]}{\sin^{2\alpha}(x)\sin(x)} & (i,j)=(2,2),\end{cases}\\
%     &\Gamma_{ij}^2=\begin{cases}
%     0 & (i,j)=(1,1),(2,2),\\
%     -\frac{1+(\alpha-1)\cos^2(x)}{\sin(x)\cos(x)} & (i,j)=(1,2),(2,1).
% \end{cases}
% \end{align*}
% % The covariant derivative $\nabla_i\partial_j:=\Gamma_{ij}^k\partial_k$ is
% % \[\nabla_1\partial_1=0,~~\nabla_1\partial_2=\nabla_2\partial_1=-\frac{1+(\alpha-1)\cos^2(x)}{\sin(x)\cos(x)}\partial_2,~~\nabla_2\partial_2=\frac{\cos(x)\left[1+(\alpha-1)\cos^2(x)\right]}{\sin^{2\alpha}(x)\sin(x)}\partial_1.\]
% Then the $N$-Bakry--Emery Ricci tensor $\Ric_{N}:=\Ric+\hes(V)-\frac{\diff V\otimes\diff V}{N-2}\geq K$ will be computed as follows.
% The Ricci tensor $\Ric:=[\partial_a\Gamma_{ij}^a-\partial_j\Gamma_{ai}^a+\Gamma_{ab}^a\Gamma_{ij}^b-\Gamma_{ib}^a\Gamma_{aj}^b]\diff x_i\otimes \diff x_j$ is
% \begin{align*}
%     \Ric=&-\frac{1}{\sin^2(x)}\left[3\alpha-1+(\alpha-1)^2\cos^2(x)\right]\diff x\otimes \diff x\\
%     &-\frac{\cos^2(x)}{\sin^{2\alpha+2}(x)}\left[3\alpha-1+(\alpha-1)^2\cos^2(x)\right]\diff y\otimes \diff y.
% \end{align*}

% The Hessian $\hes(V):=[\partial_i\partial_j-\Gamma_{ij}^k\partial_k]V\diff x_i\otimes \diff x_j$ is

% \begin{align*}
%     \hes(V)=
%     \frac{\beta}{\sin^2(x)}\diff x\otimes \diff x + \frac{\beta\cos^2(x)\left[1+(\alpha-1)\cos^2(x)\right]}{\sin^{2\alpha+2}(x)}\diff y\otimes \diff y.
% \end{align*}


% The last term is computed as
% \[-\frac{\diff V\otimes\diff V}{N-2}=-\frac{\beta^2\cos^2(x)}{(N-2)\sin^2(x)} \diff x\otimes \diff x.\]
% Therefore $\Ric_N\geq K$ holds if and only if
% \begin{equation}\label{ineq:Ricbound1}
% -3\alpha+1+\beta-K+\left\{K-(\alpha-1)^2-\frac{\beta^2}{N-2}\right\}\cos^2(x)\geq 0\end{equation}
% and
% \begin{equation}\label{ineq:Ricbound2}
%     -3\alpha+1+\beta-K+\Big\{K-(\alpha-1)^2+\beta(\alpha-1)\Big\}\cos^2(x)\geq 0
% \end{equation}
% hold, by comparing the coefficients of the tensors $\diff x\otimes \diff x$ and $\diff y\otimes \diff y$ respectively.
% Since $\cos(x)\in[0,1]$,
% the two inequalities hold for any $x\geq 0$ if and only if
% \[-3\alpha+1+\beta-K+\min\left(0,K-(\alpha-1)^2-\frac{\beta^2}{N-2},K-(\alpha-1)^2+\beta(\alpha-1)\right)\geq 0\]
% holds, which is equivalent to the inequality in the statement.


% \end{proof}}

\paragraph{The $\alpha$-Grushin hyperbolic half-plane.}

% The unweighted Riemannian (Popp's) measure is
% \[\vol=\cosh(x)|\sinh(x)|^{-\alpha}\diff x\diff y.\]
% For $\beta\geq 0$,
% we consider the weighted measure
% \[\m_\beta:=|\sinh(x)|^\beta\vol=\cosh(x)|\sinh(x)|^{\beta-\alpha}\diff x\diff y=e^{-f}\vol,\]
% where $f(x,y):=-\beta\log|\sinh(x)|$.

We make use of the coordinate chart and the notation described in \cref{section:alphaGrushin2}.

% {\red \begin{prop}
% \label{prop:NRiccialphaGrushin2}
% The $N$-Ricci curvature of the $\alpha$-Grushin open half hyperbolic plane $(\bbH_\alpha^+, \diff_{\bbH_\alpha}, \mathfrak{m}^\beta_{\bbH_\alpha})$ is given by
% \begin{align*}
%     \Ric_{N,V}=&-\frac{1}{\sinh^2(x)}\left[3\alpha-1+(\alpha-1)^2\cosh^2(x) - \beta + \frac{\beta^2}{N-2}\cosh^2(x) \right]\diff x\otimes \diff x \\
%     &-\frac{\cosh^2(x)}{\sin^{2\alpha+2}(x)}\left[3\alpha-1+(\alpha-1)^2\cosh^2(x) - \beta \left(1+(\alpha-1)\cosh^2(x)\right)\right]\diff y\otimes \diff y.
% \end{align*}
% \end{prop}

% \begin{proof}
%     Using \cref{eq:christoffelsymbols}, we get that the Christoffel symbols of this Riemannian space are given by
% \begin{align*}
% &\Gamma_{ij}^1=\begin{cases}
%     0 & (i,j)=(1,1),(1,2),(2,1),\\
%     \displaystyle \frac{\cosh(x)\left[1+(\alpha-1)\cosh^2(x)\right]}{\sinh^{2\alpha}(x)\sinh(x)} & (i,j)=(2,2),\end{cases}\\
% &\Gamma_{ij}^2=\begin{cases}
%     0 & (i,j)=(1,1),(2,2),\\
%     \displaystyle -\frac{1+(\alpha-1)\cosh^2(x)}{\sinh(x)\cosh(x)} & (i,j)=(1,2),(2,1).
% \end{cases}
% \end{align*}
% The expression for $\Ric_N$ follows from \cref{eq:pureRicci}, \cref{eq:Hessian}, and \cref{eq:NRicci} together with $V(x, y) =-\beta\log\abs{\sinh(x)}$ from \cref{def:hyperbolicmeasure}.
% \end{proof}

% A necessary and sufficient condition to have a lower bound on $\Ric_N$ can now be found.}

\begin{prop}
    Given $K \in \R$, $N \in \ointerval{-\infty}{0} \cup \linterval{2}{+\infty}$, $\alpha \geq 0$ and $\beta \geq \alpha$, the $N$-Ricci curvature of the $\alpha$-Grushin open hyperbolic half-plane $(\bbH_\alpha^+, \diff_{\bbH_\alpha}, \mathfrak{m}^\beta_{\bbH_\alpha})$ is
\begin{align*}
    \Ric_{N,V}=&-\frac{1}{\sinh^2(x)}\left[3\alpha-1+(\alpha-1)^2\cosh^2(x) - \beta + \frac{\beta^2}{N-2}\cosh^2(x) \right]\diff x\otimes \diff x \\
    &-\frac{\cosh^2(x)}{\sin^{2\alpha+2}(x)}\left[3\alpha-1+(\alpha-1)^2\cosh^2(x) - \beta \left(1+(\alpha-1)\cosh^2(x)\right)\right]\diff y\otimes \diff y.
\end{align*}
    Furthermore, it holds $\Ric_{N,V} \geq K$ if and only if
    \[
    \min\left(-K-(\alpha-1)^2,\beta-\alpha^2-\alpha\right)+\min\left(-\frac{\beta^2}{N-2},~\beta(\alpha-1)\right)\geq 0
    \]
\end{prop}

\begin{proof}
 The first part is done by \cref{e:Rictensor} with $f(x)=\frac{\abs{\sinh(x)}^\alpha}{\cosh(x)}$ and $V(x)=-\beta\log\abs{\sinh(x)}$. For the second part, we see by comparing the coefficients of the tensors $\diff x\otimes \diff x$ and $\diff y\otimes \diff y$ respectively that $\Ric_{N,V}\geq K$ holds if and only if
\begin{equation*}\label{ineq:Ricbound1-2}
-3\alpha+1+\beta+K+\left[-K-(\alpha-1)^2-\frac{\beta^2}{N-2}\right]\cosh^2(x)\geq 0\end{equation*}
and
\begin{equation*}\label{ineq:Ricbound2-2}
    -3\alpha+1+\beta+K+\left[-K-(\alpha-1)^2+\beta(\alpha-1)\right]\cosh^2(x)\geq 0.
\end{equation*}
Since $\cosh(x)\in[1,+\infty)$,
the two inequalities hold for all $x\geq 0$ if and only if
 \[\min\left(0,-3\alpha+1+\beta+K\right)+\min\left(-K-(\alpha-1)^2-\frac{\beta^2}{N-2},-K-(\alpha-1)^2+\beta(\alpha-1)\right)\geq 0\]
holds, which is equivalent to the inequality in the statement.

\end{proof}

% {\red
% \begin{theorem}
%     For the $\alpha$-Grushin hyperbolic halfplane with $\beta$-weighted measure $(\bbH_+,\g_\alpha,\m_\beta)$,
%     $\Ric_N\geq K$ holds if and only if the following inequality holds:
% \[\min\left(-K-(\alpha-1)^2,\beta-\alpha^2-\alpha\right)+\min\left(-\frac{\beta^2}{N-2},~\beta(\alpha-1)\right)\geq 0\]
% \end{theorem}
% \begin{remark}

%     As a consequence, we have the followings.
%     \begin{itemize}
%         \item For any $\alpha\geq 0$,
%         there are $\beta\geq 0$,
%         $K<0$ and $N\in[1,+\infty)$ such that $\Ric_N\geq K$.
%         \item For any $\alpha,\beta\geq0$ and $N\in[1,+\infty)$, it does not satisfy $\Ric_N\geq 0$.
%         \item For any $\alpha\geq 1$,
%         there are $\beta>0$ and $N\in(-\infty,0)$ such that $\Ric_N\geq 0$.
%     \end{itemize}
% \end{remark}

% \begin{proof}

% The metric tensor is
% % $\g_\alpha$ is
%     \[\g_\alpha=\diff x\otimes \diff x+\frac{\cosh^2(x)}{\sinh^{2\alpha}(x)}\diff y\otimes \diff y.\]
% The Christoffel symbol is
% % $\Gamma_{ij}^k=\frac{1}{2}\left[\partial_i g_{ja}+\partial_j g_{ia}-\partial_a g_{ij}\right]g^{ka}$ is
% \begin{align*}
%     &\Gamma_{ij}^1=\begin{cases}
%     0 & (i,j)=(1,1),(1,2),(2,1),\\
%     \frac{\cosh(x)\left[1+(\alpha-1)\cosh^2(x)\right]}{\sinh^{2\alpha+1}(x)} & (i,j)=(2,2),\end{cases},\\
%     &\Gamma_{ij}^2=\begin{cases}
%     0 & (i,j)=(1,1),(2,2),\\
%     -\frac{1+(\alpha-1)\cosh^2(x)}{\sinh(x)\cosh(x)} & (i,j)=(1,2),(2,1).
% \end{cases}
% \end{align*}
% The Ricci tensor is
% % $\Ric=\partial_a\Gamma_{ij}^a-\partial_j\Gamma_{ai}^a+\Gamma_{ab}^a\Gamma_{ij}^b-\Gamma_{ib}^a\Gamma_{aj}^b$ is
% \begin{align*}
%     \Ric=&-\frac{1}{\sinh^2(x)}\left[3\alpha-1+(\alpha-1)^2\cosh^2(x)\right] \diff x\otimes \diff x\\
%     &-\frac{\cosh^2(x)}{\sinh^{2\alpha+2}(x)}\left[3\alpha-1+(\alpha-1)^2\cosh^2(x)\right]\diff y\otimes \diff y.
% \end{align*}

% The Hessian is
%  % $\hes(V)$ is

% \begin{align*}
%     \hes(V)=\frac{\beta}{\sinh^2(x)}\diff x\otimes \diff x + \frac{\beta\cosh^2(x)\left[1+(\alpha-1)\cosh^2(x)\right]}{\sinh^{2\alpha+2}(x)}\diff y\otimes \diff y.
% \end{align*}


% The last term is
% \[\frac{\diff V\otimes\diff V}{N-2}=-\frac{\beta^2\cosh^2(x)}{(N-2)\sinh^2(x)}\diff x\otimes\diff y.\]
% Therefore $\Ric_{N}\geq K$ holds if and only if
% \begin{equation}\label{ineq:hyperRicbound1}
% -3\alpha+1+\beta+K+\left\{-K-(\alpha-1)^2-\frac{\beta^2}{N-2}\right\}\cosh^2(x)\geq 0,\end{equation}
% and
% \begin{equation}\label{ineq:hyperRicbound2}
%     -3\alpha+1+\beta+K+\left\{-K-(\alpha-1)^2+\beta(\alpha-1)\right\}\cosh^2(x)\geq 0.
% \end{equation}
% Since $\cosh(x)\in[1,+\infty)$,
% the two inequalities hold for any $x\geq 0$ if and only if
% \[\min\left(0,-3\alpha+1+\beta+K\right)+\min\left(-K-(\alpha-1)^2-\frac{\beta^2}{N-2},-K-(\alpha-1)^2+\beta(\alpha-1)\right)\geq 0\]
% holds, which is equivalent to the inequality in the statement.

% \end{proof}}

\paragraph{The $\infty$-Grushin half-plane.}

% The unweighted Riemannian (Popp's) measure is
% \[\vol=e^{\frac{1}{|x|}}\diff x\diff y.\]
% For $\beta,\gamma\geq 0$,
% we consider the weighted measure
% \[\m_{\beta,\gamma}:=|x|^\beta e^{-\frac{\gamma}{|x|^2}}\vol=|x|^\beta e^{-\frac{\gamma}{|x|^2}+\frac{1}{|x|}}\diff x\diff y=e^{-V}\vol,\]
% where $V(x,y):=\frac{\gamma}{|x|^2}-\beta\log|x|$,



We make use of the coordinate chart and the notation described in \cref{section:alphaGrushin3}.
% {\red
% \begin{prop}
% \label{prop:NRiccialphaGrushin3}
% The $N$-Ricci curvature of the $\infty$-Grushin open half-plane $(\G_\infty^+, \diff_{\G_\infty}, \mathfrak{m}^{\beta,\gamma}_{\G_\infty})$ is given by
% \begin{align*}
%     \Ric_{N,V}=&\left[\frac{6\gamma-1}{|x|^4}-\frac{2}{|x|^3}+\frac{\beta}{|x|^2}-\frac{1}{N-2}\left(\frac{4\gamma^2}{|x|^6}+\frac{4\beta\gamma}{|x|^4}+\frac{\beta^2}{|x|^2}\right)\right]\diff x\otimes \diff x\\
%     &+\left[\frac{2\gamma}{|x|^5}+\frac{1}{|x|^4}+\frac{\beta}{|x|^3}\right]e^{\frac{2}{|x|}}\diff y\otimes\diff y.
% \end{align*}
% \end{prop}

% \begin{proof}
%     Using \cref{eq:christoffelsymbols}, we get that the Christoffel symbols of this Riemannian space are given by
%     \begin{align*}
%     &\Gamma_{ij}^1=\begin{cases}
%     0 & (i,j)=(1,1),(1,2),(2,1),\\
%     \frac{e^{\frac{2}{|x|}}}{x|x|} & (i,j)=(2,2),\end{cases}\\
%     &\Gamma_{ij}^2=\begin{cases}
%     0 & (i,j)=(1,1),(2,2),\\
%     -\frac{1}{x|x|} & (i,j)=(1,2),(2,1).
% \end{cases}
% \end{align*}
% The expression for $\Ric_N$ follows from \cref{eq:pureRicci}, \cref{eq:Hessian}, and \cref{eq:NRicci} together with $V(x, y) =\frac{\gamma}{|x|^2}-\beta\log\abs{x}$ from \cref{def:infty_measure}.
% \end{proof}

% A necessary and sufficient condition to have a lower bound on $\Ric_N$ can now be found.
% }

\begin{prop}
    Given $K \in \R$, $N \in \ointerval{-\infty}{0} \cup \linterval{2}{+\infty}$ and $\beta,
    \gamma\geq 0$, the $N$-Ricci curvature of the $\infty$-Grushin open half-plane $(\G_\infty^+, \diff_{\G_\infty}, \mathfrak{m}^{\beta,\gamma}_{\G_\infty})$ is
\begin{align*}
    \Ric_{N,V}=&\left[\frac{6\gamma-1}{|x|^4}-\frac{2}{|x|^3}+\frac{\beta}{|x|^2}-\frac{1}{(N-2)x^6}\left(2\gamma+\beta x^2\right)^2\right]\diff x\otimes \diff x\\
    &+\left[\frac{2\gamma}{|x|^5}-\frac{1}{|x|^4}+\frac{\beta-2}{|x|^3}\right]e^{\frac{2}{|x|}}\diff y\otimes\diff y.
\end{align*}
    It holds $\Ric_{N,V} \geq K$ if and only if the following inequalities hold for any $x\geq 0$:
% {\red \[-\frac{\gamma^2}{N-2}+\left(-1+6\gamma-\frac{4\beta\gamma}{N-2}\right)t^2-2t^3+\left(\beta-\frac{\beta^2}{N-2}\right)t^4- Kt^6\geq 0,\]
% and
% \[2\gamma+t+\beta t^2-Kt^5\geq 0.\]}
\[(6\gamma-1)x^2-2x^3+\beta x^4-\frac{1}{N-2}(2\gamma+\beta x^2)^2\geq Kx^6~~\text{and}~~2\gamma-x+(\beta-2)x^2\geq Kx^5.\]
In particular, $\Ric_{\infty,V}\geq 0$ holds if and only if
\[\min\{\beta(6\gamma-1),8(\beta-2)\gamma\}\geq 1.\]
\end{prop}

\begin{proof}
% From \cref{prop:NRiccialphaGrushin3}, we see by comparing the coefficients of the tensors $\diff x\otimes \diff x$ and $\diff y\otimes \diff y$ respectively that $\Ric_N\geq K$ holds if and only if
% \begin{equation*}\label{ineq:Ricbound1-3}
% -3\alpha+1+\beta-K+\left[K-(\alpha-1)^2-\frac{\beta^2}{N-2}\right]\cosh^2(x)\geq 0\end{equation*}
% and
% \begin{equation*}\label{ineq:Ricbound2-3}
%     -3\alpha+1+\beta-K+\left[K-(\alpha-1)^2+\beta(\alpha-1)\right]\cosh^2(x)\geq 0.
% \end{equation*}
The proof is reached by using \cref{e:Rictensor} with $f(x)=e^{-1/|x|}$ and $V(x)=\frac{\gamma}{|x|^2}-\beta\log|x|$, and by comparing the coefficients of $\diff x\otimes \diff x$ and $\diff y\otimes \diff y$ respectively.
\end{proof}

\cref{thm:main1}, \cref{thm:main2}, \cref{thm:main3} follow from the Ricci computations this section, noting that the validity of the assumptions in \cref{theorem:equivalence} are verified because of the results of \cref{section:themodels}.

% {\red
% \begin{theorem}
%     For the $\infty$-Grushin halfplane with $(\beta,\gamma)$-weighted measure $(\R^2_+,\g_\infty,\m_{\beta,\gamma})$,
%     $\Ric_N\geq K$ holds if and only if the following inequalities hold for any $t\geq 0$:
% \[-\frac{\gamma^2}{N-2}+\left(-1+6\gamma-\frac{4\beta\gamma}{N-2}\right)t^2-2t^3+\left(\beta-\frac{\beta^2}{N-2}\right)t^4- Kt^6\geq 0,\]
% and
% \[2\gamma+t+\beta t^2-Kt^5\geq 0.\]
% \end{theorem}
% \begin{remark}
%     As a consequence,
%     we have the followings:
%     \begin{itemize}
%         \item If $(\R_+,\g_\infty,\m_{\beta})$ satisfies $\Ric_N\geq K$,
%         then $K\leq 0$, $N\in[-\infty,2)\cup \{+\infty\}$ and $\gamma\neq 0$.
%         \item Furthermore, if $\Ric_N\geq 0$ holds, then $\beta\neq 0$.
%     \end{itemize}
% \end{remark}
% \begin{proof}

% The metric tensor is
% \[\g_\infty= \diff x\otimes \diff x+e^{\frac{2}{|x|}}\diff y\otimes \diff y.\]
% The Christoffel symbol is
% \begin{align*}
%     &\Gamma_{ij}^1=\begin{cases}
%     0 & (i,j)=(1,1),(1,2),(2,1),\\
%     \frac{e^{\frac{2}{|x|}}}{x|x|} & (i,j)=(2,2),\end{cases}\\
%     &\Gamma_{ij}^2=\begin{cases}
%     0 & (i,j)=(1,1),(2,2),\\
%     -\frac{1}{x|x|} & (i,j)=(1,2),(2,1).
% \end{cases}
% \end{align*}
% The Ricci tensor is
% \begin{align*}
%     \Ric=\left(-\frac{1}{|x|^4}-\frac{2}{|x|^3}\right)\diff x\otimes \diff x +\frac{1}{|x|^4}e^{\frac{2}{|x|}}\diff y\otimes\diff y.
% \end{align*}
% the Hessian is

% \begin{align*}
%     \hes(V)=\frac{6\gamma}{|x|^4}+\frac{\beta}{|x|^2}\diff x\otimes\diff x + \left[\frac{2\gamma}{|x|^5}+\frac{\beta}{|x|^3}\right]e^{\frac{2}{|x|}}\diff y\otimes\diff y.
% \end{align*}
% The last term of $\Ric_N$ is
% \[-\frac{\diff V\otimes\diff V}{N-2}=-\frac{1}{N-2}\left[\frac{4\gamma^2}{|x|^6}+\frac{4\beta\gamma}{|x|^4}+\frac{\beta^2}{|x|^2}\right]\diff x\otimes\diff x.\]
% Therefore $\Ric_N\geq K$ holds if and only if
% \[-\frac{\gamma^2}{N-2}+\left(-1+6\gamma-\frac{4\beta\gamma}{N-2}\right)x^2-2x^3+\left(\beta-\frac{\beta^2}{N-2}\right)x^4- Kx^6\geq 0,\]
% and
% \[2\gamma+x+\beta x^2-Kx^5\geq 0\]
% hold for any $x\geq 0$.
% \end{proof}
% }

% \section{try: higher dimensional draft}

% On the $n$-dimensional Riemannian sphere $(\S^n, g_{\S^n})$, we introduce the following coordinate chart, the validity of which can be found in \cite{boscain-prandi-seri2016}. Fix a $(n-1)$-dimensional sub-sphere $\gamma:\S^{n-1}\to \S^n$,
% and let $N$ and $S$ be the north pole and south pole of $\S^n$ respectively, with respect to $\gamma$. For $p\in \S^n\setminus\{N,S\}$,
% we define $y := y(p) \in \S^{n-1}$ as the unique point such that $\gamma(y)$ is the perpendicular foot from $p$ to $\mathrm{Im}(\gamma)$. Furthermore, the number $x := x(p) \in \ointerval{-\pi/2}{\pi/2}$ is defined to be the (spherical) distance $\diff_{\S^{n}}(p,\gamma(y))$. The map $\S^n \to \ointerval{-\pi/2}{\pi/2}\times \S^{n-1}  : p \mapsto (x(p), y(p))$ is a well-defined coordinate chart which can be naturally extended to $N$ and $S$ under the identification
% $(\frac{\pi}{2},y^1) \sim ( \frac{\pi}{2},y^2)$ and $( -\frac{\pi}{2},y^1) \sim (-\frac{\pi}{2},y^2)$ for any $y^1,y^2\in \S^{n-1}$. It is easily shown that in this coordinates, the spherical Riemannian metric tensor $g_{\S^{n}}$ can be written as
% \[
% g_{\S^n} = \diff x \otimes \diff x + \cos^2(x) g_{\S^{n-1}}.
% \]
% % \begin{remark}
% % The coordinate system $(x,y)$ is, up to rotations, the same as the spherical coordinates $(\theta, \phi_2,\dots,\phi_{n-1}) $ which parametrises the standard sphere $\S^{n} \subseteq \R^{n+1}$ by
% % \[(\theta,\phi^2,\dots,\phi_n)\mapsto (\cos(\theta)\cos(\phi),
% % \sin(\theta)\cos(\phi),\sin(\phi)).\]
% % \end{remark}
% With this in mind, we can introduce the $\alpha$-Grushin sphere and hemisphere.
% \begin{definition}
%     For $\alpha \geq 0$, the $n$-dimensional $\alpha$-Grushin sphere $\S_\alpha$ is the sub-Riemannian structure on $\S^n$ induced from the vector field $X$ and $Y_1^\alpha,\dots,Y^\alpha_{n-1}$ given by
% \[
% X:=\partial_{x},~~Y_i^\alpha:=\frac{\abs{\sin(x)}^{\alpha}}{\cos(x)}\partial_{y_i},
% \]
% where $\partial_{x_i}$ $(i=2,\dots,n)$ are a local orthonormal frame of $\S^{n-1}$.
% Note that the induced metric is independent of the choice of an orthonormal frame of $\S^{n-1}$.
% The $\alpha$-Grushin hemisphere $\S^+_\alpha$ (resp. open hemisphere $M^+_\alpha$) is the subset of $\S_\alpha$ defined by
% \[
% \S^+_\alpha := \{ p \in \S^n \mid x \in \interval{0}{\pi/2} \} \text{ (resp. $M^+_\alpha := \{ p \in \S^2 \mid x \in (0,\pi/2] \}$)}.
% \]
% \end{definition}
% The $0$-Grushin sphere is simply the two-dimensional Riemannian sphere $\S^n$. When $\alpha > 0$, the $\alpha$-Grushin sphere is a two-dimensional almost-Riemannian structure with $\{x = 0\}$ being its set of singular points. The Grushin sphere studied in \cite{boscain-prandi-seri2016,Pan23} corresponds to the $1$-Grushin sphere. At non-singular points, this sub-Riemannian structure admits the Riemannian metric
% \begin{equation}
%     \label{eq:SalphaRiemannianMetric}
%     g_{\S_{\alpha}} = \diff x \otimes \diff x + \frac{\cos^2(x)}{\sin^{2\alpha}(x)} g_{\S^{n-1}}.
% \end{equation}
% A simple computation shows that the Popp measure, i.e. the unweighted Riemannian volume induced from \cref{eq:SalphaRiemannianMetric}, is given by
% \[
% \vol_{\S_\alpha}=\cos(x)\abs{\sin(x)}^{-\alpha}\diff x\wedge \vol_{\S^{n-1}}.
% \]
% We introduce the following weighted measure.
% \begin{definition}
% For $\beta\geq \alpha$,
% we consider the weighted measure given by
% \begin{equation}
%     \label{eq:MeasurealphaGrushinSphere}
%     \m^\beta_{\S_\alpha}:=\abs{\sin(x)}^\beta\vol_{\S_\alpha}=\cos(x)\abs{\sin(x)}^{\beta-\alpha}\diff x \vol_{\S^{n-1}}=e^{-V_{\S_\alpha}}\vol_{\S_\alpha},
% \end{equation}
% where $V_{\S_\alpha}(x,y):=-\beta\log\abs{\sin(x)}$.
% \end{definition}


% \begin{theorem}
%     For the $n$-dimensional $\alpha$-Grushin hemisphere with $\beta$-weighted measure $(\Splu,\g_\alpha,\m_\beta)$,
%     $\Ric_N\geq K$ holds if and only if the following inequality holds:
% \[\beta-(n-1)\alpha(\alpha-1)+\min\left(-K+(n-1)(\alpha-1)^2,~-\frac{\beta^2}{N-2},~\beta(\alpha-1)+(n-2)(3\alpha-1)\alpha\right)\geq 0.\]
% \end{theorem}

% \begin{remark}
% Suppose that $\Ric_N\geq K$ holds for some $\beta>0$ and $N\in[1,+\infty)$.
% Then $-\frac{\beta^2}{N-2}< 0$,
% and as a consequence,
% there is $\bar K>0$ such that $\Ric_N\geq \bar K$ holds.
%     \end{remark}
% \begin{proof}
% Fix a local orthonormal frame $\partial y,\dots,\partial x_n$ of $\S^{n-1}$,
% and we identify $g_{\S^{n-1}}=\sum_{i=2}^n\diff x_i\otimes \diff x_i$.
% The Christoffel symbol $\Gamma_{ij}^k:=\frac{1}{2}\left[\partial_i g_{ja}+\partial_j g_{ia}-\partial_a g_{ij}\right]g^{ka}$ is
% \begin{align*}
%     &\Gamma_{ij}^1=\begin{cases}
%     \frac{\cos(x)\left[1+(\alpha-1)\cos^2(x)\right]}{\sin^{2\alpha}(x)\sin(x)} & (i,j)=(2,2),(3,3),\dots,(n,n)\\
%     0 & else,
%     \end{cases}\\
%     &\Gamma_{ij}^k=\begin{cases}
%     -\frac{1+(\alpha-1)\cos^2(x)}{\sin(x)\cos(x)} & (i,j)=(1,k),(k,1),\\
%     0 & else,
% \end{cases}~~~~\text{ for }k\geq 2.
% \end{align*}
% % The covariant derivative $\nabla_i\partial_j:=\Gamma_{ij}^k\partial_k$ is
% % \[\nabla_1\partial_1=0,~~\nabla_1\partial_2=\nabla_2\partial_1=-\frac{1+(\alpha-1)\cos^2(x)}{\sin(x)\cos(x)}\partial_2,~~\nabla_2\partial_2=\frac{\cos(x)\left[1+(\alpha-1)\cos^2(x)\right]}{\sin^{2\alpha}(x)\sin(x)}\partial_1.\]
% Then the $N$-Bakry--Emery Ricci tensor $\Ric_{N}:=\Ric+\hes(V)-\frac{\diff V\otimes\diff V}{N-2}$ will be computed as follows.
% The Ricci tensor $\Ric:=[\partial_a\Gamma_{ij}^a-\partial_j\Gamma_{ai}^a+\Gamma_{ab}^a\Gamma_{ij}^b-\Gamma_{ib}^a\Gamma_{aj}^b]\diff x_i\otimes \diff x_j$ is
% \begin{align*}
%     \Ric=&-\frac{n-1}{\sin^2(x)}\left[3\alpha-1+(\alpha-1)^2\cos^2(x)\right]\diff x\otimes \diff x\\
%     &-\frac{\cos^2(x)}{\sin^{2\alpha+2}(x)}\left[-2(n-2)(\alpha-1)+3\alpha-1+(-2n+5)(\alpha-1)^2\cos^2(x)\right]g_{\S^{n-1}}.
% \end{align*}

% The Hessian $\hes(V):=[\partial_i\partial_j-\Gamma_{ij}^k\partial_k]V\diff x_i\otimes \diff x_j$ is

% \begin{align*}
%     \hes(V)=
%     \frac{\beta}{\sin^2(x)}\diff x\otimes \diff x + \frac{\beta\cos^2(x)\left[1+(\alpha-1)\cos^2(x)\right]}{\sin^{2\alpha+2}(x)}g_{\S^{n-1}}.
% \end{align*}


% The last term is computed as
% \[-\frac{\diff V\otimes\diff V}{N-2}=-\frac{\beta^2\cos^2(x)}{(N-2)\sin^2(x)} \diff x\otimes \diff x.\]
% Therefore $\Ric_N\geq K$ holds if and only if
% \begin{equation}\label{ineq:Ricbound1}
% (n-1)(-3\alpha+1)+\beta-K+\left\{K-(n-1)(\alpha-1)^2-\frac{\beta^2}{N-2}\right\}\cos^2(x)\geq 0\end{equation}
% and
% \begin{equation}\label{ineq:Ricbound2}
%     2(n-2)(\alpha-1)-3\alpha+1+\beta-K+\Big\{K+(2n-5)(\alpha-1)^2+\beta(\alpha-1)\Big\}\cos^2(x)\geq 0
% \end{equation}
% hold, by comparing the coefficients of the tensors $\diff x\otimes \diff x$ and $g_{\S^{n-1}}$ respectively.
% Since $\cos(x)\in[0,1]$ and $n\geq 2$,
% the two inequalities hold for any $x\geq 0$ if and only if
% \begin{align*}
%     \beta-K+\min\Bigg(&(n-1)(-3\alpha+1),\\
%     &(n-1)(-3\alpha+1)+K-(n-1)(\alpha-1)^2-\frac{\beta^2}{N-2},\\
%     &K+(2n-5)(\alpha-1)^2+\beta(\alpha-1)+2(n-2)(\alpha-1)-3\alpha+1\Bigg)\geq 0
% \end{align*}
% holds, which is equivalent to the inequality in the statement.


% \end{proof}

%\begin{spacing}{0.85}
  \printbibliography
%\end{spacing}

% \bibliographystyle{alpha}
% \bibliography{bibliography}
\end{document}
